% concatenated from main.tex
%2multibyte Version: 5.50.0.2960 CodePage: 1250
\documentclass{article}%
\usepackage{amssymb}%
\usepackage{amsmath, bbm, mathtools}%
\setcounter{MaxMatrixCols}{30}%
\usepackage{amsfonts}%
\usepackage{authblk}
\usepackage{graphicx}
\usepackage[a4paper, left=2.5cm, right=2.5cm, top=3cm, bottom=3cm]{geometry}
\usepackage{shuffle}
\usepackage{csquotes}
\usepackage{comment}
\usepackage{multirow}
\usepackage{subcaption}
\usepackage{float}

%TCIDATA{OutputFilter=latex2.dll}
%TCIDATA{Version=5.50.0.2960}
%TCIDATA{Codepage=1250}
%TCIDATA{CSTFile=40 LaTeX article.cst}
%TCIDATA{Created=Tuesday, August 05, 2025 20:03:03}
%TCIDATA{LastRevised=Tuesday, September 02, 2025 12:31:39}
%TCIDATA{<META NAME="GraphicsSave" CONTENT="32">}
%TCIDATA{<META NAME="SaveForMode" CONTENT="1">}
%TCIDATA{BibliographyScheme=Manual}
%TCIDATA{<META NAME="DocumentShell" CONTENT="Standard LaTeX\Blank - Standard LaTeX Article">}
%TCIDATA{Language=American English}
%BeginMSIPreambleData
\providecommand{\U}[1]{\protect\rule{.1in}{.1in}}
%EndMSIPreambleData
\newtheorem{theorem}{Theorem}
\newtheorem{acknowledgement}[theorem]{Acknowledgement}
\newtheorem{algorithm}[theorem]{Algorithm}
\newtheorem{axiom}[theorem]{Axiom}
\newtheorem{case}[theorem]{Case}
\newtheorem{claim}[theorem]{Claim}
\newtheorem{conclusion}[theorem]{Conclusion}
\newtheorem{condition}[theorem]{Condition}
\newtheorem{conjecture}[theorem]{Conjecture}
\newtheorem{corollary}[theorem]{Corollary}
\newtheorem{criterion}[theorem]{Criterion}
\newtheorem{definition}[theorem]{Definition}
\newtheorem{example}[theorem]{Example}
\newtheorem{exercise}[theorem]{Exercise}
\newtheorem{lemma}[theorem]{Lemma}
\newtheorem{notation}[theorem]{Notation}
\newtheorem{problem}[theorem]{Problem}
\newtheorem{proposition}[theorem]{Proposition}
\newtheorem{remark}[theorem]{Remark}
\newtheorem{solution}[theorem]{Solution}
\newtheorem{summary}[theorem]{Summary}
\newenvironment{proof}[1][Proof]{\noindent\textbf{#1.} }{\ \rule{0.5em}{0.5em}}

\newcommand{\floor}[1]{\left \lfloor {#1}\right \rfloor}

\newcommand{\ceil}[1]{\left \lceil {#1}\right \rceil}
\allowdisplaybreaks

\title{Signature Kernel and Schwinger-Dyson Kernel Equations as Two-Parameter Rough Differential Equations}
\author{Thomas Cass, Dan Crisan, Andrea Iannucci\thanks{Corresponding author: \texttt{andrea.iannucci22@imperial.ac.uk}}  , William F. Turner}

\affil{Department of Mathematics, Imperial College London}


\begin{document}
\maketitle
\begin{abstract}
We develop a rough-path framework for two-parameter rough differential equations on rectangular and simplicial domains, motivated by the signature kernel and Schwinger--Dyson kernel equations. The theory is formulated in spaces of jointly controlled rough paths and is based on a robust two-parameter rough integration framework. In particular, we introduce a notion of rough integration over two-dimensional simplices at low regularity extending the results of \cite{cass2023fubini} and \cite{gerasimovivcs2019hormander}.

Within this setting, we show that the signature kernel equation arises naturally as a two-parameter rough differential equation and establish well-posedness and stability. We also extend the Schwinger--Dyson kernel equation, previously formulated for bounded-variation paths, to rough driving signals, proving existence and uniqueness in appropriate controlled rough path spaces.

In the smooth rough path regime, we relate the resulting equations to PDE and integro-differential formulations. Finally, we derive and analyse a numerical scheme for the rough Schwinger--Dyson equation, including runtime and memory complexity estimates, and illustrate its performance with numerical experiments.
\end{abstract}


\noindent\textbf{MSC 2020:}
Primary 60L20, 60L10, 60L90; secondary 60L70, 35R09, 65L20.

\medskip

\noindent\textbf{Keywords:}
rough paths; two-parameter rough differential equations; controlled rough paths;
signature kernels; Schwinger--Dyson equations; rough integration.

\section{Introduction}

The theory of rough paths provides a deterministic extension of classical differential equation theory, allowing one to give meaning to equations
\[
dZ = V(Z)\, dX, \qquad Z_0 = z,
\]
driven by highly irregular signals, well beyond the Young integration regime, by enriching the driving signal $X$ with additional algebraic structure.

In several recent developments, equations of an analogous type have appeared in genuinely two-parameter form. A first motivating example is the signature kernel equation, introduced in \cite{salvi2021signature} and used in learning problems for sequential data,
\begin{equation}\label{sig_ker_intro}
K(t,v) = 1 + \int_a^t \int_c^v K(s,u)\, \langle dX_s, dY_u \rangle.
\end{equation}
A second motivating example is the Schwinger--Dyson kernel equation, which arises in the study of kernels obtained as limits of random developments of a path in the unitary group \cite{cass2024free},
\begin{equation}\label{sd_ker_intro}
K(s,t) = 1 - \int_s^t \int_s^r K(s,u)K(u,r)\, \langle dX_u, dX_r \rangle,
\end{equation}
and has also been connected to questions in quantum field theory; see \cite{crew2025quantum}. 
In both cases, the driving signals are naturally too irregular to be treated by classical two-variable integration once one leaves the bounded-variation setting.
Beyond providing motivating examples of genuinely two-parameter equations, these kernels also arise in the construction of estimators for the maximum mean discrepancy between distributions on path space; see \cite[Sections 7--8]{chevyrev2022signature} and \cite{chevyrev2016characteristic}. This observation has been used in \cite{issa2023non} and, in the bounded-variation setting, in \cite{lou2023pcf} to show that such kernels can be employed to generate time-series samples from certain target distributions. Their study in the rough setting is therefore particularly relevant, since many driving signals of interest in applications are not of bounded variation but nevertheless admit a natural enhancement to a geometric rough path. Important examples include Brownian motion, more general continuous semimartingales, and sample paths of solutions to stochastic differential equations. The connection we establish between these kernels and two-parameter rough differential equations, and, under additional regularity assumptions, partial differential equations, provides a rough-path framework in which these kernel identities can be interpreted, analysed, and approximated numerically.
Our viewpoint is that \eqref{sig_ker_intro} and \eqref{sd_ker_intro} are instances of a broader class of two-parameter rough differential equations driven by a pair of rough paths $\mathbf{X}$ and $\mathbf{Y}$.
More precisely, we study equations of the form
\begin{equation}\label{skeleton_equation_intro}
Z_{t,v}
= z
+ \int_{\mathcal{D}_{t,v}}
\Theta^g\left(Z\right)_{s,u}\,
d(\mathbf{X}_s,\mathbf{Y}_u),
\end{equation}
posed either on rectangular or simplicial domain $\mathcal{D}_{t,v}$ on which the parameters mapping $\Theta^g$ depends.
We introduce a notion of solution to \eqref{skeleton_equation_intro} as a fixed point in suitable spaces of jointly controlled paths and establish well-posedness and stability estimates for our motivating examples, yielding in particular continuous dependence on the driving rough paths.
This structural well-posedness also leads naturally to numerical approximation schemes.

On the analytical side, our development relies on recent progress in two-parameter rough integration.
Two largely independent approaches have been proposed.
On the one hand, \cite{chouk2014rough} develops an integration theory for rough sheets when the driving signal has finite $p$-variation with $2\le p<3$.
On the other hand, \cite{gerasimovivcs2019hormander} introduces the class of two-parameter controlled paths in the same $p$-variation regime and, under a smooth approximation assumption, establishes a Fubini-type theorem for rough integration over rectangles and two-dimensional simplices.
In this work we adopt the controlled-path viewpoint, as subsequently extended in \cite{cass2023fubini} to a general $p$-variation setting, including the construction of a joint rough integral and a rough Fubini theorem on rectangles that does not rely on smooth approximations.
Motivated by the structure of the Schwinger--Dyson kernel equation, we supplement this theory with a notion of rough integration over simplices at low regularity, thereby bridging the gap between the approaches of \cite{cass2023fubini} and \cite{gerasimovivcs2019hormander}.
In addition, in order to make \eqref{skeleton_equation_intro} well defined, we show that the class of two-parameter controlled paths is stable under composition with smooth functions.

Additionally, when the driving signals are smooth rough paths, introduced in \cite{bellingeri2022smooth}, the two-parameter rough differential equation viewpoint connects naturally with partial differential equations.
In particular, interpreting \eqref{sig_ker_intro} within our framework recovers the PDE derived in \cite{lemercier2024log} in a straightforward manner, and the corresponding uniqueness and approximation statements for the kernel follow from the general two-parameter RDE theory developed here.

Similarly, when $\mathbf{X}$ is a smooth rough path, the Schwinger--Dyson kernel equation admits an equivalent reformulation at the level of an integro-differential equation.
More precisely, exploiting combinatorial identities, derived  from the Schwinger--Dyson equations, we show that is possible to connect the trace component of the kernel to the remaining components of the associated two-parameter controlled path leading to a system of integro-differential equation.
This provides an alternative analytic description of the kernel in the smooth regime and may serve as a starting point for further qualitative analysis.

These combinatorial identities also lead to a numerical scheme, for which we study runtime and memory complexity.
We conclude with numerical experiments illustrating the performance of the proposed approximation schemes: the experiments compare several variants of our method and show that their agreement deteriorates as the driving signal becomes less regular, in line with the theoretical picture.
We also report runtime studies that empirically validate the complexity estimates derived above.







\begin{comment}
\paragraph{Outline of this paper}
The paper is organized as follows. In Section 2, we provide a brief review of the theory of two-parameter controlled paths. Section 3 is dedicated to the construction of the two-parameter rough integral over two-dimensional simplices and its connection to the two-parameter rough integral introduced in \cite{cass2023fubini}. In Section 4, we study two-dimensional rough differential equations. We begin by analysing the composition of two-parameter controlled rough paths with smooth functions and formalize the notion of solution to a two-parameter rough differential equation. Using this framework, we establish uniqueness for the signature kernel equation \eqref{sig_ker_intro} and examine its connection to the corresponding PDE when the driving signals $\mathbf{X}$ and $\mathbf{Y}$ are smooth rough paths.

We then derive the rough version of the Schwinger–Dyson kernel equation and prove well-posedness for the associated two-parameter rough differential equation.

In the last  section a numerical scheme is presented for this equation, including numerical experiments and comparisons with the schemes introduced in \cite{cass2024free}.
Finally, we conclude by deriving an integro-differential equation satisfied by the Schwinger–Dyson kernel in the case where $\mathbf{X}$ is a smooth rough path.


\end{comment}


\paragraph{Outline of the paper}
The paper is organized as follows. In Section~2, we briefly review and relate the theories of one- and two-parameter controlled paths, with an emphasis on the structural analogies that underlie our later constructions. We also collect the main ingredients of two-parameter rough integration used throughout the paper, including the controlled-path formulation and the joint rough integral on rectangles.

In Section~3, motivated by the Schwinger--Dyson kernel equation in the rough regime, we introduce a notion of two-parameter rough integration over two-dimensional simplices at low regularity. We then show that this simplex integral is consistent with the joint rough integral on rectangles introduced in \cite{cass2023fubini}, thereby placing both constructions within a common framework.

Section~4 is devoted to two-parameter rough differential equations. We first prove that the class of two-parameter controlled paths is stable under composition with smooth functions, which is a fundamental ingredient in defining nonlinear equations of the form \eqref{skeleton_equation_intro}. We then formalize the notion of solution to a two-parameter rough differential equation as a fixed point in a suitable jointly controlled path space and develop the corresponding well-posedness theory.

We next apply this framework to our motivating examples. For the signature kernel equation \eqref{sig_ker_intro}, we show that it arises naturally as a two-parameter rough differential equation, and we establish well-posedness and stability, including continuous dependence on the driving rough paths. In the smooth rough path regime, we also recover its associated PDE formulation. We then derive the rough Schwinger--Dyson kernel equation and prove existence and uniqueness for the corresponding two-parameter rough differential equation. In the same smooth regime, we further derive an integro-differential reformulation of the Schwinger--Dyson kernel based on the combinatorial identities developed in the paper.

In the final section, we turn to numerical approximation. We present a numerical scheme for the rough Schwinger--Dyson equation, analyse its runtime and memory complexity, and report numerical experiments illustrating the performance of the method across different regularity regimes. We also include comparisons with the schemes introduced in \cite{cass2024free}, together with runtime experiments that support the theoretical complexity estimates.























\section{Preliminaries for two-parameter controlled paths}
We use this section to develop the theory of two-parameter (jointly controlled) rough paths.To make the definitions accessible it is helpful to keep in mind the analogy with classical, one-parameter rough path theory. The primary achievement of the latter theory is well-posedness and stability for a class of equations of the form
\[
Z_t = Z_0 + \int_{0}^t V(Z_s)\, d\mathbf{X}_s ,
\]
where $\mathbf{X}$ denotes a driving rough path, allowing one to go beyond classical integration theory, which would require bounded variation of the driving signal. Interpreting these equations requires a definition of the rough path $\mathbf{X}$, its properties, and the class of paths that can be integrated against it. The structure of the equation which describe the kernel trick for the signature kernel in the setting of continuous bounded variation paths 
\begin{equation}\label{sig_ker_intro__}
K(t,v) = 1 + \int_a^t \int_c^v K(s,u)\, \langle dX_s, dY_u \rangle,
\end{equation} 
and the Schwinger–Dyson kernel trick \eqref{sd_ker_intro} suggests that the extension of these concepts to the rough path setting is to view them as particular examples of two-parameter rough differential equations of the form
\begin{equation}\label{skeleton_equation}
Z_{t,v}
= z
+ \int_{\mathcal{D}_{t,v}}
\Theta^g\left(Z\right)_{s,u}
\, d(\mathbf{X}_s, \mathbf{Y}_u),
\end{equation}
where our two core examples are obtained from the following choices:
\[
\begin{aligned}
\text{Signature kernel:}\quad
&\begin{aligned}
\mathcal{D}_{t,v} &= [a,t]\times[c,v], \\
\Theta^g(Z)(s,u) &= g(Z)_{s,u}, \\
g(z_1,z_2)[\eta,\xi] &= z_1 \langle \eta,\xi\rangle ;
\end{aligned}
\\[8pt]
\text{Schwinger-Dyson  kernel:}\quad
&\begin{aligned}
\mathcal{D}_{t,v} &= \Delta_{[t,v]}, \\
\Theta^g(Z) &= g(Z_{t, \cdot},Z)_{s,u}, \\
g(z_1,z_2)[\eta,\xi] &= -z_1 z_2 \langle \eta,\xi\rangle .
\end{aligned}
\end{aligned}
\]

The aim of the next three sections is to develop a framework for two-parameter rough differential equations that is sufficiently robust to accommodate our core examples. We will show that both examples fit this framework in a precise sense. As a consequence, the general results of rough path theory become available in this setting, including stability and uniform estimates, and the structure of the solutions naturally leads to numerical approximation schemes.

Two-parameter controlled rough paths were first introduced in \cite{gerasimovivcs2019hormander} and further developed in \cite{cass2023fubini}. Building on this theory, we introduce the additional tools needed for our programme: two-parameter controlled paths, the joint rough integral, and a rough Fubini theorem.

For background on classical rough path theory we refer to \cite{cass2022combinatorial, friz2014course}, and for a detailed treatment of smooth rough paths to \cite{bellingeri2022smooth}. The results on two-parameter controlled paths used here are taken primarily from \cite{cass2023fubini}.

\subsection{Review of one-parameter rough path theory}
Let $U, V$, $W$, and $F$ be real Banach spaces, with $U, V$ and $W$ finite-dimensional.  
We denote by $\mathrm{Hom}(U,F)$ the space of bounded linear maps from $U$ to $F$ equipped with the operator norm, and by $C([a,b],F)$ the Banach space of continuous functions $[a,b] \to F$ equipped with the topology of uniform convergence.  

For a function $f:[a,b]\to F$, we write the increment between $s$ and $t$ (with $a \le s \le t \le b$) as
\[
f_{s,t} := f_t - f_s.
\]

Finally, we denote by $\Delta_{[a,b]}$ the 2-simplex
\[
\Delta_{[a,b]} := \{ (s,t) \in [a,b]^2 : s \le t \}.
\]

\begin{definition}[Control and $\omega$-$p$-variation]
A control on $[a,b]$ is a continuous, super-additive map
$\omega:\Delta_{[a,b]}\to[0,\infty)$ vanishing on the diagonal.

Given $p\ge 1$ and $f:\Delta_{[a,b]}\to F$ set
\[
\|f\|_{p,\omega}
:=
\inf\Big\{C>0:\ \|f_{s,t}\|_F\le C\,\omega(s,t)^{1/p}
\ \text{for all }(s,t)\in\Delta_{[a,b]}\Big\},
\]
with the convention $\inf\varnothing=\infty$, and define
\[
C_\omega^p(\Delta_{[a,b]},F)
:=
\{f:\Delta_{[a,b]}\to F:\ \|f\|_{p,\omega}<\infty\}.
\]

For a path $g:[a,b]\to F$, we write $g\in C_\omega^p([a,b],F)$ if its increment map
\(
(s,t)\mapsto g_{s,t}
\)
belongs to $C_\omega^p(\Delta_{[a,b]},F)$.
\end{definition}

When the control $\omega$ is clear from context, we may omit it from the notation. We also use $C^{0,p}_\omega([a,b],V)$ to denote the closure of $C^1_\omega([a,b],V)$ in $C^p_\omega([a,b],V)$.\newline  

Denote by $U^{\otimes k}$ the 
k-fold tensor product of 
$U$, with the convention that $U^{\otimes 0} = \mathbb{R}$. We endow this spaces with the projective tensor norm.
Let $T((U))$ be the tensor algebra over $V$ with unit $\mathbf{1}$ and let $T^{\leq N}(U)$ its truncation at level $N$, we endow this last space with the $\ell^1$ norm. Let $G^N(U)$ denote the step-$N$ nilpotent group over $U$ equipped with the Carnot–Carathéodory metric, and let $\mathcal{L}^N(U)$ denote the corresponding Lie algebra.

\begin{definition}[Geometric rough paths]
Let $p\ge 1$ and let $\omega$ be a control. The set of geometric $p$-rough paths
$\Omega G_{p,\omega}([a,b],V)$ consists of all maps
\(
\mathbf{X}\in C^{0,p}_\omega\!\left(\Delta_{[a,b]},G^{\lfloor p\rfloor}(V)\right)
\)
that are multiplicative, i.e.
\[
\mathbf{X}_{s,s}=\mathbf{1}
\quad\text{and}\quad
\mathbf{X}_{s,u}=\mathbf{X}_{s,t}\otimes \mathbf{X}_{t,u},
\qquad \text{for all } a\le s\le t\le u\le b.
\]
\end{definition}

A special subclass of geometric rough paths is given by the smooth rough paths. 
These correspond to rough paths whose coordinate functionals are smooth in time, 
so that rough integrals reduce to classical integrals. This class will allow us 
to identify the signature kernel equation with a PDE and the Schwinger–Dyson 
kernel equation with an integro–differential equation.

\begin{definition}[Smooth rough paths]
We say that a rough path $\mathbf{X} \in \Omega G_{p,\omega}
([a,b],V)$ is a smooth rough path if for every
$f \in \mathrm{Hom}(T^{\le \lfloor p \rfloor}(V), \mathbb{R})$, the function
\[
t \mapsto f(\mathbf{X}_{a,t})
\]
belongs to $C^{1}_{\omega}([a,b], \mathbb{R})$.
\end{definition}

The key property linking smooth rough paths with PDE-type equations is that they admit a time derivative along the diagonal. More precisely, by Proposition~2.6 in \cite{bellingeri2022smooth}, for a smooth rough path $\mathbf{X}$ there exists an $\mathcal{L}^{\lfloor p \rfloor}(V)$-valued path $\boldsymbol{x}$ such that for every $s \in [a,b]$
\[
\boldsymbol{x}_s = \partial_t\big|_{t=s}\, \mathbf{X}_{s,t}.
\]
We refer to $\boldsymbol{x}$ as the diagonal derivate of $\mathbf{X}$.


For each, let \(\pi_k : T((U)) \to U^{\otimes k}\) denote the canonical projection onto the \(k\)-th level of the tensor algebra, and write \(\eta^{(k)} := \pi_k(\eta)\) for \(\eta \in T((U))\).  
We also denote by \(\eta^{\le k} :=\pi_{\le k} (\eta)\) (resp $\eta^{\ge k} = \pi_{\ge k}(\eta)$) the projection onto the direct sum of levels up to \(k\) (resp. levels greater than $k$). We use a similar notation for $\eta^{>k}$ (resp. $\eta^{<k}$).\newline 
Fix a basis $(e_i)_{i=1}^d$ of $U$ and let $(e^i)_{i=1}^d$ be the dual basis of $U^*$.
For a word $I=(i_1,\dots,i_n)$ with $0\le n\le N$, set
\[
e_I := e_{i_1}\otimes\cdots\otimes e_{i_n}\in U^{\otimes n},
\qquad
e^I := e^{i_1}\otimes\cdots\otimes e^{i_n}\in (U^*)^{\otimes n}.
\]
Let $\pi_n:T^{\le N}(U)\to U^{\otimes n}$ denote the grade-$n$ projection.
For $\eta\in T^{\le N}(U)$ we define its $(I)$-coordinate by
\[
\eta^{(I)} := \langle e^I,\ \pi_n \eta\rangle
= (e^{i_1}\otimes\cdots\otimes e^{i_n})(\pi_n\eta),
\qquad |I|=n.
\]
(For the empty word $\emptyset$, set $\eta^{(\emptyset)}:=\pi_0\eta\in\mathbb{R}$.)

\begin{definition}[$\mathbf{X}$-controlled paths]
Let $p \in [1,\infty)$ with $\kappa = \lfloor p \rfloor$, and let $\mathbf{X} \in \Omega G_{p, \omega}([a,b], V)$.  

Let
\[
Z : [a,b] \to \mathrm{Hom}(T^{<\kappa}(V), F),
\]
and suppose there exists 
\[
R : \Delta_{[a,b]} \to \mathrm{Hom}(T^{<\kappa}(V), F)
\]
such that
\begin{equation}\label{1D_controlled}
Z_{t} = Z_{s} \circ L_{\mathbf{X}_{s,t}^{<\kappa}} + R_{s,t}, \quad a \le s \le t \le b,
\end{equation}
where, for every $\eta \in T^{< \kappa}(V)$, 
\(
L_\eta : T((V)) \to T((V))
\) is 
the left tensor multiplication operator, defined by
\(
L_\eta(\xi) := \eta \otimes \xi.
\) 
Then we say that $Z$ is a  $\mathbf{X}$-controlled path if there exists a finite constant $C$ such that for all $a \le s \le t \le b$, any $\eta_j \in V^{\otimes j}$, $j \in [0:\kappa)$,
\begin{equation} \label{regularuty_constraint_remainder_1d}
\| R_{s,t}[\eta_j] \| \le C \, \omega(s,t)^{(\kappa - j)/p} \, \| \eta_j \|.
\end{equation}
\end{definition}
Equipped with the norm 
\[
\|Z\|_{\mathcal{D}^p_{\mathbf{X}}} = \|Z_a\| + \sum_{j=0}^k\|R^{(j)}\|_{\frac{p}{\floor{p}-k}, \omega_{\mathbf{X}}},
\]
this space becomes is a Banach space, denoted by $\mathcal{D}^p_{\mathbf{X}}([a,b], F)$.\newline

In what follows we will refer to $Z^{(0)}$ as the trace process and $R$ as the remainder process.\newline

Having introduced rough paths and controlled paths, we now turn to the construction of the rough integral.  
Given a rough path $\mathbf{X}$ defined as above and a path $Z$ controlled by $\mathbf{X}$, the rough integral of $Z$ against $\mathbf{X}$ generalizes the classical Young integral to allow for low-regularity integrators.  

A central tool in the construction is the map
\[
\Omega^Z : \Delta_{[a,b]} \to F, 
\qquad 
\Omega^Z(s,t) := Z_s\big[\mathbf{X}_{s,t}\big],
\]
associated with \(Z \in \mathcal{D}^p_{\mathbf{X}}([a,b], \mathrm{Hom}(V,F))\). This map captures the local increment structure of the integral. In fact, the rough integral is defined as the limit over a sequence of partitions of $[a,b]$ whose mesh converges to \(0\):
\[
\int_a^b Z_s \, d\mathbf{X}_s
:= \lim_{|\mathcal P| \to 0} \sum_{s_i\in \mathcal P} \Omega^Z(s_i,s_{i+1}),
\]
where the limit exists by the rough sewing argument (presented in Section~2 of \cite{cass2022combinatorial}).
The rough integral itself can be lifted to a controlled path $I_{\mathbf{X}}(Z) \in \mathcal{D}^p_{\mathbf{X}}([a,b], F)$ with $I^{(0)}_{\mathbf{X}}(Z) = \int Z d\mathbf{X}$ and $I^{(k)}_{\mathbf{X}}(Z) = Z^{(k-1)}$. Here we implicitly use the canonical identification
\[
\mathrm{Hom}\left(V^{\otimes (k-1)},\mathrm{Hom}(V,F)\right)\cong \mathrm{Hom}\left(V^{\otimes k},F\right),
\qquad
f\mapsto \left[(v_1\otimes\cdots\otimes v_k)\mapsto f(v_1\otimes\cdots\otimes v_{k-1})(v_k)\right].
\]
This identification will be used throughout the remainder of the paper. 

\subsection{Review of two-parameter rough path theory}

From the one-parameter rough path theory reviewed above, it is clear that, in order to define a notion of two-parameter (or joint) rough integration suitable for our rough kernel equations, one must first introduce an appropriate class of two-parameter controlled rough paths.

Intuitively, a two-parameter controlled rough path $Z$ is jointly controlled by the rough paths $\mathbf{X}$ and $\mathbf{Y}$ if it is controlled in one parameter by $\mathbf{X}$ and in the other by $\mathbf{Y}$. The next definition, which adapts the one in \cite{cass2023fubini},
makes this intuition precise.
\begin{definition}
[Two-parameter $\left(  \mathbf{X,Y}\right)  $-controlled paths] Let $\omega_\mathbf{X}$, $\omega_\mathbf{Y}$ be controls defined on $[a,b]$ and $[c,d]$ respectively. Let  $p$ and $q$
be in $[1,\infty)$ with $\kappa=\left\lfloor p\right\rfloor $ and
$\lambda=\left\lfloor q\right\rfloor $ and assume that $\mathbf{X} \in \Omega G_{p, \omega_{\mathbf{X}}}([a,b], V)$, $\mathbf{Y} \in \Omega G_{q, \omega_{\mathbf{Y}}}([c,d], W)$. Denote by $\boxtimes_{\mathrm{ext}}$ the external tensor product, and write
$\boxtimes$ for $\boxtimes_{\mathrm{ext}}$ in the sequel.
Let
\[
Z:\left[  a,b\right]  \times\left[  c,d\right]  \rightarrow\mathrm{Hom}\left(
T^{<\kappa}\left(  V\right)  \boxtimes_{\mathrm{ext}} T^{<\lambda}\left(  W\right)
,F\right)
\]
and suppose
\begin{align*}
R^{1}  & :\Delta_{\left[  a,b\right]  }\times\left[  c,d\right]
\rightarrow \mathrm{Hom}\left(  T^{<\kappa}\left(  V\right)  \boxtimes_{\mathrm{ext}}
T^{<\lambda}\left(  W\right)  ,F\right)  \text{ and }\\
R^{2}  & :\left[  a,b\right]  \times\Delta_{\left[  c,d\right]  }%
\rightarrow \mathrm{Hom}\left(  T^{<\kappa}\left(  V\right)  \boxtimes_{\mathrm{ext}}
T^{<\lambda}\left(  W\right)  ,F\right)
\end{align*}
are defined respectively by the identities
\begin{align}
Z_{t,u}  & = Z_{s,u}\circ\left[  L_{\mathbf{X}_{s,t}^{<\kappa}}\boxtimes_{\mathrm{ext}}
\text{id}_{T^{<\lambda}\left(  W\right)  }\right]  +R_{s,t}^{1}\left(
u\right)  \text{ and}\nonumber\\
Z_{s,v}  & =Z_{s,u}\circ\left[  \text{id}_{T^{<\kappa}\left(  V\right)
}\boxtimes_{\mathrm{ext}} L_{\mathbf{Y}_{u,v}^{<\lambda}}\right]  +R_{u,v}^{2}\left(
s\right)  .\label{constraints}%
\end{align}
Then we say that $Z$ is a two-parameter $\left(  \mathbf{X,Y}\right)
$-controlled path if there exist finite constants $C_{1}$ and $C_{2}$ such
that for all $a\leq s\leq t\leq b$ and $c\leq u\leq v\leq d$ and any $\eta\in
T^{<\kappa}\left(  V\right)  $ and $\xi\in T^{<\lambda}\left(  W\right)  $ we
have
\begin{equation}\label{regularity_condition_r1}
\left \| R_{s,t}^{1}\left(  u\right)  \left[  \eta_{j}\boxtimes_{\mathrm{ext}}
\xi \right]  \right \| \leq C_{1}\omega_{\mathbf{X}}\left(
s,t\right)  ^{\left(  \kappa-j\right)  /p}\left \| \eta_{j}%
\right \| \| \xi \|
\text{, }
\end{equation}
for all $\eta_{j}\in V^{\otimes j}$, $j\in\lbrack0:\kappa)$ and if also
\[
\left\| R_{u,v}^{2}\left(  s\right)  \left[  \eta \boxtimes_{\mathrm{ext}}
\xi _{k}\right]  \right\| \leq C_{2}\omega_{\mathbf{Y}}\left(
u,v\right)  ^{\left(  \lambda-k\right)  /q} \left\| \eta \right\| \left\| \xi_{k}\right\|,
\]
for all $\xi_{k}\in W^{\otimes k}$, $k\in\lbrack0:\lambda).$
\end{definition}
The external tensor product
\(
T^{<\kappa}(V)\boxtimes T^{<\lambda}(W)
\)
is used mainly as a bookkeeping device, keeping the \(V\)- and \(W\)-directions distinct. This is particularly convenient since \(T^{<\kappa}(V)\) and \(T^{<\lambda}(W)\) already carry their own internal tensor-algebra structures. Moreover, after choosing bases in \(T^{<\kappa}(V)\) and \(T^{<\lambda}(W)\), every element admits a unique expansion in the induced tensor-product basis, so that components may be indexed separately in the two directions.\newline
In what follows, for any $\eta\in T((V))\boxtimes T((W))$, we write
\(
\eta^{(j,k)} := (\pi_j \boxtimes \pi_k)(\eta),
\)
where $\pi_j$ and $\pi_k$ denote the canonical projections defined above.

The notion of two-parameter rough paths was first introduced in \cite{gerasimovivcs2019hormander}
for the case $p\in\lbrack2,3).$ The paper \cite{cass2023fubini} extended it to the case
\thinspace$p\geq3$ and $\mathbf{X}$ possibly different than $\mathbf{Y}$. The definitions adopted there differ slightly from those
we use. In particular, we have broadened the notion to allow for the
controlling rough paths to be different in the respective parameters of $Z.$
An important difference is that we enforce symmetry of the partial derivative
processes; the following remark captures what this means.

\begin{remark}
Even in the case where these path are the same ($\mathbf{X=Y,}$ in the
definition above) the two notions differ. To illustrate this we write
out the conditions (\ref{constraints}) in full for $p\in\lbrack2,3).$ The
first can be expanded as%
\begin{align*}
Z_{t,u}^{\left(  0,0\right)  }  & =Z_{s,u}^{\left(  0,0\right)  }%
+Z_{s,u}^{\left(  1,0\right)  }\left[  \mathbf{X}^{(1)}_{s,t}\boxtimes id_{T^{<2}(V)} \right]
+R_{s,t}^{1;\left(  0,0\right)  }\left(  u\right) \\
Z_{t,u}^{\left(  1,0\right)  }  & =Z_{s,u}^{\left(  1,0\right)  }%
+R_{s,t}^{1;\left(  1,0\right)  }\left(  u\right)  ,\\
Z_{t,u}^{\left(  0,1\right)  }  & =Z_{s,u}^{\left(  0,1\right)  }%
+Z_{s,u}^{\left(  1,1\right)  }\circ\left[  L_{\mathbf{X}^{< 2}_{s,t}}\boxtimes id_{T^{<2}(V)} \right]  +R_{s,t}^{1;\left(  0,1\right)  }\left(  u\right)  ,
\end{align*}
while the second becomes
\begin{align*}
Z_{s,v}^{\left(  0,0\right)  }  & =Z_{s,u}^{\left(  0,0\right)  }%
+Z_{s,u}^{\left(  1,0\right)  }\left[id_{T^{<2}(V)} \boxtimes \mathbf{X}^{(1)}_{s,t}  \right]
+R_{u,v}^{2;\left(  0,0\right)  }\left(  s\right) \\
Z_{s,v}^{\left(  0,1\right)  }  & =Z_{s,u}^{\left(  0,1\right)  }%
+R_{u,v}^{2;\left(  0,1\right)  }\left(  s\right)  ,\\
Z_{s,v}^{\left(  1,0\right)  }  & =Z_{s,u}^{\left(  1,0\right)  }%
+Z_{s,u}^{\left(  1,1\right)  }\circ\left[  id_{T^{<2}(V)} \boxtimes L_{\mathbf{X}^{< 2}_{u,v}%
}\right]  +R_{s,t}^{2;\left(  0,1\right)  }\left(  u\right)  .
\end{align*}
Our definition therefore enforces that \(Z_{s,\cdot}^{(1,1)}\) appears in the expansion of \(Z_{s,\cdot}^{(1,0)}\) and, simultaneously, that \(Z_{\cdot,u}^{(1,1)}\) appears in the expansion of \(Z_{\cdot,u}^{(0,1)}\). In other words, \(Z^{(1,1)}\) is the same second-order term appearing in the iterated expansions, irrespective of the order in which increments are taken in the two parameters. In the framework of \cite{gerasimovivcs2019hormander} on the other hand, the order
of these two operations is distinguished with their commutativity
subsequently imposed as a condition for the rough Fubini theorems. As our
applications all rely on being able to use these theorems it is helpful
apriori to build it into the definition of a two-parameter controlled path.
This also allows for the succinct condition (\ref{constraints}) to be used and
allows a notationally lighter treatment of the low-regularity regime.
\end{remark}
\begin{comment}
To do this, we begin by equipping each tensor power $V^{\otimes j}$ and $W^{\otimes k}$ with a fixed admissible tensor norm, the truncated tensor algebras
\(
T^{<\kappa}(V) \,,\, 
T^{<\lambda}(W)
\)
inherit the corresponding direct-sum norms
\[
\|\eta\|_{T^{<\kappa}(V)} := \sum_{j=0}^{\kappa-1} \|\eta_j\|_{V^{\otimes j}}, \qquad
\|\xi\|_{T^{<\lambda}(W)} := \sum_{k=0}^{\lambda-1} \|\xi_k\|_{W^{\otimes k}}.
\]
With these choices, the processes $R^1$ and $R^2$, which we refer to as remainders or first-level remainders, naturally belong to the spaces
\[
R^1 \in C^p_{\omega_{\mathbf{X}}}\Big([a,b], C^0\big([c,d], \mathrm{Hom}(T^{<\kappa}(V)\boxtimes T^{<\lambda}(W), F)\big)\Big),
\]
\[
R^2 \in C^q_{\omega_{\mathbf{Y}}}\Big([c,d], C^0\big([a,b], \mathrm{Hom}(T^{<\kappa}(V)\boxtimes T^{<\lambda}(W), F)\big)\Big).
\]
\end{comment}
An immediate consequence of the definition above (Lemma 3.3 in \cite{cass2023fubini}) is that the processes $R^1$ and $R^2$ are themselves respectively $\mathbf{Y}$
and $\mathbf{X}$-controlled rough paths and that there exists a map 
\[
\mathbf{R}: \Delta_{[a,b]} \times \Delta_{[c,d]} \rightarrow \mathrm{Hom}\left(
T^{<\kappa}\left(  V\right)  \boxtimes T^{<\lambda}\left(  W\right)
,F\right) 
\]
such that 
\begin{align*}
\mathbf{R} \begin{pmatrix} s,t \\ u,v \end{pmatrix} &:= R^1_{s,t}(v) - R^1_{s,t}(u) \circ\left[ \text{id}_{T^{<\lambda}\left(  V\right)}  \boxtimes  L_{\mathbf{Y}_{u,v}^{<\kappa}
 }\right] \\
    &= R^2_{u,v}(t) - R^2_{u,v}(s) \circ\left[    L_{\mathbf{X}_{s,t}^{<\kappa}}\boxtimes  \text{id}_{T^{<\lambda}\left(  W\right)
 }\right]
\end{align*}
and for which there exists a constant $C$, such that for all $a \le s \le t \le b$ and $c \le u \le v \le d$ and any $\eta\in
T^{<\kappa}\left(  V\right)  $ and $\xi\in T^{<\lambda}\left(  W\right)  $ we
have
\begin{equation}\label{remainder_remainder_norm}
\left\| \mathbf{R} \begin{pmatrix} s,t \\ u,v \end{pmatrix} [\eta_j \boxtimes \xi_k] \right\| \leq C \omega_{\mathbf{X}}\left(
s,t\right)  ^{\left(  \kappa-j\right)  /p} \omega_{\mathbf{Y}}\left(
u,v\right)  ^{\left(  \lambda-k\right)  /q}\left \| \eta_{j}%
\right \| \| \xi_k \|,
\end{equation}

for all $\eta_{j}\in V^{\otimes j}$, $j\in\lbrack0:\kappa)$ and  $\xi_{k}\in W^{\otimes k}$, $k\in\lbrack0:\lambda).$
We refer to $\mathbf{R}$ as second level remainder or remainder of the remainder to highlight the dependence above. The estimate \eqref{remainder_remainder_norm} suggests that the level-wise second order remainders can be viewed as a two-parameter map with finite $p$-variation in the first parameter and finite $q$-variation in the second. This is captured by the following definition.  
\begin{definition}[Joint $(p, q)$-variation norm]
Let $\omega_1, \omega_2$ be a controls on $[a,b]$ and $[c,d]$ and let $p,q \ge 1$. For a map $f : \Delta_{[a,b]} \times \Delta_{[c,d]} \to F$, define
\[
\|f\|_{(p,q),(\omega_1, \omega_2)} := \inf \Big\{ C > 0 : \left \|f{\begin{pmatrix}
    s, t \\ u,  v
\end{pmatrix}}\right\|_F \le C \, \omega_1(s,t)^{1/p} \omega_2(u,v)^{1/q} \ \text{for all } a \le s \le t \le b\, , \, c \le u \le v \le d  \Big\},
\]
with the convention $\inf\varnothing=\infty$.
\end{definition}


Within this framework, we introduce the following norm on the space of two-parameter $(\mathbf{X},\mathbf{Y})$-controlled paths, defined by 
\[
\begin{aligned}
\|Z\|_{\mathcal{D}_{\mathbf{X}, \mathbf{Y}}}
:= \big| Z_{a,c}  \big| + \sum_{j = 0}^{\kappa -1} \sum_{k = 0}^{\lambda - 1} \Big[
&\;
  \big\| R^{1; (j,k)}(c)
       \big\|_{\frac{p}{\kappa - j}, \omega_\mathbf{X}} + \big\| R^{2; (j,k)}(a)   \big\|_{ \frac{q}{\lambda - k}, \omega_{\mathbf{Y}}}
 + \big\| \mathbf{R}^{(j,k)} \big\|_{(\frac{p}{\kappa - j},\frac{q}{\lambda - k}), (\omega_\mathbf{X}, \omega_\mathbf{Y})} 
\Big],
\end{aligned}
\]
Additionally, for fixed $\tilde{\mathbf{X}} \in  \Omega G_{p, \omega_{\mathbf{X}}}([a,b], V)$, $\tilde{\mathbf{Y}} \in  \Omega G_{q, \omega_{\mathbf{Y}}}([c,d], W)$ and a $(\tilde{\mathbf{X}}, \tilde{\mathbf{Y}})$-controlled path $\tilde{Z}$, it is possible to define the \enquote{distance}
\[
\begin{aligned}
d^{(\mathbf{X}, \mathbf{Y}), (\tilde{\mathbf{X}}, \tilde{\mathbf{Y}})}_{[a,b] \times [c,d]}(Z, \tilde{Z})
:= 
\;\big| Z_{a,c} - \tilde{Z}_{a,c} \big| + \sum_{j = 0}^{\kappa -1} &\sum_{k = 0}^{\lambda - 1} \bigg[  \big\| R^{Z, 1; (j,k)}(c)
      - R^{\tilde{Z}, 1; (j,k)}(c)
       \big\|_{\frac{p}{\kappa - j}, \omega_\mathbf{X}} \\
&\quad
 + \big\| R^{Z, 2; (j,k)}(a)
      - R^{\tilde{Z}, 2; (j,k)}(a)
      \big\|_{\frac{q}{\lambda - k}, \omega_\mathbf{Y}} \\
&\quad
 + \big\| \mathbf{R}^{Z, (j,k)} - \mathbf{R}^{\tilde{Z}, (j,k)} \big\|_{(\frac{p}{\kappa - j},\,\frac{q}{\lambda - k}), (\omega_\mathbf{X}, \omega_\mathbf{Y})}
\bigg].
\end{aligned}
\]
We remark that the above is not a true metric as $Z$ and $\tilde{Z}$ can possibly belong to different spaces. \newline
Just like we mentioned above, when the choice of controls is clear from the context, to improve readability, we will omit the control from the norm and remove the sub and superscripts from the \enquote{distance}.\newline
In the two-parameter setting, this norm necessarily involves the second-order remainder.
Taking increments in both parameters reveals cross-parameter dependence of the first-level remainders, which produces a second-order term.
Hence controlling the increments of the two-parameter integral will require to control of the second-order remainder.
The space of two-parameters controlled rough paths equipped with the norm will be denoted as $\mathcal{D}^{p, q}_{\mathbf{X}, \mathbf{Y}}([a,b] \times [c,d], F)$ throughout this work. Whenever $\mathbf{X} = \mathbf{Y}$ we will simplify the notation by writing $\mathcal{D}^p_{\mathbf{X}}([a,b] \times [a,b], F)$.\newline

As we anticipated before, this class of two-parameter paths is defined so that their interactions with the rough paths $\mathbf{X}$ and $\mathbf{Y}$ are sufficiently regular to allow one to define a rough integral over the rectangle $[a,b]\times[c,d]$. We recall that this joint integral is precisely the one that appears in the rough formulation of the signature-kernel trick. The construction of the integral, and the main properties we use, are summarised below from \cite{cass2023fubini}.\newline
To construct such an integral, one needs an analogue of the map $\Omega$ now familiar from the one-parameter setting. 
In this case, the role is played by the map  
\[
\Omega^Z: [a,b] \times [c,d] \to F, \qquad \Omega^Z \begin{pmatrix} s,t \\ u,v \end{pmatrix} := 
Z_{s,u} \big[ \mathbf{X}^{\geq 1}_{s,t} \boxtimes \mathbf{Y}^{\geq 1}_{u,v} \big],
\]  
which encodes how the integrand \(Z\) interacts with the structure of the driving rough paths \(\mathbf{X}\) and \(\mathbf{Y}\). 
Notice that if one of the increments $s,t$ or $u,v$ is fixed, we recover the classical, one-parameter, approximating map acting on either \(\mathbf{X}_{s,t}\) or \(\mathbf{Y}_{u,v}\).

Just like in the one-parameter setting, the next step is to define an approximation of the two-parameter rough integral. 
To this end, fix the two partitions:  
\[
\mathcal{P} = \{ a = s_1 \leq s_2 \leq \dots \leq s_n = b \}, 
\qquad  
\mathcal{D} = \{ c = u_1 \leq u_2 \leq \dots \leq u_m = d \},
\]  
and define the discrete two-parameter integral  
\begin{equation}\label{two_param_sewing}
\sum_{\mathcal{P}\times \mathcal{D}} \Omega^Z 
   := \sum_{i=1}^{n-1} \sum_{j=1}^{m-1} 
   \Omega^Z \begin{pmatrix} s_{i+1}, s_i \\ u_{j+1}, u_j \end{pmatrix}.
\end{equation}  
The double sum represents the accumulation of the elementary contributions of \(\Omega^Z\) over each rectangle \([s_i,s_{i+1}]\times[u_j,u_{j+1}]\). 

One of the main results in \cite{cass2023fubini} shows that, as the mesh of the rectangular partition tends to zero, these sums converge to a continuous and additive map on the rectangle $[a,b] \times [c,d]$, which we call two-parameter (or joint) rough integral
\[
\int_{[a,b] \times [c,d]} Z_{s, u} \, d(\mathbf{X}, \mathbf{Y})  
   :=  \lim_{\left|\mathcal{P} \times \mathcal{D}\right| \to 0} 
   \sum_{\mathcal{P}\times \mathcal{D}} \Omega^Z.
\]
Moreover, this integral depends continuously on the pair of rough paths and on the controlled integrand. For convenience, we record in Appendix~\ref{appendix_technical} the precise existence estimate and the corresponding stability bound in the notation used here; see Theorems~\ref{existence_2D_theorem} and~\ref{stability_theorem_2D}.
Additionally, the structure of the two-parameter controlled rough paths allows the the two-parameter rough integral to agree with the iterated one-parameter integrals. Concretely,
using the fact that two-parameter controlled paths are controlled in both parameters to define the maps 
\begin{align*}
&Z \in \mathcal{D}^q_{\mathbf{Y}}([c,d], \mathcal{D}^p_{\mathbf{X}}([a,b], F)) \to \int_a^\cdot Z_{s,u} \mathbf{Y}_u \in \mathcal{D}^p_{\mathbf{X}}([a,b], F) \text{ and } \\
&Z \in \mathcal{D}^p_{\mathbf{X}}([a,b], \mathcal{D}^q_{\mathbf{Y}}([c,d], F)) \to \int_c^\cdot Z_{s,u} \mathbf{X}_s \in \mathcal{D}^q_{\mathbf{Y}}([c,d], F)
\end{align*}
the following identity, derived in Theorem 5.2 in \cite{cass2023fubini}, holds:
\[
\int_{[a,b] \times [c,d]} Z_{s, u} \, d(\mathbf{X}_s, \mathbf{Y}_u)  
= \int_{a}^b \int_{c}^d Z_{s, u} \, d\mathbf{Y}_u \, d\mathbf{X}_s 
= \int_c^d \int_a^b Z_{s, u} \, d\mathbf{X}_s \, d\mathbf{Y}_u.
\]
This property, which is referred to as rough Fubini theorem, allows to define a controlled path  $I_{\mathbf{X}, \mathbf{Y}}(Z) \in \mathcal{D}^{p, q}_{\mathbf{X}, \mathbf{Y}}\big([a, b] \times [c,d],  F\big)$ such that  for $j \in [1:\kappa), k\in [1:\lambda)$ and $\eta \in V^{\otimes j}$, $\xi \in W^{\otimes k}$
\begin{equation}\label{definition_I_X_Y}
\begin{aligned}
    I_{\mathbf{X}, \mathbf{Y}}(Z)_{t,v}[\mathbf{1}, \mathbf{1}] &=  \int_{[a,t] \times [c, v]} Z_{s,u} d(\mathbf{X}_s, \mathbf{Y}_u),\\
    I_{\mathbf{X}, \mathbf{Y}}(Z)_{t,v} [\eta,  \mathbf{1}] &= \int_c^v Z_{t,u} d\mathbf{Y}_u [\eta],  \\
    I_{\mathbf{X}, \mathbf{Y}}(Z)_{t,v} [ \mathbf{1}, \xi] &= \int_a^t Z_{s,v} d\mathbf{X}_s [\xi],\\
    I_{\mathbf{X}, \mathbf{Y}}(Z)_{t,v} [\eta, \xi] &=  Z_{t,v} [\eta, \xi].
\end{aligned}
\end{equation}
That two-parameter rough integrals produce two-parameter controlled paths is consistent with the fact that, as we anticipated, solutions to two-parameter rough differential equations are themselves two-parameter controlled paths. This mirrors the familiar one-parameter theory, where rough integration maps controlled paths to controlled paths, and solutions of rough differential equations are (one-parameter) controlled paths.




\section{Rough integrals on two-dimensional simplices} \label{Simplex_section}
In the previous section, we introduced a notion of joint rough integration that makes sense of the integral appearing in the signature-kernel trick. However, this construction does not directly yield a rough integral over a triangular domain, as required for the Schwinger--Dyson kernel. The obstruction is that, in the map \eqref{two_param_sewing}, the partitions in the two directions are chosen independently, and therefore do not capture the geometry of a triangle.\newline
To carry out the programme outlined above, we need to extend the two-parameter integral from rectangles to two-dimensional simplices for general two-parameter controlled rough paths. In the setting of \cite{gerasimovivcs2019hormander}, a joint integral over a triangular domain is constructed for two-parameter controlled paths admitting smooth approximations and controlled by a geometric $p$-rough path with $p\in[2,3)$. Here we broaden this result to arbitrary $p\ge 1$ and to general two-parameter controlled paths, by leveraging a rough Fubini-type theorem.

More precisely, we first show that the order of integration over a simplex can be reversed. We then prove that, given any two-parameter controlled path $Z$, one can construct a controlled rough path $\widehat{Z}$ by symmetrising $Z$ along the diagonal. This symmetrised object allows us to relate integrals over triangular domains to integrals over rectangles.
 

Throughout this section we work with a \(p\)-geometric rough path \(\mathbf{X}\), controlled by \(\omega_{\mathbf{X}}\), and a two-parameter controlled path
\[
Z \in \mathcal{D}^{p}_{\mathbf{X}}\!\left( \Delta_{[a,b]} , \operatorname{Hom}(V \boxtimes V, F)\right).
\]
For notational convenience, set \(\kappa := \lfloor p\rfloor\).

Our first objective is to study the iterated rough integrals 
\begin{equation}\label{iterated_1_d_integrals}
\int_a^b \int_a^u Z_{s, u}\, d\mathbf{X}_s \, d\mathbf{X}_u,
\qquad 
\int_a^b \int_s^b Z_{s, u}\, d\mathbf{X}_u \, d\mathbf{X}_s, 
\end{equation}
with the aim of constructing the corresponding outer integrals. For this, it is necessary to ensure that the inner integrals can be made into 
one-parameter controlled rough paths. To this end, it is convenient to introduce the auxiliary maps
\begin{align*}
Z^L &: [a,b] \to 
\operatorname{Hom}\!\left(
T^{< \kappa}(V),
\operatorname{Hom}(V, F)
\right), \\
Z^U &: [a,b] \to 
\operatorname{Hom}\!\left(
T^{< \kappa}(V),
\operatorname{Hom}(V, F)
\right),
\end{align*}
defined for \(t \in [a,b]\) by
\begin{align*}
Z^L_t[\eta](\cdot) &= Z_{t,t} \left[\Delta \eta\ \otimes ( \mathbf{1} \boxtimes \cdot)\right], \\
Z^U_t[\eta](\cdot) &= Z_{t,t} \left[\Delta \eta\ \otimes (\cdot \boxtimes  \mathbf{1})\right].
\end{align*}
Here \(\Delta\) denotes the shuffle coproduct on \(T^{<\kappa}(V)\), i.e. the unique algebra morphism
\begin{align*}
&\Delta:T^{<\kappa}(V)\to T^{<\kappa}(V)\boxtimes T^{<\kappa}(V),\\
&\Delta( \mathbf{1})= \mathbf{1}\boxtimes  \mathbf{1},\quad \Delta(\eta)=\eta\boxtimes  \mathbf{1}+ \mathbf{1}\boxtimes \eta, \ \eta\in V,
\end{align*}
see Section 2 of \cite{cass2022combinatorial} for more details.
Equivalently, for every non-empty word \(I=i_1\cdots i_n\) over the alphabet \(\{1,\dots,d\}\),
we may define \(\Delta\) by
\begin{align*}
&\Delta(\emptyset) = \emptyset \boxtimes \emptyset,\\\
&\Delta(e_{I})
=
\prod_{m=1}^n \bigl(e_{i_m}\boxtimes \emptyset +\emptyset\boxtimes e_{i_m}\bigr).    
\end{align*}


The next proposition shows how these maps arise naturally when identifying the one-parameter controlled path above the inner integrals in \eqref{iterated_1_d_integrals}.

\begin{lemma}\label{proposition_L_U_controlled}
Let $\mathbf{X}$, $Z$, $Z^L$ and $Z^U$ be defined as above, then the following paths belong to the space $\mathcal{D}^p_{\mathbf{X}}([a,b], \operatorname{Hom}(V, F))$:
\[
L_s := \int_a^s Z_{u,s} \, d\mathbf{X}_u + Z^L_{s}, 
\qquad
U_u := \int_u^b Z_{u,s} \, d\mathbf{X}_s - Z^U_{u}.  
\]
\end{lemma}
\begin{proof}
Fix $a \le s \le t \le b$. The first process admits the representation
\begin{equation} \label{LT_expansion}
L_t = \int_a^t Z_{u,t}\, d\mathbf{X}_u + Z_t^L .
\end{equation}
We begin by expanding the rough integral appearing on the right-hand side.
Using the decomposition of $Z_{u,t}$, we obtain
\begin{align*}
\int_a^t Z_{u,t}\, d\mathbf{X}_u
&= \int_a^t Z_{u,s} \circ
\left[ id_{T^{<\kappa}(V)} \boxtimes L_{\mathbf{X}_{s,t}^{<\kappa}} \right]
\, d\mathbf{X}_u
+ \int_a^t R^2_{s,t}(u)\, d\mathbf{X}_u \\
&= \int_s^t Z_{u,s} \circ
\left[ id_{T^{<\kappa}(V)} \boxtimes L_{\mathbf{X}_{s,t}^{<\kappa}} \right]
\, d\mathbf{X}_u
+ \int_a^s Z_{u,s} \circ
\left[ id_{T^{<\kappa}(V)} \boxtimes L_{\mathbf{X}_{s,t}^{<\kappa}} \right]
\, d\mathbf{X}_u  \\
&\quad + \int_a^t R^2_{s,t}(u)\, d\mathbf{X}_u .
\end{align*}
In the first integral of the last identity we expand $Z_{u,s}$ in the first parameter to get
\begin{align*}
\int_a^t Z_{u,t}\, d\mathbf{X}_u
&= Z_{s,s} \circ
\left[\mathbf{X}^{\ge 1}_{s,t} \boxtimes L_{\mathbf{X}_{s,t}^{<\kappa}}\right]
+ \int_s^t R^1_{s,u}(s) \circ
\left[ id_{T^{<\kappa}(V)} \boxtimes L_{\mathbf{X}_{s,t}^{<\kappa}} \right]
\, d\mathbf{X}_u \\
&\quad + \int_a^s Z_{u,s} \circ
\left[id_{T^{<\kappa}(V)} \boxtimes  L_{\mathbf{X}_{s,t}^{<\kappa} } \right]
\, d\mathbf{X}_u
+ \int_a^t R^2_{s,t}(u)\, d\mathbf{X}_u .
\end{align*}

Turning to the second term in \eqref{LT_expansion}, recall that
\begin{align*}
Z_t^L[\eta]
&= Z_{t,t}
\left[\Delta \eta \otimes (\mathbf{1} \boxtimes \cdot) \right] \\
&=
\left(
Z_{s,s} \circ
\left[ L_{\mathbf{X}_{s,t}^{<\kappa}} \boxtimes
L_{\mathbf{X}_{s,t}^{<\kappa}} \right]
+ R^2_{s,t}(s) \circ
\left[ L_{\mathbf{X}_{s,t}^{<\kappa}} \boxtimes id_{T^{<\kappa}(V)} \right]
+ R^1_{s,t}(t)
\right)
\left[\Delta \eta \otimes (\mathbf{1} \boxtimes \cdot) \right].
\end{align*}
Notice how $Z_{t,t}[\mathbf{1} \otimes \eta] =0$. This, combined with the expansions of the integral and $Z^L$, yields
\begin{align*}
L_t
= \int_a^s Z_{u,s} \circ
\left[ id_{T^{<\kappa}(V)} \boxtimes L_{\mathbf{X}_{s,t}^{<\kappa}} \right]
\, d\mathbf{X}_u
+ Z_{s}^L \circ
 L_{\mathbf{X}_{s,t}^{<\kappa}} 
+ R^L_{s,t}(t),
\end{align*}
where the remainder term $R^L$ is determined by the preceding expansions and is immediately seen to satisfy the regularity constraint \eqref{regularuty_constraint_remainder_1d}. This proves the claim for the first process.

The proof for $L^U$ follows by identical arguments and is therefore omitted.
\end{proof}

Having established that the two paths $L$ and $U$ belong to the same space, a natural next question is whether the corresponding rough integrals
\[
\int_a^b U_s \, d\mathbf{X}_s = \int_a^b \int_s^b Z_{s, u}\, d\mathbf{X}_u \, d\mathbf{X}_s
\quad \text{and} \quad
\int_a^b L_u \, d\mathbf{X}_u = \int_a^b \int_a^u Z_{s, u}\, d\mathbf{X}_s \, d\mathbf{X}_u
\]
coincide. In what follows, we show that this holds true.

To prepare for the proof, we introduce the \enquote{symmetrised} path 
\begin{align*}
&\widehat{Z}:\left[  a,b\right]  \times\left[  a,b\right]  \rightarrow\mathrm{Hom}\left(
T^{<\kappa}\left(  V\right)  \boxtimes T^{<\kappa}\left(  V\right)
,\mathrm{Hom}(V\boxtimes V, F)\right),\\ 
&\widehat{Z}_{s,t}\left[\eta, \xi\right] :=
\begin{cases}
Z_{s,t}\left[\eta, \xi\right], & \text{if } s \le t,\\
Z_{t,s}\left[\xi, \eta\right], & \text{if } s > t,
\end{cases}
\end{align*}
where, in the evaluation above, we use the canonical identification
\[
\mathrm{Hom}\left(
T^{<\kappa}\left(  V\right)  \boxtimes T^{<\kappa}\left(  V\right)
,\mathrm{Hom}(V\boxtimes V, F)\right)
\cong \mathrm{Hom}\left(
T^{\leq \kappa}\left(  V\right)  \boxtimes T^{\leq \kappa}\left(  V\right)
, F\right).
\] 
This construction symmetrises the controlled path $Z$ across the diagonal of
$[a,b]\times[a,b]$ and is instrumental in showing the relation  between the joint rough integral over the simplex and the joint rough integral over rectangular domains, in fact, in terms of $\hat{Z}$, the question above can be recast as
the validity of the following identity
\begin{equation}\label{identity_tringular_Fubini}
2 \int_a^b L_s \, d\mathbf{X}_s
= \int_{[a,b]\times[a,b]} \widehat{Z}_{s,u} \, d(\mathbf{X}_s,\mathbf{X}_u)
= 2 \int_a^b U_u \, d\mathbf{X}_u .
\end{equation}
The heuristic intuition behind the equality above 
comes from the classical Fubini-Tonelli theorem, for which, if $X$ is a path of bounded variation, and $Z$ is bounded and measurable, then 
\begin{align*}
&2 \int_a^b \int_s^b Z_{s,u} dX_u dX_s\\
&=  \int_a^b \int_s^b Z_{s,u} dX_u dX_s + \int_a^b \int_a^u Z_{s,u} dX_s dX_u \\
&= \int_a^b \int_a^b \left(Z_{s,u} \mathbbm{1}_{\{s \le u\}} + Z_{u,s} \mathbbm{1}_{\{s > u\}}\right) dX_u dX_s\\
&= \int_a^b \int_a^b \hat{Z}_{s,u} dX_u dX_s,  
\end{align*}
where in the third line we renamed the integration variables in the second $Z$.\newline
To address the rough case, we first need to verify that $\widehat{Z}$ belongs to the class
$\mathcal{D}^{p}_{\mathbf{X}}\!\left([a,b]\times[a,b],\, \mathrm{Hom}(V\boxtimes V, F)\right)$.
The next proposition ensures that this is the case.


\begin{proposition}
The process $ \hat{Z}$ belongs to the space  $\mathcal{D}^{p}_{\mathbf{X}}\left([a,b] \times [a,b], \mathrm{Hom}(V \boxtimes V, F)\right)$.
\end{proposition}
\begin{proof}
Consider the points $s,t \in \Delta_{[a,b]}$, $u \in [a,b]$.
In the case $u \geq t \geq s$
\begin{align*}
\hat{Z}_{t,u} &= Z_{t,u}\\
&= Z_{s,u}\circ\left[  L_{\mathbf{X}_{s,t}^{<\kappa}}\boxtimes
\text{id}_{T^{<\kappa}\left(  W\right)  }\right]  +R_{s,t}^{1}\left(
u\right)\\
&= \hat{Z}_{s,u}\circ\left[  L_{\mathbf{X}_{s,t}^{<\kappa}}\boxtimes
\text{id}_{T^{<\kappa}\left(  W\right)  }\right]  +R_{s,t}^{1}\left(
u\right),
\end{align*}
if $t \geq s > u$
\begin{align*}
\hat{Z}_{t,u} &= Z_{u,t}\\
&= Z_{u,s}\circ\left[ 
\text{id}_{T^{<\kappa}\left(  W\right)  } \boxtimes L_{\mathbf{X}_{s,t}^{<\kappa}}\right]  +\hat{R}_{s,t}^{2}\left(
u\right)\\
&= \hat{Z}_{s,u}\circ\left[  L_{\mathbf{X}_{s,t}^{<\kappa}}\boxtimes
\text{id}_{T^{<\kappa}\left(  W\right)  }\right]  +\hat{R}_{s,t}^{2}\left(
u\right),
\end{align*}
where, for every $\eta, \xi \in T^{< \kappa}(V)$, $\hat{R}^2[\eta, \xi] = R^2[ \xi, \eta]$.\newline
Assume now that $t \geq  u \geq  s$, we can write 
\begin{align*}
\hat{Z}_{t,u} &= \hat{Z}_{u,u} \circ  \left[ L_{\mathbf{X}^{< \kappa}_{u,t}} \boxtimes id_{T^{< \kappa}(V)} \right] + \hat{R}^2_{u,t}(u) \\
&= Z_{u,u} \circ  \left[ L_{\mathbf{X}^{< \kappa}_{u,t}} \boxtimes id_{T^{< \kappa}(V)} \right] + \hat{R}^2_{u,t}(u) \\
&= Z_{s,u} \circ  \left[ L_{\mathbf{X}^{< \kappa}_{s,u}} \boxtimes id_{T^{< \kappa}(V)} \right] \circ  \left[ L_{\mathbf{X}^{< \kappa}_{u,t}} \boxtimes id_{T^{< \kappa}(V)} \right] + R^1_{s,u}(u) + \hat{R}^2_{u,t}(u)\\
&= \hat{Z}_{s,u} \circ  \left[ L_{\mathbf{X}^{< \kappa}_{s,u}} \boxtimes id_{T^{< \kappa}(V)} \right] \circ  \left[ L_{\mathbf{X}^{< \kappa}_{u,t}} \boxtimes id_{T^{< \kappa}(V)} \right] + R^1_{s,u}(u) + \hat{R}^2_{u,t}(u).
\end{align*}

These expansions motivate the definition \[
R^{\hat{Z}, 1}_{s,t}(u) := \left(\hat{R}^2_{u,t}(u) +  R^{1}_{s,u}(u)\right) \mathbbm{1}_{t \geq u \geq s} + \hat{R}^2_{s,t}(u)\mathbbm{1}_{t \geq s >u} + R^1_{s,t}(u)  \mathbbm{1}_{u > t \geq s}.\]
Notice that this remainder satisfies the required regularity condition, this is seen from the fact that every element in the sum above satisfies the bound \eqref{regularity_condition_r1}.\newline
For the increments in the other parameter, we can verify via a totally analogous computation that, for every $u,v \in \Delta_{[a,b]}$, $s \in [a,b]$,
\[
\hat{Z}_{s,v} =\hat{Z}_{s,u} \circ  \left[ id_{T^{< \kappa}(V)} \boxtimes L_{\mathbf{X}^{< \kappa}_{u,v}} \right]
 + R^{\hat{Z}, 2}_{u,v}(s). \]
 with 
 \[
 R^{\hat{Z}, 2}_{u,v}(s) = \left(\hat{R}^1_{s,v}(s) +  R^{2}_{u,s}(s)\right) \mathbbm{1}_{v \geq s \geq u} + \hat{R}^1_{u,v}(s)\mathbbm{1}_{v \geq u >s} + R^2_{u,v}(s)  \mathbbm{1}_{s > v \geq u},
 \]
 where $\eta, \xi \in T^{< \kappa}(V)$, $\hat{R}^1[\eta, \xi] = R^1[ \xi, \eta]$.\newline
 This concludes the claim.
\end{proof}

\begin{proposition}
    Let $U$, $L$ be defined as above, then the following identity holds
\[
\int_a^b L_s d\mathbf{X}_s = \int_a^b U_u d\mathbf{X}_u. 
\]  
\end{proposition}
\begin{proof}
To prove the claim, it suffices to show that the identity:
\begin{equation}
2 \int_a^b L_s \, d\mathbf{X}_s
= \int_{[a,b]\times[a,b]} \hat{Z}_{s,u} \, d(\mathbf{X}_s, \, \mathbf{X}_u)
=  2\int_a^b U_u \, d\mathbf{X}_u
\end{equation}
holds.

We begin by establishing the first equality in \eqref{identity_tringular_Fubini}. The argument relies on an approximation
procedure based on the map $\Omega^L$ associated with the rough integral
$\int_a^b L_s \, d\mathbf{X}_s$.

Concretely, fix a partition $\mathcal{P}=\{s_i\}_i$ of $[a,b]$ and define the corresponding map
\[
\Omega^L_{\mathcal{P}}
:= \sum_{s_i \in \mathcal{P}} L_{s_i}\, \mathbf{X}^{\geq 1}_{s_i,s_{i+1}} .
\]
Recalling that
\[
L_s = \int_a^s Z_{u,s}\, d\mathbf{X}_u
+ Z^L_{s},
\]
we may expand $\Omega^L_{\mathcal{P}}$ using the  representation of the integral
defining $L$. This yields
\begin{align*}
&\Omega^L_{\mathcal{P}} \\
&= \sum_{(s_i,u_{j+1})\in \mathcal{P}^2\cap \Delta_{[a,b]}}
Z_{u_j,s_i} \circ
\left[\mathbf{X}^{\geq 1}_{u_j,u_{j+1}} \boxtimes
\mathbf{X}^{\geq 1}_{s_i,s_{i+1}}\right]
+ Z_{u_j,s_i} \circ
[\pi_{\geq 1} \boxtimes \pi_{\geq 1}]\circ \Delta\mathbf{X}_{s_j, s_{j+1}}\,
\mathbbm{1}_{\{s_i=u_j\}} + r^\mathcal{P} \\
&= \sum_{(s_i,u_{j+1})\in \mathcal{P}^2\cap \Delta_{[a,b]}}
Z_{u_j,s_i} \circ
\left[\mathbf{X}^{\geq 1}_{u_j,u_{j+1}} \boxtimes
\mathbf{X}^{\geq 1}_{s_i,s_{i+1}}\right]
+ Z_{u_j,s_i} \circ
\left[\mathbf{X}^{\geq 1}_{u_j,u_{j+1}} \boxtimes
\mathbf{X}^{\geq 1}_{s_i,s_{i+1}}\right]
\,\mathbbm{1}_{\{s_i=u_j\}} + r^\mathcal{P}.
\end{align*}
Where the term $r^\mathcal{P}$ is arising from the approximation of the integral. This term goes to 0 as the partition mesh goes to 0.\newline
Symmetrising over $\mathcal{P}^2$, we may rewrite this as
\begin{align*}
\Omega^L_{\mathcal{P}}
= \frac12 \sum_{(s_i,u_{j+1})\in \mathcal{P}^2}
Z_{u_j,s_i} \circ
\left[\mathbf{X}^{\geq 1}_{u_j,u_{j+1}} \boxtimes
\mathbf{X}^{\geq 1}_{s_i,s_{i+1}}\right]
+ Z_{u_j,s_i} \circ
\left[\mathbf{X}^{\geq 1}_{u_j,u_{j+1}} \boxtimes
\mathbf{X}^{\geq 1}_{s_i,s_{i+1}}\right]
\,\mathbbm{1}_{\{s_i=u_j\}} + r^\mathcal{P}.
\end{align*}
The right-hand side, excluding the term $r^\mathcal{P}$, coincides precisely with the approximating map associated with the
double integral
\[
\int_{[a,b]\times[a,b]} \hat{Z}_{s,u}\,
d\left(\mathbf{X}_s, \mathbf{X}_u\right)
\]
computed over the rectangular partition $\mathcal{P}^2$.

Consequently, for any $\varepsilon>0$, there exists a sufficiently fine partition
$\mathcal{P}$ such that $|r^\mathcal{P}(s_i)|< \frac{\epsilon}{2}$ and
\[
\left|
\int_a^b L_s \, d\mathbf{X}_s
- \frac12 \int_{[a,b]\times[a,b]} \hat{Z}_{s,u} \,
d\left(\mathbf{X}_s, \mathbf{X}_u \right)
\right|
\leq
\left|
\int_a^b L_s \, d\mathbf{X}_s - \Omega^L_{\mathcal{P}}
\right|
+ \frac12
\left|
\Omega^L_{\mathcal{P}}
- \int_{[a,b]\times[a,b]} \hat{Z}_{s,u} \,
d\left(\mathbf{X}_s, \mathbf{X}_u \right)
\right|
< \frac{\varepsilon}{2}.
\]
This proves the first equality in \eqref{identity_tringular_Fubini}.

The proof of the second equality proceeds analogously. Introduce the map
\[
\Omega^U_{\mathcal{P}}
:= \sum_{u_i \in \mathcal{P}} U_{u_i}\,
\mathbf{X}^{\geq 1}_{u_i,u_{i+1}},
\]
where
\[
U_u = \int_u^b Z_{u,s}\, d\mathbf{X}_s
- Z^U_{u}.
\]
A computation entirely parallel to the one above shows that
\begin{align*}
\Omega^U_{\mathcal{P}}
=
\frac{1}{2} \sum_{(s_i,u_j)\in \mathcal{P}^2}
Z_{s_i,u_j} \circ
\left[\mathbf{X}^{\geq 1}_{s_i,s_{i+1}} \boxtimes
\mathbf{X}^{\geq 1}_{u_j,u_{j+1}}\right]
- Z_{s_i,u_j} \circ
\left[\mathbf{X}^{\geq 1}_{s_i,s_{i+1}} \boxtimes
\mathbf{X}^{\geq 1}_{u_{j+1},u_j}\right]
\,\mathbbm{1}_{\{s_i=u_j\}} + r^\mathcal{P}.
\end{align*}
The first term coincides with the map $\Omega^{\hat{Z}}$, associated to
\[
\frac{1}{2}\int_{[a,b]\times[a,b]} \hat{Z}_{s,u}\,
d\left(\mathbf{X}_s, \mathbf{X}_u \right),
\]
while the second and third term vanish in the limit as $|\mathcal{P}|\to 0$. For the second term, this may be seen, for instance, by the change of variables
$(s,u)\mapsto (s,u-s)$.
This concludes the proof.
\end{proof}

The proof of the previous proposition shows that the symmetrisation procedure is compatible with the two-parameter rough integral on the square $[a,b]$, and that it recovers the triangular integral up to the expected factor $2$. This motivates the following definition.

\begin{definition}[Two-parameter rough integral over the simplex]
Let $Z \in \mathcal{D}^{p}_{\mathbf{X}}\!\left(\Delta_{[a,b]},\,\mathrm{Hom}(V \boxtimes V, F)\right)$ and let $\widehat{Z}$ be its symmetrisation, defined as above. We define the two-parameter rough integral of $Z$ over the simplex $\Delta_{[a,b]}$ by
\begin{equation}\label{defining_identity_2p_simplex}
\int_{\Delta_{[a,b]}} Z_{s,u}\, d\left(\mathbf{X}_s,\mathbf{X}_u\right)
:= \frac{1}{2}\int_{[a,b]\times[a,b]} \widehat{Z}_{s,u}\, d\left(\mathbf{X}_s,\mathbf{X}_u\right).
\end{equation}
\end{definition}


\begin{remark}
The fact that the integral
\[
\int_{\Delta_{[a,b]}} Z_{s,u} \, d\left( \mathbf{X}_s, \mathbf{X}_u \right)
\]
satisfies the identity \eqref{defining_identity_2p_simplex} implies that it inherits all the continuity and stability properties of the two-parameter rough integral. Alternatively, these properties can be recovered by exploiting the identity \eqref{identity_tringular_Fubini}, proved in the proposition above. Another consequence of this definition is that the controlled path $I^\Delta_{\mathbf{X}}(Z)$ associated with the simplex integral can be obtained immediately by expressing the simplex integral as an integral over the square and then applying the construction presented in the previous section. As shown in the next section, this ensures that solutions to the Schwinger–Dyson kernel equation can be characterized as fixed points in the space of two-parameter controlled paths.
\end{remark}


To improve readability, whenever no ambiguity arises from the context, we write \(I_{\mathbf X, \mathbf X}(Z)\) in place of \(I_{\mathbf X}^{\Delta}(Z)\).

\begin{remark}
A natural
    question that may arise a the end of this section is whether the rough integral over a simplex can be extended to controlled
    rough paths $\tilde Z$ controlled by two distinct rough paths $\mathbf X$ and
    $\mathbf Y$, while still agreeing with the corresponding iterated rough integrals.
    In general, such an extension is not available. Indeed, in the iterated construction
    the objects $L$ and $U$ are the inner rough integrals. For the outer integrations
    to be well-defined, these inner integrals would have to be one-parameter controlled
    rough paths controlled by $\mathbf Y$ and $\mathbf X$, respectively. This need not
    be the case for general, unrelated rough paths $\mathbf X$ and $\mathbf Y$. Hence
    such an extension can only be expected under additional compatibility assumptions, for example the case where one between $\mathbf{X}$ and $\mathbf{Y}$ has bounded variation.
    \end{remark}
\section{Differential equation for jointly controlled paths} \label{composition and stability}
We begin this section by introducing the last ingredient needed to give a rigorous meaning to the class of equations described by \eqref{skeleton_equation}, which encompasses the signature kernel trick and the Schwinger-Dyson kernel equation. 
Indeed, the formulation of \eqref{skeleton_equation} involves joint rough integrals whose integrands are obtained by applying a smooth function $g$ to a two-parameter controlled path. 
Thus, before the equation can be interpreted, one must verify that this nonlinear transformation preserves the two-parameter controlled structure. 
This is exactly what we prove below: for any smooth function $g$ and two-parameter controlled rough path $Z$, it is possible to construct a two-parameter controlled rough path above $g(Z^{(0,0)})$.

Having established this stability under smooth composition, we then introduce a notion of solution to \eqref{skeleton_equation} within the class of jointly controlled paths.
In particular, we formulate \eqref{skeleton_equation} as a fixed-point problem in an appropriate controlled-path space and define solutions as jointly controlled paths satisfying the equation in this sense.

\begin{proposition}\label{stability_under_composition}
Let $S$ be a finite-dimensional Banach space. Let $\mathbf{X}\in \Omega G_p([a,b],V)$ and $\mathbf{Y}\in \Omega G_q([c,d],W)$, and set $\kappa=\lfloor p\rfloor$ and $\lambda=\lfloor q\rfloor$. Let $g\colon F\to S$ be Fr\'echet smooth, and let $Z$ be a two-parameter $(\mathbf{X},\mathbf{Y})$-controlled path.
Then the composition $g(Z)$ defines a two-parameter controlled path
\[
g(Z) \in \mathcal{D}^{p,q}_{\mathbf{X},\mathbf{Y}}([a,b]\times[c,d], S).
\]
Moreover, $g(Z)$ is characterized by the identity
\[
g(Z)_{t,v}[\eta,\xi] = g\left(Z^{(0,0)}_{t,v}\right)[\eta,\xi],
\qquad
\text{for all } \eta \in V^{\otimes 0},\ \xi \in W^{\otimes 0}.
\]

The higher level components of $g(Z)$ are given as follows. For
$\xi_j \in V^{\otimes j}$ and $\eta_k \in W^{\otimes k}$, define
\[
g(Z_{t,v})[\xi_j,\eta_k] :=
\begin{cases}
\displaystyle
\displaystyle
\sum_{m=1}^{\kappa-1} \frac{1}{m!}
D^m g\left(Z^{(0,0)}_{t,v}\right) \circ Z_{t,v}^{\otimes m}
\circ \left(\tilde{\Delta}^m \boxtimes \mathrm{id}_{T^{<\lambda(W)}}\right)
[\xi_j,\eta_k],
& j>0,\ k=0, \\[2ex]
\displaystyle
\sum_{n=1}^{\lambda-1} \frac{1}{n!}
D^n g\left(Z^{(0,0)}_{t,v}\right) \circ Z_{t,v}^{\otimes n}
\circ \left(\mathrm{id}_{T^{<\kappa(V)}} \boxtimes \tilde{\Delta}^n\right)
[\xi_j,\eta_k],
& j=0,\ k>0, \\[2ex]
\displaystyle
\sum_{\ell=1}^{\lambda+\kappa-2} \frac{1}{\ell!}
D^\ell g\left(Z^{(0,0)}_{t,v}\right) \circ Z_{t,v}^{\otimes \ell}
\circ \widetilde{(\Delta \boxtimes \Delta)}^{\ell}
[\xi_j,\eta_k],
& j>0,\ k>0 .
\end{cases}
\]
Here, $\Delta$ denotes the shuffle coproduct, $\tilde{\Delta}$
denotes the restriction $\pi_{\geq 1} \circ \Delta$ and $\Delta^m$ denotes the $m$-fold iterate of $\Delta$.
\end{proposition}
The previous result ensures that the term $g(Z)$ is again a two-parameter controlled path, so the joint integrals in \eqref{skeleton_equation} can be defined.  
For well-posedness our two examples, however, this qualitative closure is not enough:  one must control the size of $g(Z)$ in the controlled norm and compare $g(Z)$ and $g(\tilde Z)$ when either the input $Z$ or the driving rough paths are perturbed.  
The next proposition provides precisely these quantitative bounds, showing that composition defines a (locally) Lipschitz map with respect to the distance $d$.
\begin{proposition}\label{stability_proposition}
    Let $\mathbf{X}, \tilde{\mathbf{X}} \in \Omega G_{p, \omega_\mathbf{X}}([a,b], V)$, $\mathbf{Y}, \tilde{\mathbf{Y}} \in \Omega G_{q, \omega_\mathbf{Y}}([c,d], W)$ and  $Z \in \mathcal{D}^{p,q}_{\mathbf{X},\mathbf{Y}}$, $\tilde{Z} \in \mathcal{D}^{p,q}_{\tilde{\mathbf{X}},\tilde{\mathbf{Y}}}$. The composition of $Z$ with $g$ satisfies 
    \[
    \|g(Z)\|_{\mathcal{D}^{p,q}_{\mathbf{X}, \mathbf{Y}}} \leq C_{p,q,g,\omega_\mathbf{X}, \omega_\mathbf{Y}, Z} \|Z\|_{\mathcal{D}^{p,q}_{\mathbf{X}, \mathbf{Y}}},
    \]
    moreover the following stability estimate holds
    \[
    d(g(Z), g(\tilde{Z})) \leq C_{p,q,g,\omega_\mathbf{X}, \omega_\mathbf{Y}, Z, \tilde{Z}} d(Z, \tilde{Z}).
    \]


\end{proposition}

We postpone the proof of these results to the appendix in order to maintain the flow of the exposition.\newline

With these composition and stability estimates in hand, all terms appearing in \eqref{skeleton_equation} are now well-defined within the two-parameter controlled framework. In particular, we can now give a precise meaning to two-parameter Rough Differential Equations.


\begin{definition}[Solution to a two-parameter rough differential equation]
\label{solution_2d_RDE}
Let $\mathbf{X},\mathbf{Y}$ be as above. We consider equations of the form
\[
Z_{t,v}
=
z
+
\int_{\mathcal D_{t,v}}
\Gamma^{g}(Z)_{s,u}\,d(\mathbf X_s,\mathbf Y_u),
\qquad (t,v)\in [a,b]\times[c,d],
\]
where $Z_{a,c}\in F$, and where $(g,\mathcal D_{t,v},\Gamma)$ is required to be of one of the following two types.

\smallskip
\noindent{Rectangular type.}
\[
g:F\to \mathrm{Hom}(V\boxtimes W,F),
\qquad
\mathcal D_{t,v}:=[a,t]\times[c,v], \qquad 
\Gamma^{g}(Z)_{s,u}:=g(Z)_{s,u}.
\]


\smallskip
\noindent{Triangular type.}
\[
g:F\boxtimes F\to \mathrm{Hom}(V\boxtimes W,F),
\qquad
\mathcal D_{t,v}:=\Delta_{[t,v]},
\qquad
\Gamma^{g}(Z)_{s,u}:=g(Z_{t, \cdot},Z)_{s,u}.
\]

We say that the equation is solved by a controlled path
\[
Z\in \mathcal D^{p,q}_{\mathbf X,\mathbf Y}\big([a,b]\times[c,d],F\big)
\]
if
\[
Z= z +I_{\mathbf X,\mathbf Y}\big(\Gamma^{g}(Z)\big).
\]
\end{definition}

\begin{remark}
    From our construction of the rough integral on triangular domain, the choice of $\mathcal{D}_{t,v} = \Delta_{[t,v]}$  implicitly forces the identity $\mathbf{Y} = \mathbf{X}$ in the above definition. In the same setting it is worth pointing out that, for a two-parameter controlled  the path $Z$,  the path $(s,u) \to (Z_{t,s}, Z_{s,u})$ is  a two-parameter controlled path, therefore the integral in the definition is well-defined.
\end{remark}

We devote the remainder of the paper to the analysis of two examples of two-parameter RDEs, namely the signature kernel and the Schwinger–Dyson kernel equation, which illustrate the scope of the theory developed above.
\subsection{The signature kernel}
The first of the two applications we consider concerns the signature kernel, a kernel method for sequential data that uses the signature as a feature map. This approach was introduced in \cite{salvi2021signature} and has since been applied in a variety of settings; see, for example, \cite{lemercier2021distribution, toth2020bayesian, alden2025signature}.
In what follows, we briefly review the definition of the signature kernel and show how it can be related to the solution of a two-parameter rough differential equation.
We begin by recalling the notions of the signature and the associated signature kernel.
\newline 

Let $V=W$ be a finite-dimensional Hilbert space, and let $\mathbf{X}$ and $ \mathbf{Y} $ to be defined as in the previous sections.
Denote with $\Gamma$ the Gamma function and define
\[
\beta_p
:=
p^2\!\left(1+\sum_{i=3}^\infty
\left(\frac{2}{i-2}\right)^{\frac{\floor{p}+1}{p}}\right).
\]
Without loss of generality, from this point onwards, by replacing controls $\omega_{\mathbf{X}}$ and $ \omega_{\mathbf{Y}}$ if necessary, we may assume that 
\begin{align*}
&\left|\mathbf{X}^{(n)}_{s,t}\right| \leq \frac{\omega_{\mathbf{X}}(s,t)^{n/p}}{\beta_p \Gamma(n/p + 1)} \text{ for all } a \le s \le t \le b\text{ and } n \leq \kappa, \\
&\left|\mathbf{Y}^{(n)}_{u,v}\right| \leq \frac{\omega_{\mathbf{Y}}(u,v)^{n/q}}{\beta_q \Gamma(n/q + 1)} \text{for all } c \le u \le v \le d \text{ and } n \le \lambda.\\
\end{align*}


\begin{definition}
Let $\mathbf{X}$ be a geometric $p$-rough path on $V$, with $p \notin \mathbb{N}$.  
The signature of $\mathbf{X}$ is the unique extension of $\mathbf{X}$ to a path
\[
S(\mathbf{X}) : \Delta_{[a,b]} \to T((V)),
\]
such that $S^{(n)}(\mathbf{X}) = \mathbf{X}^{(n)}$ for every non-negative integer $n \leq \kappa$ and the following properties hold:
\begin{itemize}
\item
for every integer $n > \kappa$  the $n$-th level of the signature satisfies the estimate 
\[
\left|S^{(n)}(\mathbf{X})_{s,t}\right |
\le
\frac{\omega_{\mathbf{X}}(s,t)^{n/p}}{\beta_p\,\Gamma(n/p+1)} \quad \text{for all } a \le s \le t  \le b.
\]
\item
$S(\mathbf{X})$ is multiplicative, that is,
\[
S(\mathbf{X})_{s,s} = \mathbf{1}, \quad 
S(\mathbf{X})_{s,u} = S(\mathbf{X})_{s,t}\otimes S(\mathbf{X})_{t,u}
\quad \text{for all } a \le s \le t \le u \le b.
\]
\end{itemize}
\end{definition}

For a proof of this classical result in rough path theory we refer the reader to Theorem 3.7 in \cite{lyons2007differential}.\newline

It can be seen from the previous definition that the signatures \( S(\mathbf{X}) \) and \( S(\mathbf{Y}) \) take values in the Hilbert space
\[
F := \Big\{ \eta \in T((V)) : \|\eta\| < \infty \Big\},
\]
where
\[
\|\eta\| := \left( \sum_{k=0}^\infty \left\|\eta^{(k)}\right\|^2 \right)^{1/2},
\]
where each tensor power \(V^{\otimes k}\) is endowed with the Hilbert tensor norm induced by the fixed inner product on \(V\).

Within this framework, the signature kernel is defined as the map
\[
K^{\mathbf{X}, \mathbf{Y}} : [a,b] \times [c,d] \to \mathbb{R}, 
\qquad 
K^{\mathbf{X}, \mathbf{Y}}_{t,v} := \big\langle S(\mathbf{X})_{a,t},\, S(\mathbf{Y})_{c,v} \big\rangle,
\]
where \(\langle \cdot, \cdot \rangle\) denotes the induced inner product on $F$.

An important result linking the signature kernel to joint rough integrals appears in Proposition~6.3 of \cite{cass2023fubini}, where the signature kernel is obtained as a  two-parameter rough integral of the jointly controlled path
\begin{align*}
    & \Psi : [a,b] \times [c,d] \to \mathrm{Hom}\!\left(T^{<\kappa}(V) \boxtimes T^{<\lambda}(V), \mathrm{Hom}(V \boxtimes V, \mathbb{R})\right), \\
    & \Psi_{t,v}[\eta,\xi]
      = 
        \langle S(\mathbf{X})_{a,t} \otimes \eta \otimes \cdot,\,
                S(\mathbf{Y})_{c,v} \otimes \xi \otimes \cdot \rangle,
        \qquad \eta \in T^{<\kappa}(V), \xi \in T^{<\lambda}(V) ,
\end{align*}
via the identity
\begin{equation} \label{sig_pei_function}
K^{\mathbf{X},\mathbf{Y}}_{t,v}
=
1 + \int_{[a,t]\times[c,v]} \Psi_{s,u}\, d(\mathbf{X}_s,\mathbf{Y}_u).
\end{equation}
The identity above allows to construct a two-parameter controlled rough path above $K^{\mathbf{X}, \mathbf{Y}}$. In what follows, we use the same notation for the signature kernel viewed as a controlled path and for its trace, and clarify the intended meaning whenever necessary. \newline
Equipped with the results discussed earlier in this section, one can recognize the map $\Psi$ as arising from the composition of a smooth function with this controlled path.
To see this, consider the map
\begin{equation}\label{def_g_sig}
g : \mathbb{R} \to \mathrm{Hom}(V \boxtimes V, \mathbb{R}),
\qquad
g(z)[\eta,\xi] = z \langle \eta,\xi\rangle.
\end{equation}
When composed with the signature kernel, this yields the controlled path 
\begin{equation}\label{g_composed_sig}
\begin{aligned}
&g(K^{\mathbf{X},\mathbf{Y}}) \in \mathcal{D}^{p,q}_{\mathbf{X}, \mathbf{Y}}([a,b]\times [c,d], \mathrm{Hom}(V\boxtimes V, \mathbb{R})),\\
&g(K^{\mathbf{X},\mathbf{Y}})[\eta,\xi](\tilde{\eta},\tilde{\xi})
=
\langle S(\mathbf{X}) \otimes \eta \otimes \tilde{\eta},\,
        S(\mathbf{Y}) \otimes \xi \otimes \tilde{\xi} \rangle,
\qquad
\tilde{\eta},\tilde{\xi} \in V.
\end{aligned}
\end{equation}

In particular, this shows that the integrand $\Psi$ can be expressed as a smooth functional of the signature kernel understood as the controlled path constructed above the integral \eqref{sig_pei_function}, and therefore fits naturally within the framework of two-parameter controlled differential equations.


his naturally leads to the question of the well-posedness of the two-parameter RDE 
\[
Z_{t,v} = 1 +\int_{[a,t] \times [c, v]}   g(Z)_{s,u} d\left(\mathbf{X}_s, \mathbf{Y}_u\right),
\]
which we analyse in the next theorem.

\begin{theorem}\label{uniqueness_sig_kernel}
Let \(g\) be defined by \eqref{def_g_sig}. If \(Z\) is a solution to
\begin{equation}\label{signature_kernel_RDE}
Z_{t,v}
=
1+\int_{[a,t]\times[c,v]} g(Z)_{s,u}\,d(\mathbf{X}_s,\mathbf{Y}_u)
\end{equation}
in the sense of Definition \ref{solution_2d_RDE}, then
\(
Z = K^{\mathbf{X},\mathbf{Y}}.
\)
In particular, \(K^{\mathbf{X},\mathbf{Y}}\) is the unique solution to \eqref{signature_kernel_RDE}.

Moreover, if \(\tilde{\mathbf{X}} \in \Omega G_{p,\omega_{\mathbf{X}}}([a,b],V)\) and
\(\tilde{\mathbf{Y}} \in \Omega G_{q,\omega_{\mathbf{Y}}}([c,d],V)\), then
\[
\|K^{\mathbf{X},\mathbf{Y}}-K^{\tilde{\mathbf{X}},\tilde{\mathbf{Y}}}\|_{\infty}
\leq
C_{p,q,\omega_{\mathbf{X}},\omega_{\mathbf{Y}}}
\big(
\|\mathbf{X}-\tilde{\mathbf{X}}\|_{p}
+
\|\mathbf{Y}-\tilde{\mathbf{Y}}\|_{q}
\big).
\]
\end{theorem}
\begin{proof}
The proof of the first claim is based on a local uniqueness estimate for the difference of two fixed points. To this end, fix $\tilde a \in [a,b)$ and
$\tilde c \in [c,d)$, and let
\[
\delta \le \min\{\, b-\tilde a,\; d-\tilde c \,\}.
\]
Additionally consider the values $p' \in (p, \kappa + 1)$, $q' \in (q, \lambda+1)$.\newline 
Define the map
\[
\Gamma^{\tilde a,\tilde c} :
\mathcal{D}^{p',q'}_{\mathbf{X},\mathbf{Y}}
\big([\tilde a,\tilde a+\delta]\times[\tilde c,\tilde c+\delta],\mathbb{R}\big)
\;\longrightarrow\;
\mathcal{D}^{p',q'}_{\mathbf{X},\mathbf{Y}}
\big([\tilde a,\tilde a+\delta]\times[\tilde c,\tilde c+\delta],\mathbb{R}\big)
\]
to be the controlled path constructed above the integral 
\[
z + \int_{[\tilde a,t]\times[\tilde c,v]}
g\!\left(Z\right)_{s,u}\, d(\mathbf{X}_s,\mathbf{Y}_u),
\]
according to the definition \eqref{definition_I_X_Y}. In the above the value $z \in \mathbb{R}$ is fixed.
Let $\tilde{K}$ be another fixed point to the map $\Gamma^{a, c}$ that belongs to the ball
\[
B = \left \{ Z \in \mathcal{D}^{p', q'}_{\mathbf{X}, \mathbf{Y}}([a,a+\delta] \times [c, c+\delta], \mathbb{R}): Z_{a,c} = \langle \cdot  ,  \cdot  \rangle_{T((V))},  \| Z \|_{\mathcal{D}_{\mathbf{X}, \mathbf{Y}}} \leq M  \right \},
\]
where $M > \max\left( \|K\|_{\mathcal{D}_{\mathbf{X}, \mathbf{Y}}} , \|\tilde K\|_{\mathcal{D}_{\mathbf{X}, \mathbf{Y}}} \right)$.\newline
We begin by considering the level $(k,j)$ where at least one between $k$ and $j$ is equal to $0$. From the definition of $I_{\mathbf{X}, \mathbf{Y}}(g(K)-g(\tilde{K}))$ we write the remainder in the first parameter for this path as  
\[
R^{1}_{s,t}(u)[\eta, \xi] = I_{\mathbf{X}, \mathbf{Y}}(g(K)-g(\tilde{K}))_{t,u}[\eta, \xi] - I_{\mathbf{X}, \mathbf{Y}}(g(K)-g(\tilde{K}))_{s,u} \circ \left[ L_{\mathbf{X}^{< \kappa}_{s,t}} \boxtimes id_{T^{< \lambda}(V)} \right][\eta, \xi]. 
\]
where $\eta \in V^{\otimes j}, \xi \in V^{\otimes k}$.\newline
Concretely, when $k=j=0$, 
\[
R^{1}_{s,t}(u)[\mathbf{1}, \mathbf{1}] = \int_{[a, t]\times [c,u]}\left( g(K)_{r,w} - g(\tilde{K})_{r,w} \right) \, d(\mathbf{X}_r, \mathbf{Y}_w) - I_{\mathbf{X}, \mathbf{Y}}(g(K)-g(\tilde{K}))_{s,u} \left[\mathbf{X}^{< \kappa}_{s,t} \boxtimes \mathbf{1}\right],
\]
when $\eta \in V^{\otimes j} $ with $j \ge 1$, 
\begin{align*}
&R^{1}_{s,t}(u)[\eta, \mathbf{1}]\\
&= \int_{[c,u]} \left(g(K)_{t,w} - g(\tilde{K})_{t,w} \right) \, d \mathbf{Y}_w [\eta] - \int_{[c,u]}\left( g(K)_{s,w} - g(\tilde{K})_{s,w}\right) \circ \left[L_{\mathbf{X}^{< \kappa}_{s,t}} \boxtimes id_{T^{< \lambda}(V)}\right]\, d \mathbf{Y}_w [\eta]\\
&= \int_{[c,u]} \left(g(K)_{t,w} - g(\tilde{K})_{t,w} \right) - \left( g(K)_{s,w} - g(\tilde{K})_{s,w}\right) \circ \left[L_{\mathbf{X}^{< \kappa}_{s,t}} \boxtimes id_{T^{< \lambda}(V)} \right]\, d \mathbf{Y}_w[\eta]
\end{align*}
and when $\xi \in V^{\otimes k}$ with $k \ge 1$
\begin{align*}
& R^{1}_{s,t}(u)[\mathbf{1}, \xi] = \int_{[s,t]} \left(g(K)_{r,u} - g(\tilde{K})_{r,u} \right)\, d \mathbf{X}_r [\xi] - \left( g(K)_{s,u} - g(\tilde{K})_{s,u}\right) \circ \left[\mathbf{X}^{< \kappa}_{s,t} \boxtimes \xi\right].
\end{align*}
From classical estimates on one-parameter rough integrals (e.g. Section 2 in \cite{cass2022combinatorial}) and Proposition \ref{stability_proposition} we recover for this case 
\[
\left|R^1_{s,t}(u)[\eta, \xi] \right| \leq C d(K, \tilde K)\omega_{\mathbf{X}}(s,t)^{(\kappa - j)/p}\|\eta\|_{V^{\otimes j }} \|\xi\|_{V^{\otimes k }},
\]
here, and throughout the proof, $C$ denotes a positive constant that may vary from line to line, but depends only on $g,p,q,\omega_{\mathbf{X}},\omega_{\mathbf{Y}},\|K\|_{\mathcal{D}_{\mathbf{X},\mathbf{Y}}},$ and $\|\tilde K\|_{\mathcal{D}_{\mathbf{X},\mathbf{Y}}}$.\newline
Applying the same argument to $R^2$, we find that, upon restricting to these levels, the following bound holds
\[
\left| R^2_{u,v}(s)[\eta, \xi] \right| \leq C d(K, \tilde K)\omega_{\mathbf{Y}}(u,v)^{(\lambda - k)/q}\|\eta\|_{V^{\otimes j }} \|\xi\|_{V^{\otimes k }}.
\]
For  the remainder of remainders, whenever at least one between $(j,k)$ is zero we have 
\begin{align*}
&\mathbf{R} \begin{pmatrix} s,t \\ u,v \end{pmatrix}[\eta, \xi] \\
&=  I_{\mathbf{X}, \mathbf{Y}}(g(K)-g(\tilde{K}))_{t,v}[\eta, \xi] - I_{\mathbf{X}, \mathbf{Y}}(g(K)-g(\tilde{K}))_{s,v} \circ \left[ L_{\mathbf{X}^{< \kappa}_{s,t}} \boxtimes id_{T^{< \lambda}(V)} \right][\eta, \xi]\\
&\quad - \left(I_{\mathbf{X}, \mathbf{Y}}(g(K)-g(\tilde{K}))_{t,u} - I_{\mathbf{X}, \mathbf{Y}}(g(K)-g(\tilde{K}))_{s,u} \circ \left[ L_{\mathbf{X}^{< \kappa}_{s,t}} \boxtimes id_{T^{< \lambda}(V)} \right]\right) \circ [id_{T^{< \kappa}(V)} \boxtimes L_{\mathbf{Y}^{< \lambda}_{u,v}}][\eta, \xi].
\end{align*}
We can now write 
\begin{align*}
&\mathbf{R} \begin{pmatrix} s,t \\ u,v \end{pmatrix}[\mathbf{1}, \mathbf{1}] \\
&=  I_{\mathbf{X}, \mathbf{Y}}(g(K)-g(\tilde{K}))_{t,v}[\mathbf{1}, \mathbf{1}] -  \int_{[c,v]} (g(K) -(\tilde{K}))_{s,w} d\mathbf{Y}_w \left[ \mathbf{X}^{\leq \kappa}_{s,t}\right] + \int_{[c,v]} (g(K) -(\tilde{K}))_{s,w} d\mathbf{Y}_w \left[ \mathbf{X}^{(\kappa)}_{s,t}\right] \\
&\quad - \Bigg(\int_{[s,t]}(g(K)-g(\tilde{K}))_{r,u}d\mathbf{X}_r \left[ \mathbf{Y}^{\leq \lambda}_{u,v}\right] - \int_{[s,t]}(g(K)-g(\tilde{K}))_{r,u}d\mathbf{X}_r \left[ \mathbf{Y}^{(\lambda)}_{u,v}\right] - (g(K)-g(\tilde{K}))_{s,u}\left[ \mathbf{X}^{\leq \kappa}_{s,t} \boxtimes \mathbf{Y}^{\leq \lambda}_{u,v}\right] \\ 
&\qquad \quad  + (g(K)-g(\tilde{K}))_{s,u}\left[ \mathbf{X}^{(\kappa)}_{s,t} \boxtimes \mathbf{Y}^{\leq \lambda}_{u,v}\right]  + (g(K)-g(\tilde{K}))_{s,u}\left[ \mathbf{X}^{\leq \kappa}_{s,t} \boxtimes \mathbf{Y}^{(\lambda)}_{u,v}\right] \Bigg) .
\end{align*}
We can now decompose $\mathbf{R} \begin{pmatrix} s,t \\ u,v \end{pmatrix}[\mathbf{1}, \mathbf{1}]$ as 
\[
\mathbf{R} \begin{pmatrix} s,t \\ u,v \end{pmatrix}[\mathbf{1}, \mathbf{1}] = D \begin{pmatrix} s,t \\ u,v \end{pmatrix}  + E \begin{pmatrix} s,t \\ u,v \end{pmatrix} + F \begin{pmatrix} s,t \\ u,v \end{pmatrix},
\]
where 
\begin{align*}
&D \begin{pmatrix} s,t \\ u,v \end{pmatrix} =  I_{\mathbf{X}, \mathbf{Y}}(g(K)-g(\tilde{K}))_{t,v}[\mathbf{1}, \mathbf{1}] -  \int_{[c,v]} (g(K) -(\tilde{K}))_{s,w} d\mathbf{Y}_w \left[ \mathbf{X}^{\leq \kappa}_{s,t}\right]  \\
&\quad - \Bigg(\int_{[s,t]}(g(K)-g(\tilde{K}))_{r,u}d\mathbf{X}_r \left[ \mathbf{Y}^{\leq \lambda}_{u,v}\right]  - (g(K)-g(\tilde{K}))_{s,u}\left[ \mathbf{X}^{\leq \kappa}_{s,t} \boxtimes \mathbf{Y}^{\leq \lambda}_{u,v}\right] \Bigg),\\
&E \begin{pmatrix} s,t \\ u,v \end{pmatrix} =\int_{[c,v]} (g(K) -(\tilde{K}))_{s,w} d\mathbf{Y}_w \left[ \mathbf{X}^{(\kappa)}_{s,t}\right]  - (g(K)-g(\tilde{K}))_{s,u}\left[ \mathbf{X}^{(\kappa)}_{s,t} \boxtimes \mathbf{Y}^{\leq \lambda}_{u,v}\right],\\
&F \begin{pmatrix} s,t \\ u,v \end{pmatrix} =   \int_{[s,t]}(g(K)-g(\tilde{K}))_{r,u}d\mathbf{X}_r \left[ \mathbf{Y}^{(\lambda)}_{u,v}\right] - (g(K)-g(\tilde{K}))_{s,u}\left[ \mathbf{X}^{\leq \kappa}_{s,t} \boxtimes \mathbf{Y}^{(\lambda)}_{u,v}\right].
\end{align*}
Now, applying inequality (4.13) in \cite{cass2023fubini} and using the stability of composition with smooth functions (Proposition \ref{stability_proposition}), we obtain
\[
\left|D\begin{pmatrix} s,t \\ u,v \end{pmatrix} \right|
\leq A(p,q,\mathbf{R})\, C\, d(K,\tilde{K})\, \omega_{\mathbf{X}}(s,t)^{\theta}\omega_{\mathbf{Y}}(u,v)^{\theta},
\]
where $A$ and $\theta$ are as defined in Theorem \ref{stability_theorem_2D}.\newline
Additionally, via one-parameter sewing lemma and Proposition \ref{stability_proposition} we obtain 
\[
\left|E\begin{pmatrix} s,t \\ u,v \end{pmatrix} \right| \leq C d(K, \tilde K) \omega_{\mathbf{X}}(s,t)^{p/\kappa}\omega_{\mathbf{Y}}(u,v)^{\theta}, \quad \left|D\begin{pmatrix} s,t \\ u,v \end{pmatrix} \right|\leq C d(K, \tilde K) \omega_{\mathbf{X}}(s,t)^{\theta}\omega_{\mathbf{Y}}(u,v)^{q/\lambda} . 
\]
For $\eta \in V^{\otimes j}$, with $j \ge 1$ we obtain
\begin{align*}
\mathbf{R} \begin{pmatrix} s,t \\ u,v \end{pmatrix}[\eta, \mathbf{1}] &= \int_u^v \left(g(K)_{t,w} - g(\tilde{K})_{t,w} \right) - \left( g(K)_{s,w} - g(\tilde{K})_{s,w}\right) \circ \left[L_{\mathbf{X}^{^{< \kappa}}_{s,t}} \boxtimes id_{T^{< \lambda}(V)} \right]\, d \mathbf{Y}_w[\eta]\\
& \quad-\left(g(K)_{t,u} - g(\tilde{K})_{t,u} \right) - \left( g(K)_{s,u} - g(\tilde{K})_{s,u}\right) \left[L_{\mathbf{X}^{< \kappa}_{s,t}}\eta \boxtimes \mathbf{Y}^{\ge 1}_{u,v} \right]\\
&\quad +\left(g(K)_{t,u} - g(\tilde{K})_{t,u} \right) - \left( g(K)_{s,u} - g(\tilde{K})_{s,u}\right) \left[L_{\mathbf{X}^{< \kappa}_{s,t}}\eta \boxtimes \mathbf{Y}^{(\lambda)}_{u,v} \right],
\end{align*}
which allows us to conclude, via the one-parameter sewing lemma and regularity of $\mathbf{Y}^{(\lambda)}$, that
\[
\left|\mathbf{R} \begin{pmatrix} s,t \\ u,v \end{pmatrix}[\eta, \mathbf{1}] \right|
\leq C\, d(K,\tilde{K})\,\omega_\mathbf{X}(s,t)^{(\kappa-j)/p}\omega_{\mathbf{Y}}(u,v)^{\theta}\,\|\eta\|_{V^{\otimes j}}.
\]
A similar calculation yields
\[
\left|\mathbf{R} \begin{pmatrix} s,t \\ u,v \end{pmatrix}[\mathbf{1}, \xi] \right|
\leq C\, d(K,\tilde{K})\,\omega_\mathbf{X}(s,t)^{\theta}\omega_{\mathbf{Y}}(u,v)^{(\lambda-k)/q}\,\|\xi\|_{V^{\otimes k}},
\]
whenever $\xi \in V^{\otimes k}$ with $k \geq 1$.

Since \(K\) and \(\tilde K\) are fixed points of \(\Gamma^{a,a+\delta}\),
then for every $j \in [1 :\kappa)$ and $k \in [1 :\lambda)$ it holds that
\[
K^{(j,k)} = g^{(j-1, k-1)}(K).
\]
But, since $g$ is linear, by inspection of formula \eqref{def_g_sig}, this implies that the only levels needed to show that the evaluation $\Gamma^{a, c}(K) - \Gamma^{a,c}(\tilde{K})$ is contracting for a sufficiently small $\delta$ are the levels $(j, k)$ where one between $j, k$ is zero.\newline
Since $\delta$ can be chosen uniformly, uniqueness on $[a,b]\times [c,d]$  follows by extending the solution step by step in each coordinate direction. \newline
The stability property follows from the corresponding stability estimates for joint integrals. 
More precisely, for any rectangular partition $\{(s_i,u_j)\}_{i,j}$ with sufficiently small mesh size, one can show that
\begin{align*}
&\left\| \Gamma^{a, c}(Z) - \Gamma^{a,c}(\tilde Z) \right\|_{\infty; [s_i,s_{i+1}] \times [u_j,u_{j+1}]} \\
&\quad \leq C
\Big(
|Z_{s_i,u_j}-\tilde{Z}_{s_i,u_j}| 
+ \|\mathbf{X}-\tilde{\mathbf{X}}\|_{p;[s_i,s_{i+1}]} + \|\mathbf{Y}-\tilde{\mathbf{Y}}\|_{q;[u_j,u_{j+1}]} \Big) .
\end{align*}
where $Z \in D^{p,q}_{\mathbf{X}, \mathbf{Y}}([a,b] \times [c,d])$ and $\tilde{Z} \in D^{p,q}_{\tilde{ \mathbf{X}}, \tilde{\mathbf{Y}}}([a,b] \times [c,d])$.\newline
By progressive substitution over the partition and standard $p \thinspace (\text{and }q)$-variation norm estimates, the claim follows.
\end{proof}


\subsubsection{The smooth rough path case}
Under the additional assumption that $\mathbf{X}$ and $\mathbf{Y}$ are smooth rough paths, we show that the signature kernel admits a natural interpretation as the solution of a two-parameter RDE. In this way, the associated linear PDE, first established in \cite{lemercier2024log}, emerges directly from our framework. The argument hinges on the fact that the kernel is characterized as a fixed point of the two-parameter RDE, which is precisely the structure from which the PDE arises.

 

In the smooth rough path setting considered here, Theorem ~\ref{uniqueness_sig_kernel} also yields a transparent uniqueness argument for the kernel and provides a natural starting point for approximation procedures, such as replacing the drivers by smooth (or piecewise-linear) approximations.

More precisely, let $\mathbf{X}$ and $\mathbf{Y}$ be smooth rough paths and define the diagonal derivatives 
\[
\boldsymbol{x}_t := \left.\partial_t\right|_{t=s} \mathbf{X}_{s,t},
\qquad
\boldsymbol{y}_v := \left.\partial_v |_{v=u} \mathbf{Y}_{u,v}\right.
\]
From the previous section, the signature kernel admits the representation
\[
K^{\mathbf{X}, \mathbf{Y}}_{t,v}
= 1 + \int_a^t \int_c^v g\!\left(K^{\mathbf{X}, \mathbf{Y}}\right)_{s,u}
\big[ \boldsymbol{x}_s, \boldsymbol{y}_u \big]\, du\, ds,
\]
or, equivalently,
\[
\frac{\partial^2 K^{\mathbf{X}, \mathbf{Y}}_{t,v}}{\partial t\, \partial v}
= g\!\left(K^{\mathbf{X}, \mathbf{Y}}\right)_{t,v}
\big[ \boldsymbol{x}_t, \boldsymbol{y}_v \big].
\]

Fix a basis $\{e_1,\dots,e_m\}$ of $V$, let $I \neq \emptyset$ be a word in the alphabet $\{1, \dots, m\}$, and define the maps
\[
\begin{aligned}
f \colon\; & [a,b] \times [c,d] \longrightarrow \mathbb{R}, 
& f(t,v) &:= (\pi_0 \boxtimes \pi_0)\!\left( K^{\mathbf{X},\mathbf{Y}}_{t,v} \right), \\[0.5em]
\phi \colon\; & [a,b] \times [c,d] \longrightarrow \mathrm{Hom}\!\left(T^{<\lambda}(V), \mathbb{R}\right), 
& \phi(t,v)[e_I] &:= 
\big\langle L^*_{S(\mathbf{Y})_{c,v}} S(\mathbf{X})_{a,t}
- f(t,v)\,\mathbf{1},\, e_I \big\rangle, \\[0.5em]
\psi \colon\; & [a,b] \times [c,d] \longrightarrow \mathrm{Hom}\!\left(T^{<\kappa}(V), \mathbb{R}\right), 
& \psi(t,v)[e_I] &:= 
\big\langle L^*_{S(\mathbf{X})_{a,t}} S(\mathbf{Y})_{c,v}
- f(t,v)\,\mathbf{1},\, e_I \big\rangle .
\end{aligned}
\]
Here $L^*, R^*$ denote respectively the adjoint of the left and right tensor multiplication operator on $F$. The maps $\varphi$ and $\psi$ are then extended to $T^{<\lambda}(V)$ and $T^{<\kappa}(V)$ by linearity.  \newline

Within this setting, we are now in a position to provide an alternative proof of the PDE system satisfied by the signature kernel, first derived in \cite{lemercier2024log}.
In the proof, rather than expanding the two signatures into their homogeneous components, and then extracting the terms which give the PDE, we take the rough integral equation as the primary object. The explicit characterization of each level of this equation then leads directly to the PDE.


\begin{lemma}
Let $\mathbf{X}$ and $\mathbf{Y}$ be as above. Then the signature kernel satisfies the linear PDE

\begin{equation*}
\frac{\partial^2 f}{\partial t \, \partial v}
= f \langle \boldsymbol{x}_t, \boldsymbol{y}_v \rangle
+  \phi[R^*_{\boldsymbol{x}_t}  \boldsymbol{y}_v] 
+  \psi[R^*_{\boldsymbol{y}_v} \boldsymbol{x}_t],
\end{equation*}

\begin{equation*}
\frac{\partial \phi}{\partial t}
= f \boldsymbol{x}_t
+ \phi \circ R^*_{\boldsymbol{x}_t}
+ L^*_{\psi}\boldsymbol{x}_t
- \varphi[\boldsymbol{x}_t] \mathbf{1} ,
\end{equation*}

\begin{equation*}
\frac{\partial \varphi}{\partial v}
= f \boldsymbol{y}_v
+ \psi \circ R^*_{\boldsymbol{y}_v}
+ L^*_{\phi}\boldsymbol{y}_v
- \psi[\boldsymbol{y}_v] \mathbf{1} .
\end{equation*}

\medskip

\noindent
with boundary conditions
\[
\begin{aligned}
f(a,v) &= 1, 
& \phi(a,v) &= \mathbf{0}, 
& \psi(a,v) &= \langle \mathbf{Y}_{c,v} - \mathbf{1}, \cdot \rangle, \\
f(t,c) &= 1, 
& \phi(t,c) &= \langle \mathbf{X}_{a,t} - \mathbf{1}, \cdot\rangle,
& \psi(t,c) &= \mathbf{0},
\end{aligned}
\]
where $\mathbf{0}$ is the zero in the tensor algebra.
\end{lemma}
\begin{proof}
    We begin by determining the equation for $f$. From the result above we can write $f$ as 
    \begin{align*}
    f(t,v) &= 1 + \int_{[a,t] \times [c,v]} g(K^{\mathbf{X}, \mathbf{Y}})_{s,u} d(\mathbf{X}_s, \mathbf{Y}_u).
    \end{align*}

    From \eqref{g_composed_sig}, together with the canonical identification
\[
\mathrm{Hom}\!\left(T^{< \kappa}(V)\boxtimes T^{< \lambda}(V),\mathrm{Hom}\!\left(V\boxtimes V,\mathbb{R}\right)\right)
\cong
\mathrm{Hom}\!\left(T^{\le \kappa}(V)\boxtimes T^{\le \lambda}(V),\mathbb{R}\right),
\]
it follows that, for \(\eta \in V^{\otimes j}\) and \(\xi \in V^{\otimes k}\), with
\(j \in [1:\kappa]\) and \(k \in [1:\lambda]\),
\begin{align*}
g\!\left(K^{\mathbf X,\mathbf Y}\right)_{s,u}[\eta,\xi]
&=
\begin{cases}
\langle S(\mathbf X)_{a,s}, S(\mathbf Y)_{c,u}\rangle \, \langle \eta,\xi\rangle,
& \text{if } j=k, \\[0.5em]
\displaystyle
\left\langle S(\mathbf X)_{a,s}\otimes \eta,\, S(\mathbf Y)_{c,u} \otimes \xi \right\rangle
& \text{if } j>k, \\[0.8em]
\displaystyle
\left\langle S(\mathbf X)_{a,s}\otimes \eta,\, S(\mathbf Y)_{c,u} \otimes \xi \right\rangle
& \text{if } j<k.
\end{cases}\\
&= \begin{cases}
\langle S(\mathbf X)_{a,s}, S(\mathbf Y)_{c,u}\rangle \, \langle \eta,\xi\rangle,
& \text{if } j=k, \\[0.5em]
\displaystyle
\left\langle L^*_{S(\mathbf X)_{a,s}} S(\mathbf Y)_{c,u},\,  R^*_\xi \eta \right\rangle
& \text{if } j>k, \\[0.8em]
\displaystyle
\left\langle L^*_{S(\mathbf Y)_{c,u} } S(\mathbf X)_{a,s},\,   R^*_{\eta}\xi \right\rangle
& \text{if } j<k.
\end{cases}
\end{align*}
Notice that for every non-empty word of appropriate maximum length $I, J$, with $|I| > |J|$ we have $
\left\langle S(\mathbf X)_{a,s}\otimes e_I,\, S(\mathbf Y)_{c,u} \otimes e_J\right\rangle = \psi[R^*_{e_J} e_I]$   and $\phi[R^*_{e_I} e_J] = 0$, similarly if $|I| < |J|$, $
\left\langle S(\mathbf X)_{a,s} \otimes e_I,\, S(\mathbf Y)_{c,u} \otimes e_J\right\rangle = \phi[R^*_{e_J} e_I]$ and $\varphi[R^*_{e_J} e_I] = 0$. Moreover, if $I=J$, $\psi[R^*_{e_J} e_I] = \phi[R^*_{e_I} e_J] = 0$.\newline  
This, in concert with the formula for $g\left(K^{\mathbf X,\mathbf Y}\right)$ outlined above, implies that in order to characterize $f$  we can restrict our attention to the maps
\(f, \psi\) and \(\varphi\). In fact, we have
\begin{align*}
    f(t,v) &= 1 + \int_{[a,t] \times [c,v]} g(K^{\mathbf{X}, \mathbf{Y}})_{s,u} d(\mathbf{X}_s, \mathbf{Y}_u)\\
    &= 1 + \int_a^t \int_c^v g(K^{\mathbf{X}, \mathbf{Y}})_{s,u} \left[ \boldsymbol{x}_s, \boldsymbol{y}_u\right] ds\, du\\
    &= 1+ \int_a^t \int_c^v f_{s,u} \langle \boldsymbol{x}_s, \boldsymbol{y}_u \rangle
+  \phi_{s,u}[R^*_{\boldsymbol{x}_s}  \boldsymbol{y}_u] 
+  \psi_{s,u}[R^*_{\boldsymbol{y}_u} \boldsymbol{x}_s] ds \, du.
\end{align*}
This concludes the proof for $f$.\newline
For the equations for $\phi$ and $\psi$ is sufficient to study the expressions $\langle S(\mathbf X)_{a,s} , S(\mathbf Y)_{c,u} \otimes \xi\rangle$ and $\langle S(\mathbf X)_{a,s} \otimes \eta, S(\mathbf Y)_{c,u}\rangle$ respectively. The calculations for these parts are almost identical to the ones in \cite{lemercier2024log} so we omit them.
\end{proof}

\begin{remark}
As seen above, the equations for $f, \psi$ and $\varphi$ are enough to characterize $g(K^{\mathbf{X}, \mathbf{Y}})$, $I_\mathbf{X}(g(K^{\mathbf{X}, \mathbf{Y}}))$ and $I_{\mathbf{Y}}(g(K^{\mathbf{X}, \mathbf{Y}}))$. But this implies from the same reasoning as Theorem \ref{uniqueness_sig_kernel} that they can be used to characterize all the levels in the controlled path $K^{\mathbf{X}, \mathbf{Y}}$.
\end{remark}

\begin{comment}
    and using the definition of the equation in the previous section we can write
\begin{align*}
    f(t,v) &= 1 + \int_a^t \int_c^v \sum_{i=1}^d f_{s,u} \langle \mathbf{x}^{(i)}_s, \mathbf{y}^{(i)}_u \rangle ds du + \\
    &\quad +\int_a^t \int_c^v \sum_{i=1}^d \sum_{j=i+1}^d \varphi_{s,u} [R_{\mathbf{x}^{(i)}_s} \mathbf{y}^{(i)}_u ] ds du + \int_a^t \int_c^v \sum_{j=1}^d \sum_{i=j+1}^d \psi_{s,u} [R_{\mathbf{y}^{(i)}_s} \mathbf{x}^{(i)}_u ] ds du .
\end{align*}
\end{comment}

\begin{comment}
The definition of signature allows to write for each non empty word $I = i_1...i_n$ 
\begin{align*}
    \varphi(t,v)[e_I] = \int_a^t (\varphi(s,v)[e_{i_1...i_{n-1}}]1_{n > 1} + f(s,v) 1_{n = 1} ) \mathbf{x}^{(i_n)}_s   ds, 
\end{align*}
and similarly
\begin{align*}
    \psi(t,v)[e_I] = \int_c^v (\psi(t,u)[e_{i_1...i_{n-1}}]1_{n > 1} + f(t,u) 1_{n = 1} ) \mathbf{y}^{(i_n)}_u   du. 
\end{align*}
\end{comment}


\subsection{The Schwinger-Dyson kernel}
A second example of a two-parameter RDE is given by an extension of the Schwinger--Dyson kernel equation introduced in \cite{cass2024free}. The latter arises from the study of kernels induced by path developments in the matrix Lie group of $N \times N$ anti-Hermitian matrices $\mathfrak{u}_N$. 

More precisely, consider the unitary group-valued evolution defined by
\begin{equation*}
dZ^N_t = Z^N_t \cdot M(dX_t), \qquad Z^N_0 = Id_N,
\end{equation*}
where $V \equiv \mathbb{R}^d$, $M \in \mathrm{Hom}(V, \mathfrak{u}_N)$, $X \colon [0,T] \to V$ is a path of bounded variation and $Id_N$ is the $N \times N$ identity matrix.

Define the functional $\mathcal{U}(X,M) := Z^N_T$ and let $\langle \cdot, \cdot  \rangle_{HS,N }$ denote the Hilbert-Schmidt inner product on $\mathbb{C}^{N \times N }$. One can show, under a suitable choice of probability measure $\xi_N$ on $\mathfrak{u}_N$, the existence of the limit
\begin{equation}\label{SD_kernel_definition}
K(X^1, X^2) = \lim_{N \to \infty} \frac{1}{N} \mathbb{E}\!\left[ \left\langle \mathcal{U}(X^1, M), \mathcal{U}(X^2, M) \right\rangle_{HS, N} \right],
\end{equation}
for any couple of paths with bounded variation $X^1$ and $X^2$. In the remainder of this work, we refer to $K$ as the Schwinger-Dyson kernel.

This kernel can also be characterized in two additional ways, the first one as being the solution to the quadratic functional equation
\begin{equation}\label{SD_kernel_equation_BV}
K(s,t) = 1 - \int_s^t \int_s^r K(s,u)\, K(u,r)\, \langle dX_u, dX_r \rangle 
\end{equation}
where the path $X$ is defined as the concatenation $X := X^1 * \overleftarrow{X^2}$.\newline
And the second one as being the contraction of the signature $S(X)$ against the moment sequence $\varphi(\cdot)$ of a collection of $d$-free semicircular random variables
\begin{equation}\label{expansion_K}
K(s,t) = \sum_{|L| = 0}^\infty i^{|L|} \varphi(L) S^{(L)}(X)_{s,t},
\end{equation}
where the sum is taken over all words $L$, and $i$ denotes the imaginary unit.\newline
Leveraging this representation, we can extend the definition of the Schwinger--Dyson kernel to the case where $X$ is replaced by a  geometric $p$-rough paths of arbitrarily low regularity via an approximation argument, presented in Appendix~\ref{appendix_existence_SD}. In this case equation \eqref{SD_kernel_equation_BV} reads
\begin{equation}\label{eq:uniqueness_rough_SD_1}
    K(s,t) = 1 -  \int_s^t \int_s^r f(K(s,u), K(u,r)) d\mathbf{X}_u\,d\mathbf{X}_r.
\end{equation}
where 
\begin{align*}    
f: V \times V &\to \mathrm{Hom}(V \boxtimes V), \quad f(z_1, z_2)  [\eta, \xi] = z_1 z_2 \langle \eta, \xi \rangle_V .
\end{align*}
It is now natural to ask whether this equation \eqref{eq:uniqueness_rough_SD_1}
admits a unique solution, the positive answer to this question is given in the following theorem.

\begin{theorem}
Let \(f\) be defined by
\[
f: V \times V \to \mathrm{Hom}(V \boxtimes V, \mathbb{R}), 
\qquad 
f(z_1,z_2)[\eta,\xi] := z_1 z_2 \langle \eta,\xi\rangle_V.
\]
If \(Z\) is a solution to
\begin{equation}\label{eq:uniqueness_rough_SD}
Z(s,t)
=
1-\int_s^t \int_s^r f(Z_{s,\cdot},Z)_{u,r}\,d\mathbf{X}_u\,d\mathbf{X}_r
\end{equation}
in the sense of Definition \ref{solution_2d_RDE}, then \(Z=K\). In particular, \(K\) is the unique solution to \eqref{eq:uniqueness_rough_SD} in the class \(\mathcal{D}^p_{\mathbf X}(\Delta_{[0,T]},\mathbb R^d)\). Moreover, if \(\tilde{\mathbf{X}} \in \Omega G_{p,\omega_{\mathbf{X}}}([a,b],V)\) and $\tilde K$ is the solution of \eqref{eq:uniqueness_rough_SD} driven by this rough path, then
\[
\|K-\tilde K\|_{\infty}
\leq
C_{p,q,\omega_{\mathbf{X}},\omega_{\mathbf{Y}}}
\|\mathbf{X}-\tilde{\mathbf{X}}\|_{p}.
\]
\end{theorem}
\begin{proof}
Fix $s \in [0, T]$, and $\tilde{p} \in (p, \kappa  +1 )$. The first step consists in proving local uniqueness of \eqref{eq:uniqueness_rough_SD} on a two simplex $\Delta_{[s_0,t_0]}$, where $t_0$ is sufficiently close to $s_0$.\newline
This will be obtained by using the stability properties of the solution map $\Gamma$ associated to equation \eqref{eq:uniqueness_rough_SD}
\begin{align*}
&\Gamma: B(\Delta_{[s_0, t_0 ]}) \to \mathcal{D}^p_{\mathbf{X}}(\Delta_{[s_0, t_0]}, \mathbb{R}),   \\ 
& \Gamma(Q)_{s,t} = 1- I^{\Delta_{[s, t]}}_{\mathbf{X}}(f(Q_{s, \cdot}, Q)),
\end{align*}
where
\begin{align*}
B(\Delta_{[s_0, t_0]}) := \left \{ Q \in \mathcal{D}^{\tilde{p}}_{\mathbf{X}}(\Delta_{[s_0, t_0]}, \mathbb{R}) \, : \, Q_{r,r} =  K_{r,r}  \text{ for every }  r \in [s_0, t_0],   \,  \|Q\|_{ \mathcal{D}^{\tilde{p}}_{ \mathbf{X}, \mathbf{X}}} \leq M \right \},
\end{align*}
here $M$ is an appropriately chosen constant, that we will fix later in the proof. The choice of $K_{r,r}$ as the initial condition is justified from the fact that every fixed point for the map $\Gamma$ will have the same initial condition.\newline
Assume $\tilde{K}$ is another fixed point for $\Gamma$. We will show that $\tilde K = K$ on $\Delta_{[s_0,t_0]}$.\newline
Define
\(\overline{f}(s,u,r) := f(K_{s,\cdot}, K)_{u,r} - f(\tilde{K}_{s, \cdot}, \tilde{K})_{u,r}\), which can be viewed as a two-parameter controlled path in the variables 
$(u,r)$, taking values in the space of real-valued, singly controlled rough paths in the variable 
$s$.\newline
From analogous calculations as the one performed in Lemma \ref{proposition_L_U_controlled} we have 
that 
\begin{equation}\label{formula_K_K}
\begin{aligned}
    & \left(K(s,t) - \tilde K(s,t) \right)\left[ \eta \boxtimes \xi \right]\\
    &= \int_s^t \int_u^t \overline f(s, u, r) d\mathbf{X}_r d\mathbf{X}_u \left[ \eta \right] - \int_s^t \overline f(s,s,r) d\mathbf{X}_r \circ (id_{T^{<\kappa}(V)}  \boxtimes \pi_{\geq 1}) \left[\Delta \eta\right]\\
    & \quad + \overline f (s,s,s) \circ (id_{T^{<\kappa}(V)} \boxtimes \pi_{\geq 1} \boxtimes \pi_{\geq 1}) \left[\Delta^2 \eta\right]  + \int_s^t \overline f (s,u,t)d\mathbf{X}_u \left[ \eta \boxtimes \xi \right] \\
    &\quad - \overline f(s,s,t)\circ (id_{T^{<\kappa}(V)}  \boxtimes \pi_{\geq 1} \boxtimes id_{T^{<\kappa}(V)}) \left[\Delta \eta \boxtimes \xi\right] + \overline f(s,t,t) \circ (id_{T^{<\kappa}(V)}  \boxtimes \pi_{\geq 1} \boxtimes \pi_{\geq 1})[\eta \boxtimes \Delta \xi],
\end{aligned}
\end{equation}
here in the second integral we are considering the map $\overline f(s,s,r)$  to be a one parameter controlled integral in $r$ taking values in the space of singly controlled rough integrals in the variable $s$. Similarly for the third integral, where $\overline f(s,u,t)$ is to be understood as a one parameter controlled path in $u$ taking values in the space of jointly controlled rough paths in the variables $(s,t)$.\newline

Starting from the trace, thanks to the formula \eqref{formula_K_K}, we know that the remainder in the second parameter of $K - \tilde K$ satisfies the following identity
\begin{align*}
R^{2}_{z,t}(s)[\mathbf{1}, \mathbf{1}] &= \int_z^t \int_s^u \overline f(s,u,r) d\mathbf{X}_r d\mathbf{X}_u - \int_s^z \overline f(s,u,z) d\mathbf{X}_u \left[\mathbf{X}^{\ge 1}_{z,t}\right] - \overline f(s,z,z) \left[\mathbf{1} \boxtimes \mathbf{X}^{\geq 1}_{z,t} \boxtimes \mathbf{X}^{\geq 1}_{z,t}\right]\\
&\quad +\int_s^z \overline f(s,u,z) d\mathbf{X}_u \left[\mathbf{X}^{(\kappa)}_{z,t}\right] +  \overline f(s,z,z) \left[\mathbf{1} \boxtimes \mathbf{X}^{(\kappa)}_{z,t} \boxtimes \mathbf{X}^{\geq 1}_{z,t}\right] + \overline f(s,z,z) \left[\mathbf{1} \boxtimes \mathbf{X}^{\geq 1}_{z,t} \boxtimes \mathbf{X}^{(\kappa)}_{z,t}\right],
\end{align*}
where we used the Fubini theorem for rough integrals over a simplex introduced in Section \ref{Simplex_section} to write the first integral.\newline
The regularity of the last three terms and the sewing lemma for one-parameter rough integrals (Section 2 in \cite{cass2022combinatorial}) used jointly with stability of the composition of two-parameter controlled paths with smooth functions (Proposition \ref{stability_proposition}), yields 
\[
\left|R^{2}_{z,t}(s)[\mathbf{1}, \mathbf{1}] \right| \leq C d(K, \tilde{K})\omega_{\mathbf{X}}(z,t)^{p/\kappa},
\]
where $\theta>1$ is as defined in Theorem \ref{stability_theorem_2D}, and $C$ denotes a positive constant depending only on $g,p,\omega_{\mathbf{X}},\|K\|_{\mathcal{D}_{\mathbf{X},\mathbf{X}}},$ and $\|\tilde K\|_{\mathcal{D}_{\mathbf{X},\mathbf{X}}}$. Throughout the proof, $C$ may refer to a different positive constant depending on the same parameters.\newline
The other remainder is 
\begin{align*}
    R^{1}_{w,s}(t)[\mathbf{1},\mathbf{1}]= & \int_s^t \int_u^t \overline f(s,u,r) d\mathbf{X}_r d\mathbf{X}_u  - \int_w^t \int_u^t \overline f(w,u,r) d\mathbf{X}_r d\mathbf{X}_u \left[\mathbf{X}^{< \kappa}_{w,s}\right]\\ 
    & \quad  
    - \int_w^t \overline f(w,w,r) d\mathbf{X}_r \left[\mathbf{X}_{w,s}^{< \kappa }\boxtimes \mathbf{X}^{\ge 1}_{w,s}\right] + \overline f(w,w,w)\left[ \mathbf{X}^{< \kappa}_{w,s} \boxtimes \mathbf{X}^{\ge 1}_{w,s} \boxtimes \mathbf{X}^{\ge 1}_{w,s} \right]\\
    & \quad  
    + \int_w^t \overline f(w,w,r) d\mathbf{X}_r \left[\mathbf{X}_{w,s}^{< \kappa }\boxtimes \mathbf{X}^{(\kappa)}_{w,s}\right] - \overline f(w,w,w)\left[ \mathbf{X}^{< \kappa}_{w,s} \boxtimes \mathbf{X}^{(\kappa)}_{w,s} \boxtimes \mathbf{X}^{\ge 1}_{w,s} \right]\\
    &\quad - \overline f(w,w,w)\left[ \mathbf{X}^{< \kappa}_{w,s} \boxtimes \mathbf{X}^{\ge 1}_{w,s}\boxtimes \mathbf{X}^{(\kappa)}_{w,s}  \right]
\end{align*}
The same sewing lemma argument applies for this term yielding 
\[
\left|R^{1}_{w,s}(t)[\mathbf{1}, \mathbf{1}] \right| \leq C d(K, \tilde{K})\omega_{\mathbf{X}}(w,s)^{\theta}.
\]
For the remainder of remainders for $K - \tilde{K}$, using the definition and identity \eqref{formula_K_K} again, we have the identity
\begin{align*}
& \mathbf{R} \begin{pmatrix} w,s \\ z,t \end{pmatrix}[\mathbf{1}, \mathbf{1}]\\
&= \int_z^t \int_s^u \overline f(s,u,r) d\mathbf{X}_r d\mathbf{X}_u - \int_s^z \overline f(s,u,z) d\mathbf{X}_u \left[\mathbf{X}^{\ge 1}_{z,t}\right] + \int_s^z \overline f(s,u,z) d\mathbf{X}_u \left[\mathbf{X}^{(\kappa)}_{z,t}\right]\\
& \quad - \overline f(s,z,z) \left[\mathbf{1} \boxtimes \mathbf{X}^{\geq 1}_{z,t} \boxtimes \mathbf{X}^{\geq 1}_{z,t}\right] + \overline f(s,z,z) \left[\mathbf{1} \boxtimes \mathbf{X}^{(\kappa)}_{z,t} \boxtimes \mathbf{X}^{\geq 1}_{z,t}\right] + \overline f(s,z,z) \left[\mathbf{1} \boxtimes \mathbf{X}^{\geq 1}_{z,t} \boxtimes \mathbf{X}^{(\kappa)}_{z,t}\right]\\
& \quad - \int_z^t \int_w^u \overline f(w,u,r) d\mathbf{X}_r d\mathbf{X}_u \left[\mathbf{X}^{< \kappa}_{w,s}\right] 
+ \int_z^t  \overline f(w,w,r) d\mathbf{X}_r  \left[\mathbf{X}^{< \kappa}_{w,s} \boxtimes \mathbf{X}^{\geq 1}_{w,s}\right]\\
&\quad -\int_z^t  \overline f(w,w,r) d\mathbf{X}_r  \left[\mathbf{X}^{< \kappa}_{w,s} \boxtimes \mathbf{X}^{(\kappa)}_{w,s}\right] + \int_w^z \overline f(w,u,z) d\mathbf{X}_u [ \mathbf{X}^{< \kappa}_{w,s} \boxtimes \mathbf{X}^{\ge 1}_{z,t}]\\
&\quad 
- \int_w^z \overline f(w,u,z) d\mathbf{X}_u [ \mathbf{X}^{< \kappa}_{w,s} \boxtimes \mathbf{X}^{(\kappa)}_{z,t}]
- \overline{f}(w,w,z)\left[ \mathbf{X}^{< \kappa}_{w,s} \boxtimes \mathbf{X}^{\ge 1}_{w,s} \boxtimes \mathbf{X}^{\ge 1}_{z,t}\right]\\
& \quad + \overline{f}(w,w,z)\left[ \mathbf{X}^{< \kappa}_{w,s} \boxtimes \mathbf{X}^{(\kappa)}_{w,s} \boxtimes \mathbf{X}^{\ge 1}_{z,t}\right] +  \overline{f}(w,w,z)\left[ \mathbf{X}^{< \kappa}_{w,s} \boxtimes \mathbf{X}^{\ge 1}_{w,s} \boxtimes \mathbf{X}^{(\kappa)}_{z,t}\right]\\
& \quad +\overline f(w,z,z) \left[\mathbf{X}^{< \kappa}_{w,s} \boxtimes \mathbf{X}^{\geq 1}_{z,t} \boxtimes \mathbf{X}^{\geq 1}_{z,t}\right] -\overline f(w,z,z) \left[\mathbf{X}^{< \kappa}_{w,s} \boxtimes \mathbf{X}^{(\kappa)}_{z,t} \boxtimes \mathbf{X}^{\geq 1}_{z,t}\right]\\
& \quad -\overline f(w,z,z) \left[\mathbf{X}^{< \kappa}_{w,s} \boxtimes \mathbf{X}^{\geq 1}_{z,t} \boxtimes \mathbf{X}^{(\kappa)}_{z,t}\right] .
\end{align*}
By rearranging the terms above we can express $\mathbf{R}$ as
\begin{align*}
\mathbf{R} \begin{pmatrix} w,s \\ z,t \end{pmatrix}[\mathbf{1}, \mathbf{1}]
&= A\begin{pmatrix} w,s \\ z,t \end{pmatrix} + B\begin{pmatrix} w,s \\ z,t \end{pmatrix} + D_1\begin{pmatrix} w,s \\ z,t \end{pmatrix}\\
& \quad  + D_2\begin{pmatrix} w,s \\ z,t \end{pmatrix} + D_3\begin{pmatrix} w,s \\ z,t \end{pmatrix} + D_4\begin{pmatrix} w,s \\ z,t \end{pmatrix}
\end{align*}
where 
\begin{align*}
A\begin{pmatrix} w,s \\ z,t \end{pmatrix} &:= \int_z^t \int_s^u  \overline f(s,u,r) d\mathbf{X}_r d\mathbf{X}_u  - \int_z^t \int_s^u \overline f(w,u,r) d\mathbf{X}_r d\mathbf{X}_u \left[\mathbf{X}^{< \kappa}_{w,s} \right]  - \int_s^z  \overline f(s,u,z)  d\mathbf{X}_u \left[\mathbf{X}^{\ge 1}_{z,t}\right] \\
&\quad+ \int_s^z  \overline f(w,u,z)  d\mathbf{X}_u \left[ \mathbf{X}^{< \kappa}_{w,s}\boxtimes \mathbf{X}^{\ge 1}_{z,t}  \right] - \overline f(s,z,z) \left[\mathbf{1} \boxtimes \mathbf{X}^{\geq 1}_{z,t} \boxtimes \mathbf{X}^{\geq 1}_{z,t}\right] \\
& \quad + \overline f(w,z,z) \left[\mathbf{X}^{< \kappa}_{w,s} \boxtimes \mathbf{X}^{\geq 1}_{z,t} \boxtimes \mathbf{X}^{\geq 1}_{z,t}\right],\\
B\begin{pmatrix} w,s \\ z,t \end{pmatrix}&:= - \int_z^t \int_w^s \overline f(w,u,r) d\mathbf{X}_r d\mathbf{X}_u \left[\mathbf{X}_{w,s}\right]  + \int_z^t  \overline f(w,w,r) d\mathbf{X}_r  \left[\mathbf{X}^{< \kappa}_{w,s} \boxtimes \mathbf{X}^{\geq 1}_{w,s}\right] 
\\
&\quad+ \int_w^s  \overline f(w,w,r) d\mathbf{X}_r  \left[\mathbf{X}_{w,s} \boxtimes \mathbf{X}^{\geq 1}_{w,s}\right] + \int_w^s \overline f(w,u,z) d\mathbf{X}_u [ \mathbf{X}^{< \kappa}_{w,s} \boxtimes \mathbf{X}^{\ge 1}_{z,t}] \\
&\quad - \overline{f}(w,w,z)\left[ \mathbf{X}_{w,s} \boxtimes \mathbf{X}^{\ge 1}_{w,s} \boxtimes \mathbf{X}^{\ge 1}_{z,t}\right],\\
D_1\begin{pmatrix} w,s \\ z,t \end{pmatrix} &:= \int_s^z \overline f(s,u,z) d\mathbf{X}_u \left[\mathbf{X}^{(\kappa)}_{z,t}\right] - \int_s^z \overline f(w,u,z) d\mathbf{X}_u \left[\mathbf{X}^{< \kappa}_{w,s} \boxtimes \mathbf{X}^{(\kappa)}_{z,t}\right] + \overline f(s,z,z) \left[\mathbf{1} \boxtimes \mathbf{X}^{\geq 1}_{z,t} \boxtimes \mathbf{X}^{(\kappa)}_{z,t}\right]\\
&\quad -\overline f(w,z,z) \left[\mathbf{X}^{< \kappa}_{w,s} \boxtimes \mathbf{X}^{\geq 1}_{z,t} \boxtimes \mathbf{X}^{(\kappa)}_{z,t}\right],\\
D_2\begin{pmatrix} w,s \\ z,t \end{pmatrix} &:= - \int_w^s \overline f(w,u,z) d\mathbf{X}_u \left[\mathbf{X}^{< \kappa}_{w,s} \boxtimes \mathbf{X}^{(\kappa)}_{z,t}\right] + \overline{f}(w,w,z)\left[ \mathbf{X}^{< \kappa}_{w,s} \boxtimes \mathbf{X}^{\ge 1}_{w,s} \boxtimes \mathbf{X}^{(\kappa)}_{z,t}\right],\\
D_3\begin{pmatrix} w,s \\ z,t \end{pmatrix} &:= -\int_z^t  \overline f(w,w,r) d\mathbf{X}_r  \left[\mathbf{X}^{< \kappa}_{w,s} \boxtimes \mathbf{X}^{(\kappa)}_{w,s}\right] +  \overline{f}(w,w,z)\left[ \mathbf{X}^{< \kappa}_{w,s} \boxtimes \mathbf{X}^{(\kappa)}_{w,s} \boxtimes \mathbf{X}^{\ge 1}_{z,t}\right],\\
D_4\begin{pmatrix} w,s \\ z,t \end{pmatrix} &:=  \overline f(s,z,z) \left[\mathbf{1} \boxtimes \mathbf{X}^{(\kappa)}_{z,t} \boxtimes \mathbf{X}^{\geq 1}_{z,t}\right] -\overline f(w,z,z) \left[\mathbf{X}^{< \kappa}_{w,s} \boxtimes \mathbf{X}^{(\kappa)}_{z,t} \boxtimes \mathbf{X}^{\geq 1}_{z,t}\right]
\end{align*}
For the term $A$ we can conclude via a one-parameter sewing argument in the parameter $t$ and stability of composition of two-parameter controlled paths with smooth functions that 
\[
\left|A\begin{pmatrix} w,s \\ z,t \end{pmatrix}\right|
 \leq C d(K, \tilde{K}) \omega_{\mathbf{X}}(w,s)^{\theta} \omega_{\mathbf{X}}(z,t)^{\theta}.\]
For the term $B$ we can use expression (4.13) in \cite{cass2023fubini} and stability of compositions with smooth functions to conclude that 
\[
\left|B\begin{pmatrix} w,s \\ z,t \end{pmatrix}\right|
 \leq C d(K, \tilde{K}) \omega_{\mathbf{X}}(w,s)^{\theta} \omega_{\mathbf{X}}(z,t)^{\theta}.\]
For the term $D_i, i \in [1:3]$, by using the one-parameter sewing lemma and the regularity of $\mathbf{X}^{(\kappa)}$ we recover the estimate 
\[
\left|D_i\begin{pmatrix} w,s \\ z,t \end{pmatrix}\right|
 \leq C d(K, \tilde{K}) \omega_{\mathbf{X}}(w,s)^{p/\kappa} \omega_{\mathbf{X}}(z,t)^{p/\kappa}.\]
For the term $D_4$, the definition of one-parameter controlled path and the regularity of $\mathbf{X}^{(\kappa)}$ allow to recover
\[
\left|D_4\begin{pmatrix} w,s \\ z,t \end{pmatrix}\right|
 \leq C d(K, \tilde{K}) \omega_{\mathbf{X}}(w,s)^{p/\kappa} \omega_{\mathbf{X}}(z,t)^{p/\kappa}.\]
The higher-level terms are handled by analogous estimates, which give
\begin{align*}
    d(K,\tilde{K}) \leq C
    d(K,\tilde{K})  \omega_{\mathbf{X}}(s_0, t_0)^{\frac{1}{p} - \frac{1}{\tilde{p}}}.
\end{align*}
The last term $D_4$ the desired bound 
This is enough to show uniqueness in the simplex $\Delta_{[s_0, t_0]}$ for $t_0-s_0$ small enough. We will denote this unique solution as $Q^*$.\newline
Now consider a value $M > \|Q^*\|_{\mathcal{D}^{\tilde{p}}_{ \mathbf{X}, \mathbf{X}}}$ and consider 
the set
\begin{align*}
&B(\Delta_{[s_0, t_0 + \delta]}, Q^*)\\
&:= \{ Q \in \mathcal{D}^{\tilde{p}}_{\mathbf{X}}(\Delta_{[s_0, t_0 + \delta]}, \mathbb{R}) \, : \, Q_{|\Delta_{[s_0, t_0]}} = Q^*, Q_{r,r} =  K_{r,r}  \text{ for every } r \in [s_0,t_0 + \delta],\, \|Q\|_{ \mathcal{D}^{\tilde{p}}_{ \mathbf{X}, \mathbf{X}}} \leq M \}.
\end{align*}
On such set we define the map
\begin{align*}
\Gamma &: B(\Delta_{[s_0, t_0 + \delta]}, Q^*) \longrightarrow \mathcal{D}^p_{\mathbf{X}}(\Delta_{[s_0, t_0 + \delta]}, \mathbb{R}), \\[4pt]
\Gamma(Q)_{s,t} &=
\begin{cases}
Q^*_{s,t}, & \text{if } (s,t) \in \Delta_{[s_0, t_0]}, \\[6pt]
1 -  I^{\Delta_{[s, t]}}_{\mathbf{X}}(f(Q_{s, \cdot},Q)) & \text{otherwise.}
\end{cases}
\end{align*}
Assume $\tilde{K} \in B(\Delta_{[s_0, t_0 + \delta]}, Q^*)$ is another fixed point for $\Gamma$, then we get
\begin{align*}
&\Gamma(K)_{s,t} - \Gamma(\tilde K)_{s,t}\\
&=
\begin{cases}
\mathbf{0}, & \text{if } (s,t) \in \Delta_{[s_0, t_0]}, \\[6pt]
\displaystyle
 I^{\Delta_{[s,t]}}_{\mathbf{X}}(f(K_{s, \cdot},K) - f(\tilde{K}_{s, \cdot},\tilde{K})) & \text{otherwise.}
\end{cases}
\end{align*}
Focusing on the second case,
we either have $\Delta_{[s,t]}$ or $s > t_0$ or $s < t_0$. The first case produces the estimate  
\begin{align*}
    d(K,\tilde{K}) \leq C
    d(K,\tilde{K})  \omega_{\mathbf{X}}(t_0, t_0+\delta)^{\frac{1}{p} - \frac{1}{\tilde{p}}},
\end{align*}
that can be verified via a totally analogous argument as the one present before in this proof. This is enough to conclude the uniqueness argument in this case.\newline
For the second case, we can use the decomposition $\Delta_{[s,t]} = \Delta_{[s, t_0]} + [t_0, t] \times [t_0, t] + \Delta_{[t_0, t]}$ to
write
\[ \Gamma(K)_{s,t} - \Gamma(\tilde K)_{s,t} = 
 I^{\Delta_{[t_0, t]}}_{\mathbf{X}}(f(K_{s, \cdot},K) - f(\tilde{K}_{s, \cdot},\tilde{K})) +
 I^{{[t_0, t] \times [t_0, t]}}_{\mathbf{X}, \mathbf X}(f(K_{s, \cdot},K) - f(\tilde{K}_{s, \cdot},\tilde{K})), 
\]
where the superscript ${[t_0, \cdot] \times [t_0, \cdot ]}$ is used in the second integral to clarify the domain of integration.
The identity above, in concert with now familiar estimates used previously in this proof, allows to conclude that for $s > t_0$ we have
\begin{align*}
    d(K, \tilde{K}) \leq C
    d(K, \tilde{K})  \omega_{\mathbf{X}}(t_0, t_0 + \delta)^{\frac{1}{p} - \frac{1}{\tilde{p}}}.
\end{align*}
Then by choosing $\delta$ sufficiently small, the inequality above is enough to show uniqueness on the simplex $\Delta_{[s_0, t_0 + \delta]}$. 
\newline
This is enough to conclude: given any two  solutions $K$ and $\tilde{K}$ with $\|K\|_{\mathcal{D}^{\tilde p}_{\mathbf{X}, \mathbf{X}}}, \|\tilde{K}\|_{\mathcal{D}^ {\tilde p}_{\mathbf{X}, \mathbf{X}}} \le M$, we can keep the radius of the ball fixed and iterate the argument to cover the simplex $\Delta_{[s_0, t]}$ for any $t$. Since the argument does not depend on the choice of $s_0$, the proof of uniqueness is complete. The stability estimate now follows from a close variant to the proof presented above. Here, by choosing a partition \(\mathcal P=\{s_i\}_i\) of \([0,T]\) with sufficiently small mesh size, one can show that if \(K\) and \(\tilde K\) denote the solutions of the Schwinger--Dyson kernel equation driven by \(\mathbf X\) and \(\tilde{\mathbf X}\), respectively, then
\[
|K - \tilde{K}|_{\infty, [s_i, s_{i+1}] \times [s_j, s_{j+1}] \cap \Delta_{[0, T]}} \leq C \left( |K - \tilde{K}|_{\infty, [0, s_i] \times [0, s_j] \cap \Delta_{[0, T]}} + \|\mathbf{X} - \tilde{\mathbf{X}}\|_{p; [s_i, s{i+1}]} + \|\mathbf{X} - \tilde{\mathbf{X}}\|_{p; [s_j, s{j+1}]} \right).
\]
The bound above immediately leads to the desired inequality, concluding the proof.
\end{proof}


\begin{comment}
\begin{proposition}
\begin{align}\label{SD_equation_double}
K_X(s,t) &=  1 - \sum_{i=1}^d \int_s^t \int_s^r K_X(s,u) K_X(u,r)  dX^\iota_u dX^\iota_r\\
&= 1 - \sum_{\iota=1}^d \int_s^t \int_s^r K_X(u,r) K_X(r,t)  dX^\iota_u dX^\iota_r. 
\end{align}
\end{proposition}
\begin{proof}
Following the derivation in Theorem 3.10 in \cite{cass2024free} 
    \begin{align*}
        K_X(s,t) &= \sum_{|\mathbf{I}| = 0}^\infty i^{|\mathbf{I}|} \varphi(\mathbf{I}) S^{\mathbf{I}}_{s,t} \\
        &=  1 + \sum_{\iota =1}^d \sum_{|\mathbf{I}| = 0}^\infty i^{|\mathbf{I}| + 1} \varphi(\mathbf{I}\iota) S^{\mathbf{I}\iota}_{s,t} \\
        &=  1 + \sum_{\iota=1}^d \sum_{|\mathbf{K}|, |\mathbf{L}| = 0}^\infty i^{|\mathbf{K}|, |\mathbf{L}| + 2} \varphi(\mathbf{K}) \varphi(\mathbf{L}) S^{\mathbf{K}\iota\mathbf{L}\iota}_{s,t}. 
    \end{align*}
    Now consider $\mathbf{L} = \ell_1 \dots \ell_k$
    \begin{equation}\label{signature_expansion}
    \begin{aligned}S^{\mathbf{K}\iota\mathbf{L}\iota} &= \int_{s< u<u_1<\dots < u_k<r<t} S^{\mathbf{K}}_{s,u} dX^\iota_{u} dX^{\ell_1}_{u_1} \dots dX^{\ell_k}_{u_k} dX^\iota_r \\ 
        &= \int_{s< u<r<t} S^{\mathbf{K}}_{s,u} dX^\iota_{u}  \int_{u<\ell_1< \dots< \ell_k < r}dX^{\ell_1}_{u_1} \dots dX^{\ell_k}_{u_k} dX^\iota_r\\
        &= \int_{s<u<r<t } S^{\mathbf{K}}_{s,u} S^{\mathbf{L}}_{u,r} dX^\iota_u dX^\iota_r ,
    \end{aligned}
    \end{equation}
    so that we can write 
    \begin{equation}\label{SD_classic_epression}
    \begin{aligned}
         K_X(s,t) &= 1 + \sum_{\iota=1}^d \sum_{|\mathbf{K}|, |\mathbf{L}| = 0}^\infty i^{|\mathbf{K}|, |\mathbf{L}| + 2} \varphi(\mathbf{K}) \varphi(\mathbf{L}) \int_{s<u<r<t } S^{\mathbf{K}}_{s,u} S^{\mathbf{L}}_{u,r} dX^\iota_u dX^\iota_r \\
         &= 1 - \sum_{i=1}^d \int_s^t \int_s^r K_X(s,u) K_X(u,r)  dX^\iota_u dX^\iota_r . 
    \end{aligned}
    \end{equation}
    On the other hand 
    \begin{align*}
        K_X(s,t) &= \sum_{|\mathbf{I}| = 0}^\infty i^{|\mathbf{I}|} \varphi(\mathbf{I}) S^{\mathbf{I}}_{s,t} \\
        &=  1 + \sum_{j=1}^d \sum_{|\mathbf{I}| = 0}^\infty i^{|\mathbf{I}| + 1} \varphi(j\mathbf{I}) S^{j\mathbf{I}}_{s,t} \\
        &=  1 + \sum_{j=1}^d \sum_{|\mathbf{K}|, |\mathbf{L}| = 0}^\infty i^{|\mathbf{K}|, |\mathbf{L}| + 2} \varphi(\mathbf{K}) \varphi(\mathbf{L}) S^{j\mathbf{K}j\mathbf{L}}_{s,t}. 
    \end{align*}
Similarly to \eqref{signature_expansion}, we can write
\begin{align*}
    S^{j\mathbf{K}j\mathbf{L}}_{st} = \int_{t > r > u>s} S^\mathbf{K}_{rt} dX^i_r S^\mathbf{L}_{ur} dX^i_u. 
\end{align*}
This implies that
\begin{align*}
         K_X(s,t) &= 1 + \sum_{\iota=1}^d \sum_{|\mathbf{K}|, |\mathbf{L}| = 0}^\infty i^{|\mathbf{K}|+ |\mathbf{L}| + 2} \varphi(\mathbf{K}) \varphi(\mathbf{L}) \int_{t>r>u>s } S^{\mathbf{K}}_{u,r} S^{\mathbf{L}}_{r,t} dX^\iota_u dX^\iota_r \\
         &= 1 - \sum_{\iota=1}^d \int_s^t \int_s^r K_X(u,r) K_X(r,t)  dX^\iota_u dX^\iota_r . 
    \end{align*}
\end{proof}


Let $NC(n)$ be the set partial matching of length $n$ that can be generated by the following recursive relation
\begin{align*}
NC(1) &= 1\\
NC(n) &= NC(n-1) n + NC(n-2)
\end{align*}
and for an element $\pi \in NC(n) $ define  $P(\pi)$ to be the set of pairs in $\pi$ and $U(\pi)$ the set of unmatched pairs in $\pi$.
\end{comment}
\subsubsection{The smooth rough path case}\label{sec:smooth_RP_SD}

In the last part of this section, we propose a method for associating a integro-differential equation to the Schwinger–Dyson kernel equation under the additional assumption that $\mathbf{X}$ is a smooth rough path.
The method is based on a combinatorial relationship between the   the trace component of $K$ and the remaining components of the controlled path.
Recall that we can characterize all  the components of the path $K$ via the identity
\begin{equation}\label{gub_representation}
K^{(I, J)}
(s, t) = (-1)^n\sum_{|L| = 0}^\infty i^{
|L|+n + m} \varphi(JLI) S^{(L)}
(\mathbf{X})_{s, t},
\end{equation}
where $I = i_1 \dots i_n$ and  $J = j_1 \dots j_m$ are possibly empty words over an alphabet of size $d$.
The identity above appears in the proof of Theorem \ref{theorem_existence}.
 
We first establish this combinatorial relationship in the case where $X$ is of bounded variation. The general case of a kernel driven by a path $\mathbf{X} \in \Omega G_{p,\omega}([0,T],V)$, which need not be a smooth rough path, then follows by a standard approximation argument, almost identical to that presented in Appendix \ref{appendix_existence_SD}.\newline

Before proving the following combinatorial formula, we introduce, for notational convenience, the set $NC(n)$ of non-crossing partial matchings of length $n$, defined recursively by
\begin{align*}
NC(1) &= \bigl\{\{\{1\}\}\bigr\},\\
NC(n) &= \left\{ \pi \cup \{\{n\}\} : \pi \in NC(n-1) \right\} \\
&\qquad \cup
\left\{ (\pi \setminus \{\{m\}\}) \cup \{\{m,n\}\} : \pi \in NC(n-1),\ m = \max U(\pi) \right\},
\end{align*}
where \(U(\pi)\) denotes the set of unmatched elements of \(\pi\). We also write \(P(\pi)\) for the set of pairs of \(\pi\). If \(U(\pi)=\{U_1<\cdots<U_r\}\), we list its unmatched points in increasing order.


We also recall \cite[Lemma 5.4.7, Remark 5.4.8]{anderson2010introduction} that, for every pair of words $I,J$ constructed from the alphabet $\{1,\dots,d\}$, the moment sequence $\varphi$ satisfies the identities
\begin{align}
\varphi(IJ) &= \varphi(JI), \nonumber\\
\varphi(Ii) &= \sum_{I = K i J} \varphi(K)\varphi(J), \label{SD_prop2}
\end{align}
which are known as the Schwinger–Dyson equations and give the name to the kernel presented in this section.
By the identity above we can see that the formula \eqref{gub_representation} is symmetric in the two parameters, up to a factor  $-1^{|I|}$. Therefore it suffices to derive the formula for the higher levels of $K$ in the second variable; the corresponding identity in the first variable follows by symmetry. For this reason in the next two results we treat $K$ as a one-parameter controlled path.

We are now ready to present the combinatorial relationship that constitutes the first building block of the integro-differential equation associated to the Schwinger-Dyson kernel equation.

\begin{proposition}\label{Gubinelli derivatives}
If $X$ has bounded variation, then for every non-empty word $I = i_1\dots i_n$, $n\leq \kappa-1$:
    \[
K^{(I)}
=
\sum_{\pi \in NC(n)}
\left(
\prod_{(p,q) \in P(\pi)} (-1) \delta_{i_p,i_q}
\right)
K_{\,i_{U_1},\ldots,i_{U_r}},
\]
with the convention $K^{(\emptyset)} = K_\emptyset$ and for a non-empty word $J = j_1 \dots j_m$
\begin{equation}\label{index_removal}
K_{J}(s,t) =   (-1)^{m} \int_{s < u_1 < \dots < u_m <t} K(s, u_1)K(u_1, u_2) \dots K(u_{m-1}, u_{m})K(u_{m}, t) dX^{(j_1)}_{u_1}\dots dX^{(j_m)}_{u_m}.
\end{equation}
\end{proposition}
\begin{proof}
We prove the identity by induction on $I$. 
The base case $I = \emptyset$ follows by definition. \newline
Notice that the case $|I|=1$ follows from equation \eqref{SD_prop2}, in fact, following Theorem 3.10 in \cite{cass2024free} we obtain
\begin{align*}
K(s,t) &= \sum_{|L|=0}^\infty i^{|L|}\varphi(L)S^{(L)}_{s,t} \\
&= 1 + \sum_{|L_1|, |L_2|=0}^\infty i^{|L_1|+|L_2| + 2} \varphi(L_1) \varphi(L_2) S^{(L_1 i_1 L_2 i_1)}_{s,t} \\
&= 1 - \sum_{i=1}^d \int_s^t \int_s^r K(s,u)K(u, r) dX^{(i)}_{u} dX^{(i)}_{r}.
\end{align*}
From the identity above we see that
\[
K^{(i)}(s,t) = -\int_s^t K(s,u)K(u, t) dX^{(i)}_u, 
\]
so that the case $|I|=1$ is satisfied.\newline

For the general induction step, we begin by considering a non-empty word $J = j_1 \dots j_m$ and the definition
\[
K_{J}(s,t) = (-1)^{m} \int_{s < u_1 < \dots < u_m <t} K(s, u_1)K(u_1, u_2) \dots K(u_{m}, t) dX^{(j_1)}_{u_1}\dots dX^{(j_m)}_{u_m},
\]
and apply the identity found for the induction basis to $K(u_m, t)$. From doing this we recover
\begin{align*}
K_{J}(s,t) &= (-1)^{m}
   \int_{s<u_1<\cdots<u_m<t}
   K(s,u_1)\cdots K(u_{m-1},u_{m})\,
   dX^{(j_1)}_{u_1}\cdots dX^{(j_m)}_{u_m} \\
&\quad
   + (-1)^{m+1}
   \sum_{j_{m+1}=1}^d
   \int_s^t
   \int_{s<u_1<\cdots<u_m<u_{m+1}<r}
   K(s,u_1)\cdots K(u_m,u_{m+1}) K(u_{m+1},u) \\
&\hspace{6em}
   dX^{(j_1)}_{u_1}\cdots dX^{(j_{m+1})}_{u_{m+1}}
   \, dX^{(j_{m+1})}_u \\
&=
   - \int_s^t
     K_{j_1\dots j_{m-1}}(s,u)\,
     dX^{(j_m)}_u
   + \sum_{j_{m+1}=1}^d
     \int_s^t
     K_{j_1\dots j_{m+1}}(s,u)\,
     dX^{(j_{m+1})}_u.
\end{align*}
Notice how the case $J = \emptyset$ reduces to the base case just analysed.\newline
Now, consider the combinatorial formula holds for any word $I = i_1...i_n$, plugging the above identities for $K_J$ into the induction hypothesis, and after appropriately relabelling the letters forming $J$, we obtain
\begin{align*}
K^{(I)}
&=
\sum_{\pi \in NC(n)}
\Bigg(
\prod_{(p,q)\in P(\pi)} (-1)\delta_{i_p,i_q}
\Bigg)
\Bigg[
 - \int_s^t
   K_{i_{U_1}\dots i_{U_{r-1}}}(s,u)\,
   dX^{(i_{U_r})}_u  \\
&\qquad\qquad
 + \sum_{i_{U_{r+1}}=1}^d
   \int_s^t
   K_{i_{U_1}\dots i_{U_{r+1}}}(s,u)\,
   dX^{(i_{U_{r+1}})}_u
\Bigg].
\end{align*}


The non-matching partial matchings of order $n+1$ can be constructed from $NC(n)$. In fact, according to the definition for $NC(n+1)$, this can be done for every  $\pi \in NC(n)$ by either leaving the last index unpaired or pairing it with the last unmatched index. This implies that we can rewrite the integrand with respect to $X^{(i_{n+1})}$ in the previous expression as
    \[
K^{(Ii_{n+1})}
=
\sum_{\pi \in NC(n+1)}
\left(
\prod_{(p,q) \in P(\pi)}  (-1)\delta_{i_p,i_q}
\right)
K_{\,i_{U_1},\ldots,i_{U_r}},
\]
concluding the proof.
\end{proof}


\begin{remark}\label{remark_smooth_numerical}
Observe that, by the definition in \eqref{index_removal} and the representation \eqref{gub_representation}, it follows immediately that for every word $I$, $K_I \in \mathcal{D}^p_{\mathbf{X}}(\Delta_{[0,T]}, \mathbb{R})$.
    This implies that the argument in previous proposition can be used to compute the higher levels of this term, $K^{^{(i_{n+1} \dots i_\ell)}}_{i_1 \dots i_n}$ for an arbitrary word $i_1 \dots i_n$ and $i_{n+1} \dots i_\ell$ with adequate maximum length. We have
    \begin{align*}
        K_{i_1,\ldots,i_{n}}^{(i_{n+1} \dots i_\ell)} = \sum_{\pi \in NC(\ell)_n}
\left(
\prod_{(p,q) \in P(\pi)} (-1) \delta_{i_p,i_q}
\right)
K_{\,i_{U_1},\ldots,i_{U_m}},
    \end{align*}
where \(NC(\ell)_n\) denotes the set of non-crossing partial matchings on \(\{1,\dots,\ell\}\) such that no pair joins two indices in \(\{1,\dots,n\}\).
\end{remark}
Following the convention introduced for the components of a controlled rough path, we denote by $K_n$ the level-wise analogue of the coordinate-wise notation $K_I$, where $n \in \mathbb{N}$ and $I$ is a word of length $|I| = n$. More precisely 
\[
K_n: \Delta_{[0, T]} \to \mathrm{Hom}(V^{\otimes n}, \mathbb{R}), \qquad K_n(s,t)[e_I] = K_I(s,t)\, \text{for }|I|=n .
\]
\newline
To distinguish the jointly controlled rough path $K$ from the auxiliary process with coordinates $K_I$ we will use the notation $K_\bullet$ to denote the latter.\newline
The proposition above and the subsequent remark lead immediately to the next proposition, that gives us the system used to construct the integro-differential equation associated to the Schwinger-Dyson kernel equation when $\mathbf{X}$ is a smooth rough path.

\begin{proposition}\label{numerical_scheme_rep}
For every \(n\in [1:\kappa+\zeta]\), let \(\ell=\lceil n/2\rceil\). Then the homogeneous components of the Schwinger--Dyson kernel satisfy
\begin{equation}\label{alternative_representation_levelwise}
K_0(s,t)=1 +  \int_s^t  K_{1}(s,u) d\mathbf{X}_u du,
\qquad
K_n(s,t)
=
-\int_s^t \sigma_n\bigl(K_{\ell-1}(s, \cdot),K_{n-\ell}(\cdot, t)\bigr)_u\,d\mathbf X_u .
\end{equation}
Here, for every word \(I=i_1\cdots i_n\) over the alphabet \(\{1,\dots,d\}\), we set
\(
\tilde I:=i_1\cdots i_{\ell-1},\,
\hat I:=i_{\ell+1}\cdots i_n,
\)
and,given a basis \(e_1,\dots,e_d\) of \(V\), define
\begin{align*}
&\sigma_n:
\mathrm{Hom}\bigl(V^{\otimes(\ell-1)},\mathbb R\bigr)\boxtimes
\mathrm{Hom}\bigl(V^{\otimes(n-\ell)},\mathbb R\bigr)
\to
\mathrm{Hom}\bigl(V,\mathrm{Hom}(V^{\otimes n},\mathbb R)\bigr),\\
&\sigma_n(\eta,\xi)[e_j](e_I)
=
\sigma_n(\eta,\xi)[e_j](e_{\tilde I}\otimes e_{i_\ell}\otimes e_{\hat I})
=
\eta(e_{\tilde I})\,\delta_{i_\ell j}\,\xi(e_{\hat I}),
\end{align*}
for every word \(I\) of length \(n\), and extend linearly on \(V^{\otimes n}\).
\end{proposition}

\begin{proof}
The first identity follows from the same approximation argument used in the proof of Theorem~\ref{theorem_existence}.

For the second identity, we first assume that $X$ has bounded variation. Let $J = j_1 \dots j_n$ be a non-empty word. By the definition of $K_J$ in Proposition~\ref{Gubinelli derivatives}, we have
\begin{align*}
K_{J}(s,t) 
&= (-1)^{n} \int_{s < u_1 < \dots < u_n < t} 
K(s, u_1) K(u_1, u_2) \dots K(u_{n-1}, u_n) K(u_n, t)\,
dX^{(j_1)}_{u_1} \dots dX^{(j_n)}_{u_n}.
\end{align*}
Setting $\ell = \lceil n/2 \rceil$, and using the definitions of 
$K_{j_1 \dots j_{\ell-1}}(s,u)$ and $K_{j_{\ell+1} \dots j_n}(u,t)$, this expression can be rewritten as
\[
K_{J}(s,t)
= - \int_s^t K_{j_1 \dots j_{\ell-1}}(s,u)\,
K_{j_{\ell+1} \dots j_n}(u,t)\, dX^{(j_\ell)}_u.
\]

The general case $\mathbf{X} \in \Omega G_p([0, T], V)$ then follows by a standard approximation argument.
\end{proof}

With these results at hand we are now ready to finalise the construction of the integro-differential equation.\newline
We begin by fixing a non-negative integer $\zeta$, which determines the level of truncation in the tensor algebra $T^{\le \kappa + \zeta}(V)$ where the solution $K$ takes values. To ensure consistency of the scheme, we require $\zeta$ to satisfy
\[
\max\left(\floor{\frac{\kappa + \zeta}{2}}, 1\right) + \kappa - 1 \le \kappa + \zeta,
\]
which is equivalent to
\[
\zeta \ge
\begin{cases}
\kappa - 3, & \text{if } \kappa \ge 4,\\
0, & \text{otherwise}.
\end{cases}
\]

As we will see, this choice guarantees that the system of equations presented later in this part of the work is closed. \newline
Proposition \ref{Gubinelli derivatives} and the following remark allows us to define the linear maps that will constitute the vector fields appearing in the equation
\[
\begin{aligned}
    \psi&: \Delta_{[0, T]} \to \mathrm{Hom}( T^{\le \ceil{\frac{ \kappa + \zeta}{2}}-1}(V), \mathrm{Hom}(T^{\le \kappa}, \mathbb{R})), \\
    \phi&: \Delta_{[0, T]} \to \mathrm{Hom}( T^{\le \kappa + \zeta - \ceil{\frac{ \kappa + \zeta}{2}}}(V), \mathrm{Hom}(T^{\le \kappa}, \mathbb{R})),
\end{aligned}
\]
where, for every word $I$, $J$ of admissible lengths, the following holds for a fixed basis of $V$
\begin{align*}
\psi_{s,t}(e_I)[e_J]  &= 
\sum_{\pi \in NC(IJ)_{|I|}}
\left(
\prod_{(p,q) \in P(\pi)} (-1) \delta_{i_p,i_q}
\right)
K(s,t)[{e_{\,i_{U_1},\ldots,i_{U_m}}}],\\
\phi_{s,t}(e_I)[e_J]  &= (-1)^{|J|}
\sum_{\pi \in NC(IJ)_{|I|}}
\left(
\prod_{(p,q) \in P(\pi)} (-1) \delta_{i_p,i_q}
\right)
K(s,t)[{e_{\,i_{U_1},\ldots,i_{U_m}}}].  
\end{align*}

Under the hypothesis that $\mathbf{X}$ is a smooth path, we can then use the expressions above to rephrase 
equation \eqref{alternative_representation_levelwise} as 
\begin{equation*}
\begin{aligned}
K_{0}(s,t) 
&= 1 +  \int_s^t  K_{1}(s,u) [\boldsymbol{x}_u]  du,\\
K_{n}(s,t) 
&= -  \int_s^t \sigma_n( \psi_{s, \cdot}, \phi_{\cdot, t})_u  [\boldsymbol{x}_u]  du.
\end{aligned}
\end{equation*}

From these identities, by taking derivatives in the two parameters, we can then recover the following result
\begin{lemma}
  Let $\mathbf{X} \in \Omega G([0,T], V)$ be a smooth rough path. Then the Schwinger-Dyson kernel equation satisfies the integro-differential equation
  \begin{align*}
      \frac{\partial^2 K_{0}}{\partial s \partial t} &=  \frac{\partial K_{1}(s,t) }{\partial s} [\boldsymbol{x}_t], \\
      \frac{\partial^2 K_{n}}{\partial s \partial t} &= \frac{\partial \sigma_n( \psi_{s, \cdot}, \phi_{\cdot, t})_s}{\partial t}\,[\boldsymbol{x}_s] - \frac{\partial \sigma_n( \psi_{s, \cdot}, \phi_{\cdot, t})_t}{\partial s} [\boldsymbol{x}_t] \\
      &\qquad - \int_s^t \frac{\partial^2 \sigma_n( \psi_{s, \cdot}, \phi_{\cdot, t})_u}{\partial s \partial t} [\boldsymbol{x}_u] du,
  \end{align*}
  with conditions $K_{\bullet}(s,s) = id_{T^{\leq \kappa + \zeta}(V)}$ for every $s \in [0, T]$.
\end{lemma}

\begin{remark}
It is now clear that extending the dimension of the tensor algebra in which the solution takes values is unavoidable in this system. Indeed, in order to close the system, one must be able to compute the equations satisfied by the components at level $\kappa + \zeta$.
\end{remark}

\section{Numerical method for Schwinger-Dyson kernel}\label{numerical_method_section}

\subsection{Methodology}
The setup constructed in the previous section, alongside Propositions \ref{Gubinelli derivatives} and \ref{numerical_scheme_rep}, provides a natural framework for constructing a numerical scheme for the Schwinger-Dyson kernel equation. This scheme is defined on the grid $\mathcal{P}^2 \cap \Delta_{[s,t]}$, where $\mathcal{P}$ is a finite partition of $[0,T]$ and $\mathcal{P}^2 := \mathcal{P} \times \mathcal{P}$. In what follows, for any indices \(i,j\) such that \((t_i,t_j)\in \mathcal{P}^2 \cap \Delta_{[s,t]}\), we will use the shorthand $K_\bullet(i,j) = K_\bullet(t_i, t_j)$. We recall that the initial conditions are
\[
K_{\emptyset}(i,i) = 1, \qquad K_n(i,i) = 0 \quad \text{for all } i \in \mathcal{P},\, n \in [1: \kappa +  \zeta],
\]
where $\zeta$ is defined as in subsection \ref{sec:smooth_RP_SD}. Let $D$ the total dimension of the tensor algebra $T^{\le \kappa + \zeta}(V) \equiv T^{\le \kappa + \zeta}(\mathbb{R}^d)$.\newline
To update the value $K_\bullet(i,j+1)$, we can leverage the identities in Proposition~\ref{numerical_scheme_rep}. Explicitly, for the level-$0$ component, we set
\begin{align} \label{trace_numerical}
K_\emptyset(i,j+1)
= 1 + \sum_{|\omega| = 1}^{\kappa} \sum_{m=i+1}^{j+1} 
K^{(\omega)}(i,m)\, \mathbf{X}^{(\omega)}_{m-1,m},
\end{align}
where the first sum is to be taken over words $\omega$ from an alphabet of size $d$. 

For $I = i_1 \dots i_n$, $|I| \geq 1$, with $\ell = \ceil{n/2}$, $\tilde{I} = i_1 \dots i_{\ell -1}$, and $\hat{I} = i_{\ell + 1} \dots i_{n}$, we have
\[
K_{I}(i,j+1)
= \sum_{m=i+1}^{j+1}
\sum_{|\omega| = 0}^{\kappa-1}
\sum_{(\alpha,\beta)\in Sh^{-1}(\omega)}
(-1)^{|\beta| + 1}
K^{(\alpha)}_{\tilde{I}}(i,m)\,
K^{(\beta)}_{\hat{I}}(m,j+1)\,
\mathbf{X}^{(i_\ell\omega)}_{m-1,m},
\]
where $K^{(\alpha)}_I, K^{(\beta)}_{\hat{I}}$ are defined as in Remark \ref{remark_smooth_numerical} and the operation \(Sh^{-1}(\omega)\) denotes the inverse shuffle of the word \(\omega\). More precisely, for a word \(I=i_1\cdots i_n\), we define
\[
Sh^{-1}(I)
:=
\{(J,L):\ I\in J \Delta^{*} L\},
\]
where \(\Delta^{*}\) denotes the product dual to the shuffle coproduct defined in Section \ref{Simplex_section}.\newline
The expression above can be further decomposed as
\begin{align}
K_{I}(i,j+1)
&= \sum_{m=i+1}^{j}
\sum_{|\omega| = 0}^{\kappa-1}
\sum_{(\alpha,\beta)\in Sh^{-1}(\omega)}
(-1)^{|\beta| + 1}
K^{(\alpha)}_{\tilde{I}}(i,m)\,
K^{(\beta)}_{\hat{I}}(m,j+1)\,
\mathbf{X}^{(i_\ell\omega)}_{m-1,m}
\label{term_b}\\
&\quad + \sum_{|\omega| = 0}^{\kappa-1}
\sum_{(\alpha,\beta)\in Sh^{-1}(\omega)}
(-1)^{|\beta| + 1}
K^{(\alpha)}_{\tilde{I}}(i,j+1)\,
K^{(\beta)}_{\hat{I}}(j+1,j+1)\,
\mathbf{X}^{(i_\ell\omega)}_{j,j+1}.
\label{term_A}
\end{align}

Using the combinatorial identities of Proposition~\ref{Gubinelli derivatives} and the subsequent remark to rewrite $K^{(\omega)}, K^{(\alpha)}_I$ and $K^{(\beta)}_{\hat{I}}$ in terms of elements of $K_\bullet$, the relations \eqref{trace_numerical}, \eqref{term_b} and \eqref{term_A} can be restated in the compact form
\[
(Id_D - A_j) K_{\bullet} = b_{i,j},
\]
so that the numerical solution is obtained by solving a linear system where $b_{i,j} \in \mathbb{R}^D$ and $A_j \in \mathbb{R}^{D \times D}$. Note that the matrix $A_j$ does not depend on $i$, due to the initial conditions imposed on the scheme.\newline

We conclude by presenting the runtime and memory costs of the algorithm, which we derive in detail in Appendix~\ref{appendix:algorithmic_complexity}.
Let \(M>N\in\mathbb N\). We assume that the driving path is sampled on a fine dyadic grid with \(2^M\) points, while the numerical scheme is implemented on a coarser dyadic grid with \(2^N\) points. Therefore, each increment of the rough path on the coarse grid is obtained by combining \(2^{M-N}\) fine-grid increments of the sampled path. Finally let $\tilde{D}$ be the total dimension of the tensor algebra $T^{\le \kappa }(\mathbb{R}^d)$. Then, the algorithm above has runtime complexity
\[
\mathrm{Cost}_{\mathrm{runtime}} \leq C\left(
2^{3N}\,(2d)^{\kappa + \zeta -1} (d+1)^\kappa \;+\; 2^{2N} D^3 + \tilde{D}(2^N - 1)d^{\tilde{D}} 2^{M-N}
\right),
\]
and memory complexity
\[
\mathrm{Memory} \leq C\left(2^{2N} D\right),
\]
where \(C>0\) is a constant independent of all the parameters appearing above.

\subsection{Numerical experiments}
In this part of the work, we perform numerical experiments using the methods developed above.

The first experiment investigates how different choices of \(\kappa\) affect the performance of the algorithm as the regularity of the driving signal decreases. For a given choice of \(\kappa\), we denote the corresponding method by \(SDK\kappa\).

We also include in the comparison the randomized CDE method for the Schwinger--Dyson kernel, denoted by \(SDKr\), introduced in \cite{cass2024free} for the computation of the kernel in the bounded-variation setting. The parameters used for this method are matrix dimension \(200\) and \(250\) Monte Carlo simulations.

Simulated paths of three-dimensional fractional Brownian motion on \([0,1]\) with various Hurst indices \(H\) serve as input to the different numerical methods. Since no closed-form solution is available, the discrepancies reported below should be interpreted as empirical measures of agreement between numerical schemes, rather than as errors with respect to an exact reference solution. More precisely, for each of \(50\) independent sample paths of the driving noise, each consisting of \(4096\) piecewise-interpolated increments, we compute the corresponding numerical solutions and, for each pair of methods, evaluate the terminal-time discrepancy
\[
\left|K^{(\emptyset),\text{method 1}}_{0,1}-K^{(\emptyset),\text{method 2}}_{0,1}\right|.
\]
We then report the mean absolute error and the standard deviation of this quantity over the \(50\) sample paths. For the methods \(SDK2\) and \(SDK3\), the rough path is constructed over blocks of \(32\) increments. This procedure provides a simple terminal-time diagnostic for assessing inter-method consistency and for quantifying how decreasing input regularity affects the discrepancy between schemes.

We emphasize that some of the schemes considered below are tested on driving signals whose regularity lies outside the regime in which the corresponding methods are theoretically justified. In such cases, these schemes are used only as heuristic numerical benchmarks. Therefore, the observed agreement or discrepancy between methods should be interpreted as an empirical comparison, and not as a conclusion supported by the available convergence theory.

\begin{table}[H]
\centering
\begin{tabular}{c|c|c|c}
\hline
$H$ & Method pair & MAE & STD \\
\hline
\multirow{6}{*}{0.85} 
 & SDK3 vs SDK2 & $2.02\times 10^{-11}$ & $1.76\times 10^{-11}$ \\
 & SDK3 vs SDK1 & $5.95\times 10^{-6}$  & $1.21\times 10^{-6}$ \\
 & SDK3 vs SDKr & $2.38\times 10^{-5}$  & $4.85\times 10^{-6}$ \\
 & SDK2 vs SDK1 & $5.95\times 10^{-6}$  & $1.21\times 10^{-6}$ \\
 & SDK2 vs SDKr & $2.38\times 10^{-5}$  & $4.85\times 10^{-6}$ \\
 & SDK1 vs SDKr & $1.79\times 10^{-5}$  & $3.64\times 10^{-6}$ \\
\hline
\multirow{6}{*}{0.50} 
 & SDK3 vs SDK2 & $2.68\times 10^{-9}$  & $5.97\times 10^{-10}$ \\
 & SDK3 vs SDK1 & $1.84\times 10^{-4}$  & $1.22\times 10^{-5}$ \\
 & SDK3 vs SDKr & $7.36\times 10^{-4}$  & $4.89\times 10^{-5}$ \\
 & SDK2 vs SDK1 & $1.84\times 10^{-4}$  & $1.22\times 10^{-5}$ \\
 & SDK2 vs SDKr & $7.36\times 10^{-4}$  & $4.89\times 10^{-5}$ \\
 & SDK1 vs SDKr & $5.52\times 10^{-4}$  & $3.67\times 10^{-5}$ \\
\hline
\multirow{6}{*}{0.255} 
 & SDK3 vs SDK2 & $2.91\times 10^{-7}$  & $4.41\times 10^{-8}$ \\
 & SDK3 vs SDK1 & $1.98\times 10^{-3}$  & $1.57\times 10^{-4}$ \\
 & SDK3 vs SDKr & $7.96\times 10^{-3}$  & $6.32\times 10^{-4}$ \\
 & SDK2 vs SDK1 & $1.98\times 10^{-3}$  & $1.57\times 10^{-4}$ \\
 & SDK2 vs SDKr & $7.96\times 10^{-3}$  & $6.32\times 10^{-4}$ \\
 & SDK1 vs SDKr & $5.97\times 10^{-3}$  & $4.75\times 10^{-4}$ \\
\hline
\end{tabular}
\caption{Pairwise MAE and STD between the numerical methods for different Hurst parameters \(H\).}
\label{tab:method_errors_updated}
\end{table}

From Table~\ref{tab:method_errors_updated}, we observe that, as the regularity of the input signal decreases, both the MAE and the STD between the different methods increase. We also note that the discrepancy between \(SDK3\) and \(SDK2\) is smaller than the discrepancy between either of these methods and the remaining ones.

Moreover, when the Hurst index decreases from \(0.5\) to \(0.255\), the pair \(SDK3\) vs \(SDK2\) exhibits the largest relative increase in both MAE and STD, as can be seen by comparing the corresponding ratios in Table~\ref{tab:method_errors_updated}. This behaviour is consistent with the fact that the schemes \(SDK_\kappa\) have increasing approximation order. Heuristically, this suggests that the discrepancy between \(SDK2\) and \(SDK1\) should be larger than that between \(SDK3\) and \(SDK2\).

Moreover, the methods \(SDK_\kappa\) are designed to handle kernels driven by paths of regularity strictly less than \(\kappa+1\). Therefore, as the regularity decreases, the discrepancy between successive schemes is expected to increase, which is consistent with the more pronounced effect observed as \(H\) approaches \(0.255\).

Below, we also plot the runtime cost for a kernel driven by three-dimensional Brownian motion, first as the grid size increases and then as the block size increases. The numerical scheme used in these simulations is \(SDK2\).

\begin{figure}[H]
  \centering
  \begin{subfigure}[t]{0.48\textwidth}
    \centering
    \includegraphics[width=\linewidth]{plots/runtime_scaling_length.png}
    \caption{Runtime as a function of the grid size \(2^M\).}
    \label{fig:test1}
  \end{subfigure}%
  \hfill
  \begin{subfigure}[t]{0.48\textwidth}
    \centering
    \includegraphics[width=\linewidth]{plots/runtime_scaling_block.png}
    \caption{Runtime as a function of \(2^{M-N}\), the number of fine-grid increments used to construct each rough-path increment.}
    \label{fig:test2}
  \end{subfigure}
  \caption{Computational times for a kernel driven by three-dimensional Brownian motion, obtained by varying (a) the grid size and (b) the block size.}
  \label{fig:overall}
\end{figure}

We observe that, in accordance with our estimates, the runtime scales exponentially with both the grid size and the block size.

The next plot, in which we fix the grid size to \(M=12\) and the block size to \(M-N=5\), illustrates the polynomial scaling of the runtime complexity as the dimension \(d\) increases. The numerical scheme used in this experiment is again \(SDK2\).

\begin{figure}[H]
\centering
\includegraphics[scale=0.5]{plots/runtime_scaling_dimension} 
\caption{Runtime as a function of the dimension \(d\) for a kernel driven by Brownian motion.}
\label{fig:runtime_dimension}
\end{figure}



\appendix
\section{Technical results on two-parameter controlled paths}\label{appendix_technical}
We begin this section by recalling two of the main results in \cite{cass2023fubini}.
\begin{theorem}[Two-parameter rough integral on rectangular domains]\label{existence_2D_theorem}
Let 
\[
Z \in \mathcal{D}^{p, q}_{\mathbf{X}, \mathbf{Y}}\big([a, b] \times [c,d], \mathrm{Hom}(V \boxtimes W, F)\big).
\] 
Then the limit
\[
\int_{[a,b] \times [c,d]} Z_{s, u} \, d(\mathbf{X}, \mathbf{Y})  
   :=  \lim_{\left|\mathcal{P} \times \mathcal{D}\right| \to 0} 
   \sum_{\mathcal{P}\times \mathcal{D}} \Omega^Z
\]
exists and satisfies the estimate
\begin{align*}
\Bigg|
\int_{[a,b]\times[c,d]} Z_{s,u}\, d(\mathbf{X}_s,\mathbf{Y}_u)
- \Omega^Z\begin{pmatrix} a, b \\ c, d \end{pmatrix}
\Bigg|
&\le C_{p,q,T}\Big[
      A(\theta,T,\mathbf{R})\,
      \omega_{\mathbf{X}}(a,b)^{\theta}\,
      \omega_{\mathbf{Y}}(c,d)^{\theta} \\
&\qquad
      + B_X(c,d)\,\omega_{\mathbf{X}}(a,b)^{\theta}
      + B_Y(a,b)\,\omega_{\mathbf{Y}}(c,d)^{\theta}
\Big].
\end{align*}

where $\theta = \min(\frac{\floor{p} + 1}{p}, \frac{\floor{q} + 1}{q})$ and 
\begin{align*}
&A(p, q, \mathbf{R}) 
= \left( \sup_{j \in [0:\kappa)} \| \mathbf{X}^{(j+1)} \|_{p_j} \right)
  \left( \sup_{k \in  [0:\lambda)} \| \tilde{\mathbf{X}}^{(k+1)} \|_{\tilde{p}_k} \right)
  \left( \sup_{(j,k)\in [0:\kappa) \times [0:\lambda)} \| \mathbf{R}^{(j,k)} \|_{(q_j,\tilde{q}_k)} \right),\\
&B_{X}(c,d)
= 
\sup_{(j,k)\in [0:\kappa) \times [0:\lambda)}\left[ 
  \|\mathbf{X}^{(j+1)}\|_{\frac{p}{j+1}}\,
  \|\mathbf{Y}^{(k+1)}\|_{\frac{q}{k+1}}\,
  \big\| R^{1, (j,k)}
       \big\|_{\frac{p}{\kappa - j}}\,
  \omega_{\mathbf{Y}}(c,d)^{\frac{k+1}{q}}
\right],\\
&B_{Y}(a,b)
= \sup_{(j,k)\in [0:\kappa) \times [0:\lambda)}
\left[
   \|\mathbf{X}^{(j+1)}\|_{\frac{p}{j+1}}\,
  \|\mathbf{Y}^{(k+1)}\|_{\frac{q}{k+1}}\,
  \big\| R^{2,(j,k)}\!
       \big\|_{\frac{q}{\lambda - k}}\,
  \omega_{\mathbf{X}}(a,b)^{\frac{j+1}{p}}
\right].
\end{align*}



\end{theorem}




\begin{theorem}[Stability of two-parameter rough integrals]\label{stability_theorem_2D}
Let $\mathbf{X}, \tilde{\mathbf{X}} \in  \Omega G_{p, \omega_{\mathbf{X}}}([a,b], V)$ and $\mathbf{Y}, \tilde{\mathbf{Y}} \in  \Omega G_{q, \omega_\mathbf{Y}} ([c,d], W)$. Consider the controlled paths
\[
Z \in \mathcal{D}^{p, q}_{\mathbf{X}, \mathbf{Y}}\big([a, b] \times [c,d], F\big), \, \tilde{Z}\in \mathcal{D}^{p, q}_{\tilde{\mathbf{X}}, \tilde{\mathbf{Y}}}\big([a, b] \times [c,d], F\big).
\] 
Then the following estimate holds
\begin{align*}
&\left| \int_{[a,b] \times [c,d]} Z_{s, u} \, d(\mathbf{X}_s, \mathbf{Y}_u) 
      - \int_{[a,b] \times [c,d]} \tilde{Z}_{s, u} \, d(\tilde{\mathbf{X}}_s, \tilde{\mathbf{Y}}_u) \right|\\
&\qquad \leq C_{\omega_\mathbf{X}, \omega_\mathbf{Y}, p, q, Z, \tilde{Z}} 
\, \Big( D_{\omega_{\mathbf{X}}, \omega_{\mathbf{Y}}} d( Z, \tilde{Z} ) + \|\mathbf{X} - \tilde{\mathbf{X}}\|_{p} + \|\mathbf{Y} - \tilde{\mathbf{Y}}\|_{q} \Big),
\end{align*}
where $D_{\omega_{\mathbf{X}}, \omega_{\mathbf{Y}}} := \omega_{\mathbf{X}}(a,b)^{1/p} \omega_{\mathbf{Y}}(c,d)^{1/q}$.

\end{theorem}


Below we present the proof of Propositions \ref{stability_under_composition} and \ref{stability_proposition}.\newline
We only treat the case $j>0$ and $k>0$, since the remaining cases follow
directly from the classical one-parameter theory of controlled rough paths for  the first level remainders and are much simpler to prove for the second order remainders. To improve readability we replace $id_{T^{\kappa}(V)}$ and $id_{T^{\lambda}(W)}$ by $id$ whenever it is clear from the context. Moreover we introduce the notation $\Delta^{\boxtimes}_{>0} := \widetilde{\Delta \boxtimes \Delta}$ .  We also allow the remainder terms to change from line to line. \newline
\begin{proof}[Proof of Proposition 
\ref{stability_under_composition}]
Notice that, from the symmetry of $D^\ell g$, we can write
\begin{align}
&\sum_{\ell=1}^{\lambda+\kappa-2} \frac{1}{\ell!}
D^\ell g\left(Z^{(0,0)}_{t,v}\right)
\circ Z^{\otimes \ell}_{t,v}
\circ (\Delta^{\boxtimes}_{>0})^\ell \label{original_expression_expansion_g}\\
&=
\sum_{i = 1}^{\kappa -1} \sum_{m = 0}^{\lambda -1} \frac{1}{i!} \frac{1}{m!} D^{i+m}g(Z^{(0,0)}_{t,v}) \circ Z^{\otimes i + m}_{t,v} \circ \left( (\tilde{\Delta} \boxtimes \Delta)^{i} \otimes ( id \boxtimes \tilde{\Delta})^{m} \right). \nonumber
\end{align}
Expanding $(Z_{t,v})^{\otimes i+m}$ in the first parameter yields 
\begin{align}
&\sum_{i = 1}^{\kappa -1} \sum_{m = 0}^{\lambda -1} \frac{1}{i!} \frac{1}{m!} D^{i+m}g(Z^{(0,0)}_{t,v}) \circ (Z_{s,v} \circ [L_{\mathbf{X}^{<\kappa}_{s,t}} \otimes id])^{\otimes i + m} \circ \left((\tilde{\Delta} \boxtimes \Delta))^{i} \otimes ( id \boxtimes \tilde{\Delta})^{m} \right) + A_{s,t}(v),
\end{align}
where the term $A_{s,t}(v)$ absorbs all the terms in the expansion that contain at least a term $R^1_{s,t}(v)$.\newline
We now expand the derivatives of $g$ around $Z^{(0,0)}_{s,v}$ in the first term above.
This yields
\begin{align*}
&\sum_{i = 1}^{\kappa -1} \sum_{n = i}^{\kappa -1} \sum_{m = 0}^{\lambda -1} \frac{1}{n!} \frac{1}{m!} \binom{n}{i} D^{n+m}g(Z^{(0,0)}_{s,v})\\
& \qquad \circ (Z_{s,v} \circ [L_{\mathbf{X}^{<\kappa}_{s,t}} \otimes id])^{\otimes i + m} \circ \left((\tilde{\Delta} \boxtimes \Delta))^{i} \otimes ( id \boxtimes \tilde{\Delta})^{m} \otimes (Z^{(0,0)}_{t,v} - Z^{(0,0)}_{s,v})^{\otimes n-i}\right) + B_{s,t}(v).
\end{align*}
Here $B_{s,t}(v)$ collects the Taylor expansion remainder for the expansion above.\newline
Expanding the difference $Z^{(0,0)}_{t,v} - Z^{(0,0)}_{s,v}$, we obtain
\begin{align*}
&\sum_{i = 1}^{\kappa -1} \sum_{n = i}^{\kappa -1} \sum_{m = 0}^{\lambda -1} \frac{1}{n!} \frac{1}{m!} \binom{n}{i} D^{n+m}g(Z^{(0,0)}_{s,v})\\
& \qquad \circ (Z_{s,v})^{n+m} \circ \left([L_{\mathbf{X}^{<\kappa}_{s,t}} \otimes id]^{\otimes i + m} \circ \left((\tilde{\Delta} \boxtimes \Delta))^{i} \otimes ( id \boxtimes \tilde{\Delta})^{m} \right) \otimes \left[\mathbf{X}_{s,t} \boxtimes \mathbf{1} \right]^{\otimes n-i}\right) + C_{s,t}(v).
\end{align*}
Where the term $C$ groups the terms in which at least a term $R^{1, (0,0)}$ appears.\newline
From the estimate
\begin{align*}
&\Bigg\| \sum_{n = 1}^{\kappa -1}  \frac{1}{n!}  D^{n+m}g(Z^{(0,0)}_{s,v})\circ (Z_{s,v})^{n} \circ\\
&\qquad \left(\sum_{i = 1}^{n}\binom{n}{i} [L_{\mathbf{X}^{<\kappa}_{s,t}} \otimes id]^{\otimes i } \circ (\tilde{\Delta} \boxtimes \Delta)^{i}  \otimes \left[\mathbf{X}_{s,t} \boxtimes \mathbf{1} \right]^{\otimes n-i} - (\tilde{\Delta} \boxtimes \Delta)^{n}  \otimes [L_{\mathbf{X}^{<\kappa}_{s,t}} \otimes id]^{\otimes n } \right)[\eta_j, \xi_k]\Bigg\|\\
&\leq C \omega_\mathbf{X}(s,t)^{(\kappa - j)/p} \|\eta_j\| \|\xi_k \|,
\end{align*}
valid for every $m \leq \lambda -1$, $\eta_j \in V^{\otimes j}$, $\xi_k \in W^{\otimes k}$ and exploiting the symmetry of the derivatives of g, the previous identity reduces to
\begin{align}
&\sum_{n = 1}^{\kappa -1} \sum_{m = 0}^{\lambda -1} \frac{1}{n!} \frac{1}{m!} D^{n+m}g(Z^{(0,0)}_{s,v}) \circ (Z_{s,v} \circ [L_{\mathbf{X}^{<\kappa}_{s,t}} \otimes id])^{\otimes n + m} \circ \left((\tilde{\Delta} \boxtimes \Delta))^{n} \otimes ( id \boxtimes \tilde{\Delta})^{m} \right) + E_{s,t}(v).
\end{align}
So that we can rewrite the original expression \eqref{original_expression_expansion_g} as 
\[
\sum_{\ell=1}^{\lambda+\kappa-2} \frac{1}{\ell!}
D^\ell g\left(Z^{(0,0)}_{s,v}\right)
\circ Z^{\otimes \ell}_{s,v}
\circ (\Delta^{\boxtimes}_{>0})^\ell \circ \left[L_{\mathbf{X}^{< \kappa}_{s,t}} \boxtimes id \right] + R^{g, 1}_{s,t}(v)\\
\]
The remainder $R^{g, 1}_{s,t}(v)$ collects all higher-order terms arising from:
(i) the remainder in the controlled expansion of $Z$,
(ii) the Taylor remainder of $g$, (iii) the remainder in the expansion for $Z^{(0,0)}_{t,v} - Z^{(0,0)}_{s,v}$
and (iv) the truncation error in the shuffle expansion. Its explicit form is
given by
\begin{align*}
R^{g, 1}_{s,t}(v)
&= A_{s,t}(v) + B_{s,t}(v) + C_{s,t}(v) + E_{s,t}(v),
\end{align*}
where
\begin{align*}
A_{s,t}(v)
&:=
\sum_{\ell=1}^{\lambda+\kappa-2}
\sum_{i=0}^{\ell}
\frac{1}{\ell!}
D^\ell g\left(Z^{(0,0)}_{t,v}\right)
\circ
\left(
Z^{\otimes(i-1)}_{s,v} \circ \left[L_{\mathbf{X}^{< \kappa}_{s,t}} \boxtimes id \right]
\otimes
R^{1}_{s,t}(v)
\otimes
Z^{\otimes(\ell-i)}_{t,v}
\right)
\circ
(\Delta^{\boxtimes}_{>0})^\ell,
\\[1ex]
B_{s,t}(v)
&:=
\sum_{i = 1}^{\kappa -1} \sum_{m = 0}^{\lambda -1} \frac{1}{(\kappa-1)!} \frac{1}{m!} \binom{\kappa-1}{i} \int_0^1 (1-r)^\kappa D^{\kappa+m}g((1-r)Z^{(0,0)}_{s, v} + rZ^{(0,0)}_{t, v}) dr\\
& \qquad \circ \left((Z_{s,v} \circ [L_{\mathbf{X}^{<\kappa}_{s,t}} \otimes id])^{\otimes i + m} \circ \left((\tilde{\Delta} \boxtimes \Delta)^{i} \otimes ( id \boxtimes \tilde{\Delta})^{m}\right) \otimes (Z^{(0,0)}_{t,v} - Z^{(0,0)}_{s,v})^{\otimes \kappa-i}\right),
\\[1ex]
C_{s,t}(v) &:= \sum_{i = 1}^{\kappa -1} \sum_{n = i}^{\kappa -1} \sum_{m = 0}^{\lambda -1} \sum_{p=0}^{n-i} \frac{1}{n!} \frac{1}{m!} \binom{n}{i} D^{n+m}g(Z^{(0,0)}_{s,v})\\
& \qquad \circ (Z_{s,v} \circ [L_{\mathbf{X}^{<\kappa}_{s,t}} \otimes id])^{\otimes i + m} \\
&\qquad \circ \left((\tilde{\Delta} \boxtimes \Delta))^{i} \otimes ( id \boxtimes \tilde{\Delta})^{m} \otimes \left(Z_{s,v} \circ \left[\mathbf{X}^{\geq 1 }_{s,t} \boxtimes id \right]\right)^{\otimes p-1} \otimes R^{1, (0,0)}_{s,t}(v) \otimes \left(Z^{(0,0)}_{t,v} - Z^{(0,0)}_{s,v}\right)^{\otimes n-i-p} \right) \\[1ex]
E_{s,t}(v)
&:=
\sum_{i = 1}^{\kappa -1} \sum_{n = i}^{\kappa -1} \sum_{m = 0}^{\lambda -1}  \frac{1}{n!} \frac{1}{m!} \binom{n}{i} D^{n+m}g(Z^{(0,0)}_{s,v})\circ Z^{\otimes m+n}_{s,v} \circ \Bigg( \left[ id \boxtimes \tilde{\Delta}\right]^{\otimes m} \otimes \\
&\qquad \left(\sum_{i = 1}^{n}\binom{n}{i} [L_{\mathbf{X}^{<\kappa}_{s,t}} \otimes id]^{\otimes i } \circ (\tilde{\Delta} \boxtimes \Delta)^{i}  \otimes \left[\mathbf{X}_{s,t} \boxtimes \mathbf{1} \right]^{\otimes n-i} - (\tilde{\Delta} \boxtimes \Delta)^{n}  \otimes [L_{\mathbf{X}^{<\kappa}_{s,t}} \otimes id]^{\otimes n } \right)\Bigg).
\end{align*}
It is not hard to see that the each term of this remainder satisfies the regularity constraints in equation \eqref{regularity_condition_r1}, moreover, the expressions above satisfy for a $K > 0$
\[
\big\| R^{g,1, (j,k)}\!\left(v\right)
       \big\|_{\frac{p}{\kappa - j}; \omega_\mathbf{X}} \leq K \|Z\|_{\mathcal{D}^{p,q}_{\mathbf{X}, \mathbf{Y}}}
\]
so that we can conclude for the expansion in the first parameter.\newline
The verification of the second constraint in \eqref{constraints} follows by
identical arguments and is therefore omitted.\newline
\end{proof}


\begin{proof}[Proof of Proposition \ref{stability_proposition}]
To prove the  continuity property we only need to check that $R^{g,1}$ and $R^{g,2}$ are respectively  $\mathbf{Y}$ and $\mathbf{X}$-controlled paths. Since the proof is the analogous for both processes we will only consider the path $R^{g,1}$.\newline
We begin by recalling the decomposition of $R^{g,1}$ derived in the previous proof, in which we are writing
\[
R^{g, 1} = A + B + C + E,
\]
 in order to show that $R^{g,1}$  is controlled by $\mathbf{Y}$ it is sufficient to show that each term in the above sum is controlled by $\mathbf{Y}$.\newline 
Starting with $A$, fix $u$ and expand the factors in the $v$--variable around $u$. This yields
\begin{align*}
A_{s,t}(v)
&=
\sum_{\ell=1}^{\lambda+\kappa-2}\sum_{i=0}^{\ell}\frac{1}{\ell!}\,
D^\ell g\left(Z^{(0,0)}_{t,v}\right)
\circ
\left(
Z^{\otimes(i-1)}_{s,u} \circ [L_{\mathbf{X}^{<\kappa}_{s,t}} \boxtimes id] \otimes R^1_{s,t}(u)\otimes Z^{\otimes(\ell-i)}_{t,u}
\right)
\circ
\left[id \boxtimes L_{\mathbf{Y}^{<\lambda}_{u,v}}\right]^{\otimes\ell}
\circ (\Delta^{\boxtimes}_{>0})^\ell
\\
&\quad + \hat{\mathbf{R}}^{A}\begin{pmatrix} s,t \\ u, v \end{pmatrix},
\end{align*}
where $\hat{\mathbf{R}}^{A}\begin{pmatrix} s,t \\ u, v \end{pmatrix}$ collects the remaining contributions
coming from terms containing one (or more) remainder terms in the expansions of
$Z_{s,v}$, $Z_{t,v}$, and $R^1_{s,t}(v)$.

We now focus on the leading term and reindex the sums as
$i=0,\dots,\lambda+\kappa-2$ and $\ell=i,\dots,\lambda+\kappa-2$:
\begin{equation}\label{pre_Taylor_g_Z_second_rem}
\begin{aligned}
&\sum_{i=0}^{\lambda+\kappa-2}\sum_{\ell=i}^{\lambda+\kappa-2}\frac{1}{\ell!}\,
D^\ell g\left(Z^{(0,0)}_{t,v}\right)
\circ
\left(
 Z^{\otimes(i-1)}_{s,u} \circ [L_{\mathbf{X}^{<\kappa}_{s,t}} \boxtimes id]\otimes R^1_{s,t}(u)\otimes Z^{\otimes(\ell-i)}_{t,u}
\right)
\circ
\left[id \boxtimes L_{\mathbf{Y}^{<\lambda}_{u,v}}\right]^{\otimes\ell}
\circ (\Delta^{\boxtimes}_{>0})^\ell \\
& \qquad +\hat{\mathbf{R}}^{A}\begin{pmatrix} s,t \\ u, v \end{pmatrix}.
\end{aligned}
\end{equation}

Next, we expand $D^\ell g$ around $Z^{(0,0)}_{t,u}$. For each $\ell$ we obtain
\begin{align*}
&D^\ell g\left(Z^{(0,0)}_{t,v}\right)
=
\sum_{m=\ell}^{\lambda+\kappa-2}\frac{1}{(m-\ell)!}\,
D^m g\left(Z^{(0,0)}_{t,u}\right)
\circ
\left(Z^{(0,0)}_{t,v}-Z^{(0,0)}_{t,u}\right)^{\otimes(m-\ell)}
+ \mathrm{Rem}_{\ell}(u,v),
\end{align*}
and we absorb the Taylor remainder $Rem_\ell(u,v)$ into
$\hat{\mathbf{R}}^{A}\begin{pmatrix} s,t \\ u, v \end{pmatrix}$. It immediate to see that this term satisfies the constraint \eqref{remainder_remainder_norm}. Substituting this Taylor in \eqref{pre_Taylor_g_Z_second_rem} yields
\begin{align*}
&=
\sum_{i=0}^{\lambda+\kappa-2}
\sum_{\ell=i}^{\lambda+\kappa-2}
\sum_{m=\ell}^{\lambda+\kappa-2}
\frac{1}{\ell!(m-\ell)!}\,
D^m g\left(Z^{(0,0)}_{t,u}\right)
\\
&\qquad\circ
\left(
Z^{\otimes(i-1)}_{s,u}\otimes R^{1}_{s,t}(u)\otimes Z^{\otimes(\ell-i)}_{t,u}
\otimes
\left(Z^{(0,0)}_{t,v}-Z^{(0,0)}_{t,u}\right)^{\otimes(m-\ell)}
\right)
\circ
\left[id \boxtimes L_{\mathbf{Y}^{<\lambda}_{u,v}}\right]^{\otimes\ell}
\circ (\Delta^{\boxtimes}_{>0})^\ell
\\
&\qquad + \hat{\mathbf{R}}^{A}\begin{pmatrix} s,t \\ u, v \end{pmatrix}.
\end{align*}

Finally, reindex the triple sum by fixing $m$ and summing over $\ell\le m$.
Using $\binom{m}{m-\ell}=\frac{m!}{\ell!(m-\ell)!}$, we obtain
\begin{align*}
&=
\sum_{i=0}^{\lambda+\kappa-2}
\sum_{m=i}^{\lambda+\kappa-2}
\sum_{\ell=i}^{m}
\frac{1}{m!}\binom{m}{m-\ell}\,
D^m g\left(Z^{(0,0)}_{t,u}\right)
\\
&\qquad\circ
\left(
Z^{\otimes(i-1)}_{s,u}\otimes R^{1}_{s,t}(u)\otimes Z^{\otimes(\ell-i)}_{t,u}
\otimes
\left(Z^{(0,0)}_{t,v}-Z^{(0,0)}_{t,u}\right)^{\otimes(m-\ell)}
\right)
\circ
\left[id \boxtimes L_{\mathbf{Y}^{<\lambda}_{u,v}}\right]^{\otimes\ell}
\circ (\Delta^{\boxtimes}_{>0})^\ell
\\
&\qquad + \hat{\mathbf{R}}^{A}\begin{pmatrix} s,t \\ u, v \end{pmatrix}.
\end{align*}
We can now expand the term $Z^{(0,0)}_{t,v} - Z^{(0,0)}_{t,u}$ and include the terms that present at least one remainder in $\hat{\mathbf{R}}^{A}\begin{pmatrix} s,t \\ u, v \end{pmatrix}$.\newline
\begin{align*}
&=
\sum_{i=0}^{\lambda+\kappa-2}
\sum_{m=i}^{\lambda+\kappa-2}
\sum_{\ell=i}^{m}
\frac{1}{m!}\binom{m}{m-\ell}\,
D^m g\left(Z^{(0,0)}_{t,u}\right)
\\
&\qquad \circ \left(
\left(
Z^{\otimes(i-1)}_{s,u}\otimes R^{1}_{s,t}(u)\otimes Z^{\otimes(\ell-i)}_{t,u}
\right)
\circ 
\left[id \boxtimes L_{\mathbf{Y}^{<\lambda}_{u,v}}\right]^{\otimes\ell}
\circ (\Delta^{\boxtimes}_{>0})^\ell \otimes
\left(Z_{t,u} \circ [ \mathbf{1} \boxtimes \mathbf{Y}^{\geq 1}_{s,t}  ]\right)^{\otimes(m-\ell)} \right)
\\
&\qquad + \hat{\mathbf{R}}^{A}\begin{pmatrix} s,t \\ u, v \end{pmatrix}.
\end{align*}
    
At this stage, we extract the principal contribution by replacing the inner
$\ell$--sum with the corresponding $m$--fold term, and we include the resulting
difference in the remainder:
\begin{align*}
&=
\sum_{i=0}^{\lambda+\kappa-2}
\sum_{m=i}^{\lambda+\kappa-2}
\frac{1}{m!}\,
D^m g\left(Z^{(0,0)}_{t,u}\right)
\circ
\left(
Z^{\otimes(i-1)}_{s,u}\otimes R^1_{s,t}(u)\otimes Z^{\otimes(m-i)}_{t,u}
\right)
\circ
(\Delta^{\boxtimes}_{>0})^{m}
\circ
\left[id \boxtimes L_{\mathbf{Y}^{<\lambda}_{u,v}}\right]
\\
&\qquad + \mathbf{R}^{A}\begin{pmatrix} s,t \\ u, v \end{pmatrix},
\end{align*}
where $\mathbf{R}^{A}\begin{pmatrix} s,t \\ u, v \end{pmatrix}$ absorbs both
$\hat{\mathbf{R}}^{A}\begin{pmatrix} s,t \\ u, v \end{pmatrix}$ and the truncation error arising from the
replacement of the inner $\ell$--sum by
$(\Delta^{\boxtimes}_{>0})^{m}\circ[id \boxtimes L_{\mathbf{Y}^{<\lambda}_{u,v}}]$. The latter term is the second level analogous of the term $C$ in the expansion of the first-level remainder, and can be seen to satisfy inequality \eqref{remainder_remainder_norm}.

We now turn to the terms $B_{s,t}(v)$ and $C_{s,t}(v)$, defined respectively as 
\begin{align*}
B_{s,t}(v)
&:=
\sum_{i = 1}^{\kappa -1} \sum_{m = 0}^{\lambda -1} \frac{1}{(\kappa-1)!} \frac{1}{m!} \binom{\kappa-1}{i} \int_0^1 (1-r)^\kappa D^{\kappa+m}g((1-r)Z^{(0,0)}_{s, v} + rZ^{(0,0)}_{t, v}) dr\\
& \qquad \circ \left((Z_{s,v} \circ [L_{\mathbf{X}^{<\kappa}_{s,t}} \boxtimes id])^{\otimes i + m} \circ \left((\tilde{\Delta} \boxtimes \Delta)^{i} \otimes ( id \boxtimes \tilde{\Delta})^{m}\right) \otimes (Z^{0,0}_{t,v} - Z^{0,0}_{s,v})^{\otimes \kappa-i}\right), \\
C_{s,t}(v) &:= \sum_{i = 1}^{\kappa -1} \sum_{n = i}^{\kappa -1} \sum_{m = 0}^{\lambda -1} \sum_{p=0}^{n-i} \frac{1}{n!} \frac{1}{m!} \binom{n}{i} D^{n+m}g(Z^{(0,0)}_{s,v})\\
& \qquad \circ (Z_{s,v} \circ [L_{\mathbf{X}^{<\kappa}_{s,t}} \otimes id])^{\otimes i + m} \\
&\qquad \circ \left((\tilde{\Delta} \boxtimes \Delta))^{i} \otimes ( id \boxtimes \tilde{\Delta})^{m} \otimes \left(Z_{s,v} \circ \left[\mathbf{X}^{\geq 1 }_{s,t} \boxtimes id \right]\right)^{\otimes p-1} \otimes R^{1, (0,0)}_{s,t}(v) \otimes \left(Z^{(0,0)}_{t,v} - Z^{(0,0)}_{s,v}\right)^{\otimes n-i-p} \right). 
\end{align*}
A procedure analogous to the one used for the term $A$ shows that
\begin{align*}
    &B_{s,t}(v) = B_{s,t}(u) \circ [id \boxtimes L_{\mathbf{X}^{< \kappa}}] + \mathbf{R}^{B}\begin{pmatrix} s,t \\ u, v \end{pmatrix}, \\
    &C_{s,t}(v) = C_{s,t}(u) \circ [id \boxtimes L_{\mathbf{X}^{< \kappa}}] + \mathbf{R}^{C}\begin{pmatrix} s,t \\ u, v \end{pmatrix}. \\
\end{align*}
Where $\mathbf{R}^{B}\begin{pmatrix} s,t \\ u, v \end{pmatrix}$ and $\mathbf{R}^{C}\begin{pmatrix} s,t \\ u, v \end{pmatrix}$
satisfy the bound \eqref{remainder_remainder_norm}.

We finally treat the term $E$. Set
\begin{align*}
E_{s,t}(v)
&:=
\sum_{i = 1}^{\kappa -1} \sum_{n = i}^{\kappa -1} \sum_{m = 0}^{\lambda -1}  \frac{1}{n!} \frac{1}{m!} \binom{n}{i} D^{n+m}g(Z^{(0,0)}_{s,v})\circ Z^{\otimes m + n}_{s,v} \circ \Bigg(  \left[ id \boxtimes \tilde{\Delta}\right]^{\otimes m}   \\
&\qquad \otimes \left(\sum_{i = 1}^{n}\binom{n}{i} [L_{\mathbf{X}^{<\kappa}_{s,t}} \otimes id]^{\otimes i } \circ (\tilde{\Delta} \boxtimes \Delta)^{i}  \otimes \left[\mathbf{X}^{\ge 1}_{s,t} \boxtimes \mathbf{1} \right]^{\otimes n-i} - (\tilde{\Delta} \boxtimes \Delta)^{n}  \otimes [L_{\mathbf{X}^{<\kappa}_{s,t}} \otimes id] \right)\Bigg)\\ 
&= \sum_{i = 1}^{\kappa -1} \sum_{n = i}^{\kappa -1} \sum_{m = 0}^{\lambda -1}  \frac{1}{n!} \frac{1}{m!} \binom{n}{i} D^{n+m}g(Z^{(0,0)}_{s,v})\circ Z^{\otimes m + n}_{s,v} \circ \Bigg(  \left[ id \boxtimes \tilde{\Delta}\right]^{\otimes m} \otimes  G(\mathbf{X}^{\geq 1}_{s,t})\Bigg)
\end{align*}
proceeding in a totally analogous way to what we did above, we obtain expanding $Z$ in the second parameter and grouping the terms that present at least one remainder term $R^2$ in $\hat{\mathbf{R}}^{E}$,
\begin{align*}
&=
\sum_{i=1}^{\kappa-1} \sum_{n=i}^{\kappa-1} \sum_{m=0}^{\lambda-1} 
\frac{1}{n!} \frac{1}{m!} \binom{n}{i}
D^{n+m} g(Z^{(0,0)}_{s,v}) \circ Z^{\otimes m+n}_{s,u} \\
&\quad \circ
\Bigg(
    \left[ id \boxtimes L_{\mathbf{Y}^{<\lambda}_{u,v}} \right]
    \circ \left[ id \boxtimes \tilde{\Delta} \right]^{\otimes m}  \otimes
    \left[ id \boxtimes L_{\mathbf{Y}^{<\lambda}_{u,v}} \right]^{\otimes n}
    \circ G(\mathbf{X}^{\ge 1}_{s,t})
\Bigg) \\
&\quad
+ \hat{\mathbf{R}}^{E} \begin{pmatrix} s,t \\ u,v \end{pmatrix} .
\end{align*}

From definition of $R^2$ it can be seen that the term $\hat{\mathbf{R}}^{E}\begin{pmatrix} s,t \\ u, v \end{pmatrix}$ satisfies \eqref{remainder_remainder_norm}.\newline
Subsequently, by expanding $g$ around $Z^{(0,0)}_{s,u}$ and absorbing the Taylor remainder in $\hat{\mathbf{R}}^{E}\begin{pmatrix} s,t \\ u, v \end{pmatrix}$
\begin{align*}
        = \sum_{i = 1}^{\kappa -1} \sum_{n = i}^{\kappa -1} \sum_{m = 0}^{\lambda -1} \sum_{l=m}^{\lambda -1} &\frac{1}{n!} \frac{1}{l!} \binom{n}{i} \binom{l}{m} D^{n+l}g(Z^{(0,0)}_{s,u})\circ Z^{\otimes l + n}_{s,u}\\ 
        & \circ \Bigg( \left( \left[id \boxtimes L_{\mathbf{Y}^{< \lambda}_{u,v}}\right] \circ \left[ id \boxtimes \tilde{\Delta}\right]\right)^{\otimes m} 
         \otimes  \left(\left[id \boxtimes L_{\mathbf{Y}^{< \lambda }_{u,v}}\right]\right)^{\otimes n} \circ G(\mathbf{X}^{\geq 1}_{s,t})\otimes \left[ \mathbf{1} \boxtimes \mathbf{Y}^{\geq 1}_{u,v} \right]^{\otimes l-m}\Bigg) \\
         \qquad + \hat{\mathbf{R}}^{E}\begin{pmatrix} s,t \\ u, v \end{pmatrix}.
\end{align*}
Again, the newly absorbed term satisfies \eqref{remainder_remainder_norm}.\newline
Finally proceeding as we did in the previous proof in order to recover $C$ we obtain the desired identity i.e.
\begin{align*}
E_{s,t}(v) = E_{s,t}(u)\circ [id \boxtimes L_{\mathbf{Y}^{< \lambda}_{u,v}}] + \mathbf{R}^{E}\begin{pmatrix} s,t \\ u, v \end{pmatrix}, 
\end{align*}
where $\mathbf{R}^{E}\begin{pmatrix} s,t \\ u, v \end{pmatrix}$  is the sum of this newly found second level remainder term and $\hat{\mathbf{R}}^{E}\begin{pmatrix} s,t \\ u, v \end{pmatrix}$. This allows to conclude by noting that $\mathbf{R}^{E}\begin{pmatrix} s,t \\ u, v \end{pmatrix}$ respects the bound \eqref{remainder_remainder_norm} and that, from construction, there exists a $K > 0$ such that 
\[
\big\| \mathbf{R}^{F, (j,k)} \big\|_{(\frac{p}{\kappa - j},\frac{q}{\lambda - k}), (\omega_\mathbf{X}, \omega_\mathbf{Y})} \leq K \| Z \|_{\mathcal{D}_{\mathbf{X}, \mathbf{Y}}},
\] 
for all $j \in [0:\kappa)$, $k \in [0:\lambda)$, $F \in \{A, B, C, E\}$.
\end{proof}

\section{Existence of the rough Schwinger-Dyson equation}\label{appendix_existence_SD}


In this section, we present the assumptions and the proof of existence for the rough Schwinger--Dyson kernel and its related equation. This result is a variant of Theorem 3.10 in \cite{cass2022combinatorial}, extended to the case where the driving signal is a geometric rough path of arbitrarily low regularity. 

Throughout the remainder of this section, we use $\mathfrak{h}_N$ to denote the space of $N \times N$ Hermitian matrices, for which the identity $\mathfrak{u}_N = i\,\mathfrak{h}_N$ holds. Moreover, we denote by $U(N, \mathbb{C})$ the unitary group and $\mathcal{P}(W)$ the space of Borel probability measures on a finite-dimensional Banach space $W$. Recall that every path in $X \in C^1_{ \omega}([a,b], V)$ can be lifted to the  geometric $p$-rough path
\[
(s,t) \to \left(1, X_{s,t}, \int_{s < u_1 < u_2 < t} dX_{u_1} \otimes dX_{u_2}, \dots, \int_{s < u_1 < \dots < u_\kappa < t} dX_{u_1} \otimes \dots \otimes dX_{u_\kappa} \right),
\]
where each integral is defined in the Riemann-Stieltjes sense, we refer to this as the natural lift of $X$.\newline
Within this framework, we are now in a position to state and prove the following existence theorem.
\begin{theorem}\label{theorem_existence}
Let $\mathbf X \in G\Omega_p(V)$. For each $N\in\mathbb N$ and each $1\le i\le d$, let
\(
A_i^N:\Omega\to\mathfrak h_N
\)
be a Hermitian random matrix defined on a probability space $(\Omega,\mathcal F,\mathbb P)$. Assume that the following hold.

First, for every $k\in\mathbb N$,
\[
\sup_{N\in\mathbb N}\sup_{1\le i\le d}\sup_{1\le m,l\le N}
\mathbb E\big[|A_i^N(m,l)|^k\big]\le c_A(k)<\infty,
\]
and moreover
\[
\mathbb E[A_i^N(m,l)]=0,
\qquad
\mathbb E[|A_i^N(m,l)|^2]=1.
\]

Second, there exists $K>0$ such that, for every $r\in\mathbb N$,
\[
\sup_{N\in\mathbb N}\sup_{0\le s<t\le T}
\mathbb E\!\left[\frac{1}{N^{r/2}}\|A_i^N(m,l)\|_{\mathrm{op}}^r\right]
\le K^r\Gamma\!\left(\frac{r}{2p}+1\right).
\]

Finally, assume that the sigma-algebras
\[
\bigvee_{i=1}^d \sigma(A_i^N(m,l)),
\qquad 1\le m\le l\le N,
\]
are independent, and that
\[
\mathbb E\!\left[A_i^N(m,l)A_j^N(m,l)\right]=0,
\qquad i\ne j.
\]

Define
\[
M^N\in\operatorname{Hom}(\mathbb R^d,\mathfrak h_N),
\qquad
M^N(\eta)=\sum_{i=1}^d A_i^N\,\eta^{(i)}.
\]
Let $Z^N(\mathbf X)\in\mathcal D^p([0,T],U(N;\mathbb C))$ denote the solution of
\begin{equation} \label{eq:36}
dZ^N(s,t)=\frac{i}{\sqrt N}M^N(d\mathbf X_t)\,Z^N(s,t),
\qquad
Z^N(s,s)=I_N.
\end{equation}

Then, for every $0\le s\le t\le T$ and for $\xi_N \in \mathcal{P}(\mathfrak{h}^d_N)$ denoting the law of the hermitian matrices $A^N_i$, the limit
\[
K(s,t)=\lim_{N\to\infty}\frac1N\,
\mathbb E_{\xi_N}\!\left[\operatorname{Tr}(Z^N(s,t))\right]
\]
exists and is given by
\begin{equation}\label{rough_SD_rep}
K(s,t)=\sum_{|L|=0}^\infty i^{|L|}\varphi(L)\,S^{(L)}(\mathbf X)_{s,t}.
\end{equation}
Moreover, \(K\) solves the two-parameter RDE
\begin{equation}\label{SD_with_f}
K(s,t)=1-\int_s^t\int_s^r f(K(s,u),K(u,r))\,d\mathbf X_u\,d\mathbf X_r,
\end{equation}
where
\[
f:\mathbb R^2\to \mathrm{Hom}(V\boxtimes V, \mathbb{R}),
\qquad
f(z_1,z_2)[\eta,\xi]=z_1z_2\langle \eta,\xi\rangle_V.
\]
\end{theorem}

\begin{proof}
Let \((X^n)_{n\ge1}\) be a sequence of paths with bounded variation whose canonical lifts
\(\mathbf X^n\) converge to \(\mathbf X\) in the \(p\)-variation topology. Then, without loss of generality we can assume that, for every non-empty word \(I\) and all \(n\),
\begin{equation}\label{eq:unif-signature-bound}
\left|S^{(I)}(X^n)_{s,t}\right|
\;\le\;
\frac{\omega_{\mathbf{X}}(s,t)^{\,|I|/p}}{\beta\,\Gamma\!\left(|I|/p+1\right)}.
\end{equation}
By Lyons’ extension theorem, the same estimate holds for \(S(\mathbf X)\).

Let \(Z^{N,n}\) be the solution to \eqref{eq:36} with driving signal \(\mathbf X^n\). Expanding as in Theorem~2.1 in \cite{cass2024free} yields
\[
\mathbb E\!\left[\frac{1}{N}\operatorname{Tr}(Z^{N,n}_{st})\right]
=
\sum_{|I|=1}^{\infty}
\frac{i^{\,|I|}}{N^{|I|/2+1}}\;
\mathbb E\!\left[\operatorname{Tr}(A_I^N)\right]\,
S^{(I)}(X^n)_{st}.
\]
Taking absolute values and using \(|\operatorname{Tr}B|\le d\,\|B\|_{\mathrm{op}}\), together with
\eqref{eq:unif-signature-bound}, we obtain the uniform (in \(N,n\)) bound
\begin{equation}\label{eq:series-bound}
\mathbb E\!\left[\left|\tfrac{1}{N}\operatorname{Tr}(Z^{N,n}_{s,t})\right|\right]
\;\le\;
\sum_{|I|=1}^{\infty}
\frac{1}{N^{|I|/2+1}}\;
\mathbb E\!\left[\left|\operatorname{Tr}(A_I^N)\right|\right]\,
\frac{\omega(s,t)^{\,|I|/p}}{\beta\,\Gamma\!\left(|I|/p+1\right)}
\;<\;\infty.
\end{equation}
(Under the standing assumptions on \(A_i^N\), moment bounds and
independence across distinct indices, the series on the right is finite,
provides an \(N,n\)-independent \(L^1\) dominator.)

Consequently, by Tonelli’s theorem we may interchange expectation and the (absolutely
convergent) series, and by dominated convergence (using \eqref{eq:series-bound} as the
dominator) we may pass to the limits \(n\to\infty\) and \(N\to\infty\). Therefore,
\[
\begin{aligned}
\lim_{N\to\infty}\mathbb E\left[Z^N_{s,t}\right]
&= \lim_{N\to\infty}\mathbb E\!\left[\lim_{n\to\infty} Z^{N,n}_{s,t}\right]
 = \lim_{N\to\infty}\lim_{n\to\infty}\mathbb E\left[Z^{N,n}_{s,t}\right] \\
&= \lim_{n\to\infty}\sum_{|I|=0}^{\infty} i^{\,|I|}\,\varphi(I)\,S^{(I)}(X^n)_{s,t}
 = \sum_{|I|=0}^{\infty} i^{\,|I|}\,\varphi(I)\,S^{(I)}(\mathbf X)_{s,t},
\end{aligned}
\]
which is the desired identity \eqref{rough_SD_rep}.\newline
From this expansion is clear that $K \in \mathcal{D}^p_{\mathbf{X}}(\Delta_{[a,b]}, \mathbb{R})$. An entirely analogous argument to the one above allows us to characterize all higher-level components in the truncated tensor algebra $T^{<\kappa}(V) \boxtimes T^{<\kappa}(V)$ through the formula
\[
K^{(I, J)}(s,t)
= (-1)^{|I|}\sum_{|L|=0}^{\infty} i^{\,|JLI|}\,\varphi(JLI)\, S^{(L)}(\mathbf{X})_{s,t}.
\]
As a consequence, the sequence of Schwinger-Dyson kernel approximations $\{K^n\}$, once lifted to a $p$-controlled rough path, satisfies
\[
d_{\Delta_T}^{\mathbf{X}^n, \mathbf{X}}(K^n, K) \to 0 \qquad \text{as } n \to \infty.
\]
Recall that, as derived in Section 3 of  \cite{cass2024free}, for each $\mathbf{X}^n$, the Schwinger-Dyson kernel satisfies the equation
\[
K^n(s,t) = 1 -  \int_s^t \int_s^r K^n(s,u) K^n(u,r) \langle dX^n_u, \; dX^n_r \rangle.
\]
The fact that the kernel $K$ solves equation~\eqref{SD_with_f} then follows from this convergence together with the stability properties of two-parameter rough integrals on two-simplices, derived from the stability of joint rough integrals on rectangles (Theorem \ref{stability_theorem_2D}). This completes the proof.
\end{proof}



\section{Algorithm complexity} \label{appendix:algorithmic_complexity}

\begin{proposition}
    The set $NC(n)$ constructed over the alphabet $1, \dots, d$, contains 
    \[
    \sum_{i = 0}^{\floor{n/2}} \frac{n -2i +1}{n+1} \binom{n+1}{i} d^{n - i}
    \] 
    elements.
\end{proposition}
\begin{proof}
Denote by $A_{n,m}$ the number of words appearing in $NC(n)$ with length equal to $m$. We know that
\[
A_{0, 0} = 1, \qquad A_{0, m} = 0 \, \text{ for any } \, m \geq 1. 
\]
Moreover, it is clear from the definition of $NC(n)$ that 
\begin{align}\label{empty_SD_words_recursion}
A_{n, 0} = A_{n-1, 1}.
\end{align}
For $m \geq 1$ we also have the following relationship
\begin{align} \label{SD_words_recursion}
    A_{n, m} = d A_{n-1, m-1} + A_{n-1, m+1},
\end{align}
in fact, we can obtain a word in $NC(n)$ with length $m$ by either appending a letter at the end of a word with length $m-1$ in $NC(n-1)$
or by removing the last letter to each word with length $m+1$ in the set $NC(n-1)$.\newline
Additionally notice that for $m\geq1$, either $m$ has the same parity of $n$ or $A_{n,m} = 0$.\newline
The next step is to show that for $m \equiv n \, mod(2)$
\[
A_{n, m} = d^{\frac{n+m}{2}} \left( \binom{n}{\frac{n+m}{2}} - \binom{n}{\frac{n+m}{2} + 1}  \right).
\] 
We proceed by induction over $n$. For the base case we have 
\[
A_{0,0} = 1, \qquad A_{0,m} = 0 \text{ for }\, m \geq 1. 
\]
For the induction step, suppose the identity holds for $n-1$. Fix $m \geq 1$ and define 
$u :=  \frac{m +n}{2}$, now using the identity in equation \eqref{SD_words_recursion} and the induction hypothesis allows to obtain 
\begin{align*}
    A_{n,m} &= d A_{n-1, m-1} + A_{n-1, m+1} \\
            &= d d^{u-1} \left( \binom{n-1}{u-1} - \binom{n-1}{u} \right) + d^u\left( \binom{n-1}{u} - \binom{n-1}{u+1} \right) \\
            &= d^u\left( \binom{n-1}{u-1} - \binom{n-1}{u+1} \right) \\
            &= d^u\left( \binom{n}{u} - \binom{n}{u+1} \right)
\end{align*}
where the last step uses Pascal' identity.
This concludes for the case $m \geq 1$. 
For the case $m = 0$, we have by equation \eqref{empty_SD_words_recursion} and Pascal's identity that
\begin{align*}
    A_{n,0} &=  A_{n-1, m+1} \\
            &=  d^{n/2}\left( \binom{n-1}{n/2} - \binom{n-1}{n/2+1} \right). 
\end{align*}
Now it is clear that the cardinality of $NC(n)$ satisfies the identities 
\begin{align*}
    \sum_{m \geq 0} A_{n,m} = \sum_{m \geq 0} d^{\frac{n+m}{2}} \left( \binom{n}{\frac{n+m}{2}} - \binom{n}{\frac{n+m}{2} + 1}  \right),
\end{align*}
defining $i := \frac{n - m}{2}$ we get
\begin{align*}
    &= \sum_{i=0}^{\floor{n/2}} d^{n-i} \left( \binom{n}{n-i} - \binom{n}{n - i + 1}  \right)\\
    &= \sum_{i = 0}^{\floor{n/2}} \frac{n -2i +1}{n+1} \binom{n+1}{i} d^{n - i},
\end{align*}
concluding the proof.
\end{proof}

This rule can be generalized to cover the case where pairings between the first $n$ letters are ruled out. 

\begin{proposition}\label{identity_cardinality_NC}
For $\ell \geq n$, the cardinality of the set $NC(\ell)_n$ over the alphabet $\{1,\dots,d\}$ is
\[
|NC(\ell)_n|
= d^n \sum_{i=0}^{\lfloor (t+n)/2 \rfloor}
\left( \binom{t}{i} - \binom{t}{i-n-1} \right) d^{\,t-i},
\]
with $t := \ell - n$.
\end{proposition}

\begin{proof}
The proof of the identity follows in a very similar way as the formula for the cardinality of $NC(n)$ so we will just sketch the proof for this case.\newline
Denote by $A_{\ell,m}$ the number of words appearing in $NC(\ell)_n$ with length equal to $m$. It is easy to see that
\[
A_{n, m} = \begin{cases}
d^n & m = n,\\
0 & \text{ otherwise,}
\end{cases}
\]
For $\ell > n$ and $m =0$, we have the recursion 
\begin{align*}
A_{\ell, 0} = A_{\ell-1, 1}.
\end{align*}
For $\ell > n$ and $m \geq 1$, we have the following relationship
\begin{align*}
    A_{\ell, m} = d A_{\ell-1, m-1} + A_{\ell-1, m+1}.
\end{align*}
Just like the previously analysed case, a simple induction argument allows to conclude the proof.
\end{proof}

\begin{corollary}
For every $t \in (0 : n]$, we have the following upper and lower bounds for $NC(\ell)_n$:
\[
 \frac{1}{2} d^n (d+1)^t
\leq | NC(\ell)_n |
\leq d^n (1+d)^t.
\]
Moreover,  the right inequality holds for every $t > 0$.
\end{corollary}

\begin{proof}
The proof of the right inequality immediately from the definition of $NC(\ell)_n$
\begin{align*}
NC(\ell)_n
&\leq d^n \sum_{i=0}^{\lfloor (t+n)/2 \rfloor}
\binom{t}{i} d^{\,t-i} \\
&= d^n \sum_{i=0}^{\min\left(t, \lfloor (t+n)/2 \rfloor\right)}
\binom{t}{i} d^{\,t-i} \\
&\leq d^n \sum_{i=0}^{t}
\binom{t}{i} d^{\,t-i} \\
&\leq d^n (1+d)^t.
\end{align*}
For the left inequality, define the binomially distributed random variable $V \sim B(t, \frac{1}{d+1})$. Then
\begin{align*}
    |NC(\ell)_n| &= (d+1)^t\left( d^n \mathbb{P}(V \leq \floor{(n+t)/2}) - d^{-1} \mathbb{P}(V \leq \floor{(n+t)/2} - n - 1) \right),
\end{align*}
so that if $t \leq n$, recalling that $\frac{1}{d+1} \leq \frac{1}{2}$ then
\begin{align*}
    NC(\ell)_n &\geq (d+1)^t\left( d^n \mathbb{P}(V \leq \floor{(n+t)/2})\right)\\
    &\geq \frac{1}{2} (d+1)^t  d^n .
\end{align*}
This concludes the proof.

\end{proof}

\begin{corollary}
For every $1 \le i \le m \le j \le 2^N - 1$, let $\mathrm{Poly}$ denote the number of operations required to compute
\[
\sum_{|\omega| = 0}^{\kappa-1} \sum_{(\alpha, \beta) \in Sh^{-1}(\omega)} (-1)^{|\beta|}  
K^{(\alpha)}_{\tilde{I}}(i,m)\, K^{(\beta)}_{\hat{I}}(m,j+1).
\]
for every word $\gamma$ with length $|I| \leq \kappa + \zeta$.\newline
Then 
\[
c
(2d)^{\kappa + \zeta-1} (d+1)^\kappa \leq  \mathrm{Poly} \leq C
(2d)^{\kappa + \zeta-1} (d+1)^\kappa,
\]
where the constants $C, c \geq 0$ do not depend on the other parameters appearing in the estimate above.
\end{corollary}

\begin{proof} 
The uniform upper bound for $|NC(\ell)_n|$ yields
\begin{align*}
&\sum_{k=1}^{\kappa + \zeta}
\sum_{n=0}^{\kappa-1}
\sum_{j=0}^{n} 
\binom{k-1}{j}\, 
\left| NC\left(\ceil{k/2} + j - 1\right)_{\ceil {k/2}-1}\right|
\, 
\left|NC\left(k- \ceil{k/2}  + n - j\right)_{k -\ceil{k/2}}\right| \\
&\leq \sum_{k=1}^{\kappa + \zeta}
\sum_{n=0}^{\kappa-1}
\sum_{j=0}^{n} 
\binom{k-1}{j}\, d^{k-1}(d+1)^{n} \\
&\leq \sum_{k=1}^{\kappa + \zeta}
\sum_{n=0}^{\kappa-1}
2^{k-1} d^{k-1}(d+1)^{n} \\
&\leq
C (2d)^{\kappa + \zeta-1} (d+1)^\kappa.
\end{align*}
For the lower bound, we get
\begin{align*}
&\sum_{k=1}^{\kappa + \zeta}
\sum_{n=0}^{\kappa-1}
\sum_{j=0}^{n} 
\binom{k-1}{j}\, 
\left| NC\left(\ceil{k/2} + j - 1\right)_{\ceil{k/2}-1}\right|
\, 
\left|NC\left(k- \ceil{k/2}  + n - j\right)_{k -\ceil{k/2}}\right| \\
&\geq \sum_{j=0}^{\kappa-1}
\binom{\kappa+\zeta-1}{j}\,
\left|NC\left(\ceil{\frac{\kappa+\zeta}{2}} + j - 1\right)_{\ceil{\frac{\kappa+\zeta}{2}}-1} \right|
\, \left|
NC\left(\kappa + \zeta -\ceil{\frac{\kappa+\zeta}{2}}  + \kappa - j\right)_{\kappa + \zeta-\ceil{\frac{\kappa+\zeta}{2}}} \right| \\
&\geq (d+ 1)^{\kappa} \left( \frac{1}{2} d^{\ceil{\frac{\kappa + \zeta}{2}}-1} \right) \left( \frac{1}{2} d^{{\kappa + \zeta - \ceil{\frac{\kappa + \zeta}{2}}}} \right) \sum_{j=0}^{\kappa-3} \binom{\kappa+\zeta-1}{j} \\
&\geq c (d+ 1)^{\kappa} d^{\kappa + \zeta -1} \sum_{j=0}^{\kappa-3} \binom{\kappa+\zeta-1}{j} \\
&\geq c 2^{\kappa + \zeta -  3} d^{\kappa+\zeta -1} (d+ 1)^{\kappa},
\end{align*}
where the  constant $c$ above is independent of the parameters of interest.\newline
This concludes the proof.
\end{proof}

\begin{corollary}\label{corollary_trace_lvl}

For every $1 \le i \le m \le 2^N - 1$, let $\mathrm{Poly}_{\emptyset}$ denote the number of operations required to compute
\[
\sum_{|\omega| = 1}^{\kappa} \sum_{m=i+1}^{j+1} K^{(\omega)}(i,m) \, \mathbf{X}^{(\omega)}_{m-1, m}.
\]

Then, for $d > 1$, there exists $\alpha^* \in [0,1/2]$ that maximizes the expression on the left, such that
\[
c\left(d^{1 - \alpha^*} \frac{2^{H(\alpha^*)}}{\sqrt{\kappa + \zeta}}\right)^{\kappa + \zeta}
\;\leq\; \mathrm{Poly}_{\emptyset} \;\leq\; C(d+1)^{\kappa+\zeta},
\]
where $H(\alpha) = -\alpha \log_2 \alpha - (1-\alpha) \log_2 (1-\alpha)$ is the binary entropy function and some constants $C,c \geq 0$ independent from $\kappa, \zeta$ and $d$.

In the special case $d=1$, we have
\[
\frac{2^{\kappa + \zeta}}{\kappa + \zeta + 1}  \leq \mathrm{Poly}_{\emptyset} \leq C\left(2^{\kappa + \zeta}\right).
\]

  
\end{corollary}
\begin{proof}
One can see that, 
\begin{align*}
&\text{Poly}_{\emptyset} = \sum_{k = 1}^{2\kappa - 2} |NC(k)| \\
& \quad = \sum_{k = 1}^{\kappa +\zeta}  \sum_{i = 0}^{\floor{k/2}} \frac{k -2i +1}{k+1} \binom{k+1}{i} d^{k - i}\\
& \quad \leq \sum_{k = 1}^{\kappa + \zeta} d^{k} ( 1 + 1/d)^{k+1} \leq C (d+1)^{\kappa+\zeta}.
\end{align*}
Moreover, for $d > 1$, recalling that
\[
\sum_{i=0}^{\floor{k/2}} \binom{k+1}{i} d^{-i} \geq \sup_{0 \leq \alpha \leq 1/2 }   \binom{k+1}{\floor{k\alpha}} d^{-\floor{k\alpha}} \geq C_\alpha \frac{2^{H(\alpha) k }}{\sqrt{k}} d^{-\floor{k\alpha}},
\]
where $H(\alpha) = -\alpha \log_2\alpha -(1-\alpha)\log_2(1-\alpha)$,
we also have
\begin{align*}
\text{Poly}_{\emptyset} &= \sum_{k = 1}^{\kappa + \zeta}  \sum_{i = 0}^{\floor{k/2}} \frac{k -2i +1}{k+1} \binom{k+1}{i} d^{k - i} \\
&\geq \sup_{0 \leq \alpha \leq 1/2}\sum_{k = 1}^{\kappa + \zeta} c_\alpha d^{k - \floor{k\alpha}} \binom{k}{\floor{k\alpha}} \\
&\geq \sup_{0 \leq \alpha \leq 1/2}\sum_{k = 1}^{\kappa + \zeta} C_\alpha d^{k - \floor{k\alpha}} \frac{2^{H(\alpha)k}}{\sqrt{k}}  \\
& \geq c \left(d^{(1 - \alpha^*)} \frac{2^{H(\alpha^*)k}}{\sqrt{\kappa + \zeta}}\right)^{\kappa + \zeta},
\end{align*}
where $\alpha^*$ is one the maximizers of the penultimate term.\newline
For $d = 1 $, using the fact that \[
\frac{k -2i +1}{k+1}  \geq \frac{1}{k+1}, 
\] we obtain
\begin{align*}
&\text{Poly}_{\emptyset} = \sum_{k = 1}^{\kappa +  \zeta} |NC(k)| \\
& \quad = \sum_{k = 1}^{\kappa + \zeta} \sum_{i = 0}^{\floor{k/2}} \frac{k -2i +1}{k+1} \binom{k+1}{i}\\ 
& \geq \sum_{k = 1}^{\kappa + \zeta}  \frac{1}{k+1} 2^{k} \\
& \geq \frac{2^{\kappa + \zeta}}{\kappa + \zeta -1}.
\end{align*}
This concludes the proof.
\end{proof}

\begin{comment}

For the polynomial appearing in .... we have 

Assume that the signature of $\mathbf{X}$ is provided,
We are now ready to prove the algorithmic complexity of the scheme proposed in ... 
Let $D = \sum_{k=0}^\kappa d^k $ be the dimension of the truncated tensor algebra $T^{ 2\kappa - 2}(\mathbb{R}^d)$.\newline
we recall that the algorithm consists in sequentially filling the entries of the matrix $K^{sol}_{i,j} \in \mathbb{R}^{2^N \times 2^N}$, where each entry of the matrix belongs to $T^{ 2\kappa - 2}(\mathbb{R}^d)$.\newline
the initial condition involves the principal diagonal which is filled with $id_{T^{ 2\kappa - 2}(\mathbb{R}^d)}$ and the succesive iterations fill the lower diagonals.\newline
The update rule for the level zero is given by:
\begin{align*}
    K(i,j+1)
    &= 1+\sum_{|\omega| = 1}^{\kappa} \sum_{m=i+1}^{j+1}  K_{\omega}(i,m)  \mathbf{x}^{(\omega)}_{m-1, m}
\end{align*}
which requires a computational cost $O(\sum_{n = 0}^\kappa |\mathcal{P}_{0, n}|)$.\newline
For the other levels we are computing 
\begin{align*}
    K^{\gamma}(i,j+1) &=  \sum_{m=i+1}^{j}\sum_{(\alpha, \beta) \in Sh^{-1}(\omega)} \sum_{|\omega| = 1}^{\kappa} (-1)^{|\beta|}  K_\alpha^{\tilde{\gamma}}(i,m) K_\beta^{\hat{\gamma}}(m,j+1)  \mathbf{x}^{(\gamma_\ell, \omega)}_{m-1, m} \\
    & \qquad + \sum_{(\alpha, \beta) \in Sh^{-1}(\omega)} \sum_{|\omega| = 1}^{\kappa} (-1)^{|\beta|}  K_\alpha^{\tilde{\gamma}}(i,j+1) K_\beta^{\hat{\gamma}}(j+1,j+1)  \mathbf{x}^{(\gamma_\ell, \omega)}_{j, j+1}.
\end{align*}
the first term in the right side of the equality can be obtained via $O(\sum_{k=0}^ {2\kappa - 2/2} \sum_{n = k}^{\kappa + k} |\mathcal{P}_{k, n}|)$ operations and the term on the 

\end{comment}



\subsection{Computational complexity}
\label{subsec:complexity}

We analyze the cost of the solver in two stages: (i) compile-time preprocessing
(which depends only on the algebraic specifications and is reused across paths),
and (ii) runtime dynamic programming (DP) for a given sequence of increments.

\paragraph{Notation.}
Let
\begin{itemize}
  \item $2^N \in \mathbb{N}$ be the time grid size,
  \item $D$ be the total dimension of the tensor algebra $T^{\leq \kappa + \zeta}(\mathbb{R}^d)$,
  \item $\mathrm{Poly}$ be the number of monomials in the compiled polynomial representation,
  \item $\mathrm{Poly}_{\mathrm{lin}}$ be the number of right-linear monomials, i.e.\ terms with $\tilde{\gamma} = \alpha =\emptyset$,
  \item $\mathrm{Poly}_{\emptyset}$ be the number of terms in the empty-word rule used for row replacement.
\end{itemize}

% ------------------------------------------------------------
\subsubsection{Compile-time preprocessing}
% ------------------------------------------------------------

The compile-time container performs two main precomputations.

\paragraph{(1) Empty-word rule.}
For the trace level computation, expressed in the inner sum of equation \eqref{trace_numerical}, loop over all non-empty basis words and expands each word via
non-crossing pairings. From Corollary \ref{corollary_trace_lvl} we know the number of operation for this step is equal to
\begin{equation}
  \label{eq:empty-rule-cost}
  \mathrm{Cost}_{\emptyset}
  \; = \;\mathcal{O}\left(
    \sum_{w \neq \emptyset} |w| \, NC(|w|)\right)
  \;=\; \mathcal{O} \left(
  \mathrm{Poly}_{\emptyset}\right).
\end{equation}
This cost is paid once per compiled solver and is amortized across all runtime executions.

\paragraph{(2) Right-linear part extraction.}
Extracting the right-linear part of the polynomial, which is used to compute $A_j$, can be implemented by a boolean mask applied to arrays of length $\mathrm{Poly}$, hence
\begin{equation}
  \label{eq:lin-extract-cost}
  \mathrm{Cost}_{\mathrm{lin\_extract}}
  \;=\;\mathcal{O}\left(
  \mathrm{Poly}\right).
\end{equation}


\subsubsection{Runtime algorithm}


The runtime solver computes a triangular DP table $S[i][j]$ for $0 \le i \le j \le 2^N$.
The number of DP states is $\mathcal{O}(2^{2N})$.

\paragraph{Per-state cost.}
Fix $(i,j)$ with $0 \le i \le j \le 2^N-1$. Throughout, we will use a constant $C$ that is independent of the parameters of interest and is allowed to change from line to line.\newline
The computation consists of four main steps.

\paragraph{(A) Building $A_j$.}
For a fixed increment $ \mathbf{X}_{j, j+1}$, the linear operator $A_j$ is assembled from the right-linear monomials following equation \eqref{term_A}.\newline
Thus, ignoring dense allocation, the arithmetic work is $\mathcal{O}(\mathrm{Poly}_{\mathrm{lin}})$.
However, the implementation stores $A_j$ as a dense $D\times D$ array, so allocating and zeroing costs $\mathcal{O}(D^2)$.
Since $A_j$ depends only on $j$ (not on $i$) and is cached, this cost is paid only once per $j$
\begin{equation}
  \label{eq:A-build-cost}
   \mathrm{Cost}(A_j)
  \;\leq\;
  C \!\left(
    D^2 + \mathrm{Poly}_{\mathrm{lin}}
  \right),
\end{equation}
for some constant $C \geq 0$ independent of $\mathrm{Poly}_{\mathrm{lin}}$ and $D$.
\paragraph{(B) Building the right-hand side $b_{i,j}$.}
The right hand side is a sum of polynomial evaluations requring the coordinates of the controlled path and the rough path increments.
Each call performs $\mathcal{O}(\mathrm{Poly})$ multiplications plus a scatter-add into $\mathbb{R}^D$.
For fixed $(i,j)$, the algorithm performs one such evaluation plus a sum over $m=i+1,\dots,j$, i.e.\ $\mathcal{O}(2^N)$ calls.
Therefore,
\begin{equation}
  \label{eq:b-build-cost}
  \mathrm{Cost}(b_{i,j})
  \;=\;
  C 2^N \,\mathrm{Poly}.
\end{equation}

\paragraph{(C) Empty-word enforcement (row replacement).}
The empty-word rule, that specifies the trace level identity,  modifies a single row by a scatter-add over $\mathrm{Poly}_{\emptyset}$ entries.
This costs $\mathcal{O}(\mathrm{Poly}_{\emptyset})$ per state, and in practice is lower order compared with
\eqref{eq:b-build-cost} and the linear solve below.

\paragraph{(D) Linear solve.}
Each DP state solves a dense linear system
\[
(Id_D - A_j)\, x = b_{i,j} \in \mathbb{R}^D,
\]
using a direct dense solver. This costs
\begin{equation}
  \label{eq:solve-cost}
  \mathrm{Cost}_{\mathrm{solve}}(i,j)
  \;=\;
  \mathcal{O}(D^3).
\end{equation}

\paragraph{Total runtime cost.}
Summing \eqref{eq:b-build-cost} and \eqref{eq:solve-cost} over the $\mathcal{O}(2^{2N})$ DP states yields
\begin{equation}
  \label{eq:total-runtime}
  \mathrm{Cost}_{\mathrm{runtime}}
  \;\leq\;
  C\!\left(
    2^{2N} (2^N\,\mathrm{Poly} + D^3)
  \right)
  \;\leq\;
  C\left(
    2^{3N}\,\mathrm{Poly} \;+\; 2^{2N}\,D^3
  \right).
\end{equation}
Including the cached assembly of $A_j$ from \eqref{eq:A-build-cost} adds the lower-order term
$C2^N\,(D^2 + \mathrm{Poly}_{\mathrm{lin}})$, however it is worth noting that the addition of this term is not going to alter the estimate above.

\paragraph{Memory complexity.}
The DP table stores $\mathcal{O}(2^{2N})$ tensor elements in $\mathbb{R}^D$, hence
\begin{equation}
  \label{eq:memory}
  \mathrm{Memory}
  \;\leq\;
  C2^{2N} D,
\end{equation}
plus the compiled polynomial storage $\mathcal{O}(\mathrm{Poly})$ and (optionally) the cached dense matrices
$C2^N D^2$ if all $A_j$ are cached simultaneously.

\begin{remark}
In the previous section we assumed that the rough path $\mathbf{X}$ is known, which is rarely the case in practice.  \newline
If the path is sampled on a partition with $2^M$ points, we group the points into blocks of size $2^{M -N}$ and use the corresponding increments to construct the rough path, resulting in $2^{N}$ effective points. The runtime cost therefore increases by
\[ c\left(\tilde{D}(2^N - 1)d^{\tilde{D}} 2^{M-N}\right) \leq 
 \mathrm{SigCompute} \leq C\left(\tilde{D}(2^N - 1)d^{\tilde{D}} 2^{M-N}\right),
\]
where $\tilde{D}$ is the dimension of the tensor algebra $T^{\leq \kappa}(\mathbb{R}^d)$ and $C, c \geq 0$ are independent of all the parameters in the expression above.
\end{remark}

\subsubsection*{Funding}
The work of TC was supported in part by UK Research and Innovation (UKRI) through the Engineering and Physical Sciences Research Council (EPSRC) via a Programme Grant [Grant No. UKRI1010: High order mathematical and computational infrastructure for streamed data that enhance contemporary generative and large language models].\\
DC was partially supported by the European Research Council (ERC) under the European Union’s Horizon 2020 Research and Innovation Programme (ERC, Grant Agreement No 856408).\\
AI  acknowledges support from the Imperial College London Schrödinger Scholarship.\\
WT acknoledges the support by the EPSRC Centre for Doctoral Training in Mathematics of Random Systems: Analysis, Modelling and Simulation (EP/S023925/1).

\bibliographystyle{alpha}
\bibliography{biblio}
\nocite{*}
\addcontentsline{toc}{chapter}{Bibliography}                                                                                           \begin{comment}
\begin{proposition}
The equation of equation
    \begin{equation}\label{SD_equation}
Z_{s,t} = 1 - \sum_{\iota=1}^d \int_s^t \int_s^r Z_{u,r} dX^\iota_u dX^\iota_r
+ \sum_{\iota,j=1}^d 
\int_s^t \int_s^u 
\phi^\iota_{u,t}\; Z_{z,u}\; \sigma^j_{s,z}\; 
dX^\iota_u \, dX^j_z,
\end{equation}
where
\[
\phi^\iota_{u,t} := \int_u^t Z_{u,r}\, dX^\iota_r, 
\qquad
\sigma^j_{s,z} := \int_s^z Z_{v,z}\, dX^j_v,
\] is bounded.
\end{proposition}
\begin{proof}
A contracting ball argument shows that for any couple $(s,t)$ with $a \leq s \leq t \leq b$, with $|t-s| \leq \delta$ the equation \eqref{SD_equation} admits a unique local and continuous solution $Z^{s,t}$, which is bounded above by $M>0$ for any admissible pair $(s,t)$.\newline
For a value $\epsilon > 0$, denote by $\mathcal{P}^\epsilon$  the partition of $[a,b]$  defined as
\begin{equation*}
t_0 := a \quad \quad t_i := \sup\Big\{ z>s_{i-1} : \|X\|_{[z, t_{i-1}]} \leq \epsilon \Big\} \wedge b. 
\end{equation*}
We now will prove by contradiction that 
\[
\sup_{s \in [a,t_i]}\|K(s, t_i)\| \leq M + i +1.
\]
For $t_0$, there is nothing to prove, reasoning by induction we know that 
\begin{equation}\label{auxiliary_difference_boundedness}
K(s,t_{i+1}) - K(s,t_{i}) = - \sum_{\iota=1}^d 
\int_{t_i}^{t_{i+1}} \int_s^r 
Z_{u,r} 
\, dX^\iota_u \, dX^\iota_r  + 
\sum_{\iota,j=1}^d 
\int_{t_i}^{t_{i+1}} \int_s^r \int_s^u \int_s^z
Z_{u,r}Z_{v,z}Z_{z,u} 
\, dX^j_v \, dX^j_z \, dX^\iota_u \, dX^\iota_r,
\end{equation}
and pick $s* = \sup\{ s \leq t_{i+1} \,: |K(s,t_{i+1})| > M + i + 1\}$, we know that by boundedness of local solution we cannot have $s*=t$.\newline
Plugging this into equation \eqref{auxiliary_difference_boundedness} we obtain
\begin{align*}
 \left|K(s^*,t_{i+1}) - K(s^*,t_{i})\right| \leq  \sum_{\iota=1}^d 
\int_{t_i}^{t_{i+1}} \int_s^r 
Z_{u,r} 
\, dX^\iota_u \, dX^\iota_r  + 
\sum_{\iota,j=1}^d 
\int_{t_i}^{t_{i+1}} \int_s^r \int_s^u \int_s^z
Z_{u,r}Z_{v,z}Z_{z,u} 
\, dX^j_v \, dX^j_z \, dX^\iota_u \, dX^\iota_r,   
\end{align*}
\end{proof}


\begin{proposition}
    The equation
    \begin{equation}\label{uniqueness_thm_equation_statement}
Z_{s,t} = 1 - \sum_{\iota=1}^d \int_s^t \int_s^r Z_{s,u} Z_{u,r} dX^\iota_u dX^\iota_r.
\end{equation}
admits a unique solution.
\end{proposition}
\begin{proof}
Notice that by the structure of equation \eqref{uniqueness_thm_equation_statement}, if any solution exists, it has to be bounded.\newline
Fix $s \in [a, b]$, by choosing $t_0$ so that $|t_0-s|$ is sufficiently small it is easy to see that the solution to \eqref{SD_equation_double} is unique in this interval, we will call this solution $Z^*$.\newline
Fix $M > \|Z^*\|_{\infty}$ and consider 
the set
\[
B(\Delta_{[s, t_0 + \delta]}, Z^*) := \{ Z \in C^0(\Delta_{[s, t_0 + \delta]}, \mathbb{R}) \, s.t. \, Z_{|\Delta_{[s, t_0]}} = Z^*, Z_{r,r} = 1 \text{ for every } r \in [s, t_0+ \delta], \|Z\|_{ \infty} \leq M \}
\]
and define the map
\begin{align*}
&\Gamma: B(\Delta_{[s, t_0 + \delta]}, Z^*) \to B(\Delta_{[s, t_0 + \delta]}, Z^*),   \\
& \Gamma(Z) = 
\begin{cases}
Z^*_{s, t} & \text{ when } (s,t) \in \Delta_{[s, t_0]}, \\
1 - \sum_{\iota=1}^d \int_s^t \int_s^r Z_{s,u}Z_{u,r} dX^\iota_u dX^\iota_r
 & \text{ otherwise.}
\end{cases}
\end{align*}
For $Z, \tilde{Z} \in B(\Delta_{[s, t_0 + \delta]}, Z^*)$ we get
\[
\Gamma(Z)_{s,t} - \Gamma(\tilde{Z})_{s,t}
=
\begin{cases}
0, & \text{if } (s,t) \in \Delta_{[s,t_0]}, \\[1.2em]
\displaystyle
- \sum_{\iota=1}^d 
\int_{t_0}^{t} \int_s^r 
Z_{s,u}\left(Z_{u,r} - \tilde{Z}_{u,r}\right)
\, dX^\iota_u \, dX^\iota_r \\[0.8em]
\displaystyle
\quad - \sum_{\iota=1}^d 
\int_{t_0}^{t} \int_s^r 
Z_{u,r}\left(Z_{s,u} - \tilde{Z}_{s,u}\right)
\, dX^\iota_u \, dX^\iota_r,
& \text{otherwise.}
\end{cases}
\]
Focusing on the second case, we have
\begin{align*}
    \left\| \Gamma(Z) - \Gamma(\tilde{Z})\right\|_{\infty} \leq_{b,a, \|X\|_{1}} \delta (M^2 + M + 1)
    \left\| Z - \tilde{Z}\right\|_{\infty} 
\end{align*}
so, choosing $\delta$ appropriately produces a contraction on the simplex $\Delta_{[s, t_0 + \delta]}$. 
\newline
This is enough to conclude: given any two bounded solutions $Z$ and $\tilde{Z}$ with $\|Z\|_{\infty}, \|\tilde{Z}\|_{\infty} \le M$, we can keep the radius of the ball fixed and iterate the argument to cover the simplex $\Delta_{[s, t]}$ for any $t$. Since the argument does not depend on the choice of $s$, the proof is complete.



\end{proof}

For $\kappa \geq 4$, we have for $n=1$ and $n=2$ the classic expansion introduced by Will.
For $n \geq 2$
\begin{align*}
    K_{s,t}^{i_1\dots i_n}
    &= 1 - \int_s^t
        K_{s,u_1}^{i_1 \dots i_{\lfloor 2n/3 \rfloor}}
        K_{u_1, t}^{i_{\lfloor 2n/3 \rfloor+1}  \dots i_n}
        d\mathbf{X}^{i_{\lfloor 2n/3 \rfloor}}_{u_1}
        \\
        &= 1 - \int_s^t K_{u_1, t}^{{i_{\lfloor 2n/3 \rfloor+1}}  \dots i_n}
        d\mathbf{X}^{i_{\lfloor 2n/3 \rfloor}}_{u_1}\\ 
        & \qquad + \int_s^t \int_s^{u_1}
        K_{s,u_2}^{i_1 \dots i_{\lfloor n/3 \rfloor} + 1} K_{u_2, u_1}^{{i_{\lfloor n/3 \rfloor} + 1} \dots {i_{\lfloor 2n/3 \rfloor}}}
        K_{u_1, t}^{{i_{\lfloor 2n/3 \rfloor+1}}  \dots i_n}
        d\mathbf{X}^{i_{\lfloor n/3 \rfloor}}_{u_2} d\mathbf{X}^{i_{\lfloor 2n/3 \rfloor}}_{u_1}
\end{align*}
Notice that in each term we only integrate  Gubinelli derivatives up to order 
\(\frac{\kappa - 2}{3}\). Consequently, in order to compute the 2D Fubini 
rough integral, we must know the terms up to order 
\(\frac{4}{3}\kappa - \frac{5}{3}\), which implies that the system closes whenever 
\(\kappa \leq 5\).

However, given the effective \enquote{smoothness} of the solution to the equation, we can 
in fact exploit the Gubinelli derivatives up to level \(\kappa+j\). This allows us to 
recover that we only need to know terms up to order 
\(\frac{\kappa + j}{3} + \kappa - 1\). Therefore, the system closes if and only if
\[
    \frac{\kappa + j}{3} + \kappa - 1 \leq \kappa + j,
\]
that is,
\[
    j \geq \frac{\kappa - 5}{2}.
\]
This allows to close, provided an appropriate choice for $j$, any system driven a signal of arbitrary large $p$-variation.
\end{comment}

\end{document}
